The sections of this chapter describe the files of \code{Share}, in the
order in which they are compiled. The modules checked so far have unit
tests in \code{Testing/share}, built and run by \code{make check}: each
program prints PASS or FAIL for every check and exits with an error if any
check failed.

\begin{longtable}{L{5cm} L{9cm}}
  \toprule
  \textbf{Test} & \textbf{Module} \\
  \midrule
  \endhead
  \bottomrule
  \endfoot
  \code{test\_date} & \code{mod\_date} (Section~\ref{sec:moddate}) \\
  \code{test\_spbarcoord} & \code{mod\_spbarcoord}
    (Section~\ref{sec:spbarcoord}) \\
  \code{test\_projections} & \code{mod\_projections}
    (Section~\ref{sec:projections}) \\
  \code{test\_hgt} & \code{mod\_hgt} (Section~\ref{sec:hgt}) \\
  \code{test\_vertint} & \code{mod\_vertint} (Section~\ref{sec:vertint}) \\
  \code{test\_humid} & \code{mod\_humid} (Section~\ref{sec:humid}) \\
  \code{test\_interp} & \code{mod\_interp} (Section~\ref{sec:interp}) \\
  \code{test\_spline} & \code{mod\_spline} (Section~\ref{sec:spline}) \\
  \code{test\_memutil} & \code{mod\_space}, \code{mod\_memutil}
    (Section~\ref{sec:memutil}) \\
  \code{test\_kdtree2} & \code{mod\_kdtree2\_priority},
    \code{mod\_kdtree2} (Section~\ref{sec:kdtree2}) \\
  \code{test\_kdinterp} & \code{mod\_kdinterp} (Section~\ref{sec:kdinterp}) \\
  \code{test\_nchelper} & \code{mod\_nchelper} (Section~\ref{sec:nchelper}) \\
  \code{test\_ncstream} & \code{mod\_ncstream} (Section~\ref{sec:ncstream}) \\
  \code{test\_async\_netcdf} & \code{mod\_async\_netcdf}
    (Section~\ref{sec:async_netcdf}) \\
\end{longtable}


%==============================================================================
\section{Basic Numeric Types --- \texttt{mod\_intkinds.F90} and
         \texttt{mod\_realkinds.F90}}
\label{sec:kinds}
%==============================================================================

\subsection{Overview}

The two modules \code{mod\_intkinds} and \code{mod\_realkinds} form the
foundation of RegCM's numeric type system. Every other module in the model
depends --- directly or indirectly --- on one or both of them. Together they:

\begin{itemize}
  \item define portable, machine-independent kind parameters for integers
        and reals at four widths each;
  \item expose a single \emph{working-precision} alias \code{rkx} that is
        resolved at compile time, allowing the entire model to be built in
        either single or double precision by toggling one preprocessor flag;
  \item provide IEEE-754 special-value constants (\code{nan}, \code{inf})
        and the utility predicates \code{is\_nan} and \code{is\_inf} in a
        portable, compiler-agnostic manner.
\end{itemize}

Neither module contains subroutines or data variables (with the sole exception
of \code{init\_realkinds} under NAG Fortran, described below). Both use
\code{public} visibility, so all names they export are directly accessible to
any \code{use}ing module.

\subsection{Dependency Graph}

\begin{lstlisting}
mod_intkinds       (no dependencies)
mod_realkinds
    |-- [optional] iso_fortran_env   (F2008 mode)
    |-- [optional] ieee_arithmetic   (F2008 mode)
mod_constants
    |-- mod_realkinds
        ...
\end{lstlisting}

\code{mod\_intkinds} has no Fortran-level dependencies at all; under
\code{F2008} it uses the intrinsic \code{iso\_fortran\_env} module purely to
import the named kind constants \code{int8}, \code{int16}, \code{int32}, and
\code{int64}.

%------------------------------------------------------------------------------
\subsection{\texttt{mod\_intkinds} --- Integer Kind Parameters}
\label{subsec:intkinds}

\subsubsection{Purpose}

\code{mod\_intkinds} provides four named integer kind values (\code{ik1},
\code{ik2}, \code{ik4}, \code{ik8}) covering 1-byte to 8-byte widths, plus
a convenience constant \code{bigint} equal to \code{HUGE(integer(ik4))}.
Having named kinds avoids relying on the default integer kind, which is
implementation-defined, and prevents silent truncation when mixing 32-bit
and 64-bit integers in index or counter arithmetic.

\subsubsection{Compile-Time Modes}

The module supports two implementations selected by the preprocessor flag
\code{F2008}:

\begin{longtable}{L{3.5cm} L{5.5cm} L{5cm}}
  \toprule
  \textbf{Name} & \textbf{\texttt{F2008} mode} & \textbf{Fallback (Fortran 95/2003)} \\
  \midrule
  \endhead
  \bottomrule
  \endfoot
  \code{ik1} & \code{int8}  & \code{selected\_int\_kind(R=2)}  \\
  \code{ik2} & \code{int16} & \code{selected\_int\_kind(R=4)}  \\
  \code{ik4} & \code{int32} & \code{selected\_int\_kind(R=9)}  \\
  \code{ik8} & \code{int64} & \code{selected\_int\_kind(R=18)} \\
\end{longtable}

The \code{F2008} path binds directly to the \code{iso\_fortran\_env} named
constants, which are guaranteed by the standard to map to the corresponding
ISO C fixed-width types. The fallback path uses \code{selected\_int\_kind(R)}
which requests the smallest available kind that can represent integers with at
least $R$ decimal digits (e.g.\ $R=9$ gives at least $\pm10^9$, safely
covering the 32-bit signed range).

\subsubsection{Kind Parameters}

\begin{longtable}{L{2.5cm} L{2cm} L{3cm} L{3cm} L{4.5cm}}
  \toprule
  \textbf{Name} & \textbf{Width} & \textbf{Range} & \textbf{Typical use} & \textbf{Note} \\
  \midrule
  \endhead
  \bottomrule
  \endfoot
  \code{ik1} & 8-bit  & $[-128,\,127]$                     & Flags, masks           & Rarely used directly \\
  \code{ik2} & 16-bit & $[-32768,\,32767]$                 & Compact index fields   & \\
  \code{ik4} & 32-bit & $[\approx{-2.1\times10^9},\,\approx{2.1\times10^9}]$ & Loop indices, counts & Default in most loops \\
  \code{ik8} & 64-bit & $[\approx{-9.2\times10^{18}},\,\approx{9.2\times10^{18}}]$ & Large file offsets, NetCDF record counts & \\
\end{longtable}

\subsubsection{\texttt{bigint} --- Maximum 32-bit Signed Integer}

\begin{lstlisting}
integer(ik4), parameter :: bigint = 2147483647   ! = HUGE(integer(ik4))
\end{lstlisting}

\code{bigint} holds the value $2^{31}-1 = 2\,147\,483\,647$, the largest
representable \code{ik4} integer. It is defined via the preprocessor macro
\code{\_\_SYSTEM\_INTMAX\_32\_\_} (which also exists as a 64-bit variant
\code{\_\_SYSTEM\_INTMAX\_64\_\_} for documentation purposes but is not exposed
as a named parameter). \code{bigint} is used in the rest of the model wherever
a sentinel ``not-set'' or ``unlimited'' integer value is needed.

%------------------------------------------------------------------------------
\subsection{\texttt{mod\_realkinds} --- Real Kind Parameters and IEEE Utilities}
\label{subsec:realkinds}

\subsubsection{Purpose}

\code{mod\_realkinds} provides three real kind parameters (\code{rk4},
\code{rk8}, \code{rk16}) and the critical working-precision alias \code{rkx},
together with IEEE-754 special values (\code{nan}, \code{inf} and their
precision-specific variants) and two generic interfaces (\code{is\_nan},
\code{is\_inf}) for portable floating-point diagnostics.

\subsubsection{Compile-Time Modes}

Three orthogonal preprocessor flags interact in this module:

\begin{longtable}{L{4cm} L{10cm}}
  \toprule
  \textbf{Flag} & \textbf{Effect} \\
  \midrule
  \endhead
  \bottomrule
  \endfoot
  \code{F2008} &
    Uses \code{iso\_fortran\_env} for kind constants (\code{real32},
    \code{real64}, \code{real128}) and \code{ieee\_arithmetic} for
    special-value construction and IEEE class queries.
    Without this flag the module falls back to \code{kind(1.0)},
    \code{selected\_real\_kind}, and octal-literal bit patterns. \\
  \code{NAGFOR} &
    Active only when \code{F2008} is also set. NAG Fortran does not allow
    \code{ieee\_value} in a \code{parameter} initialiser, so \code{nan}
    and \code{inf} are declared as non-constant variables initialised by
    the auxiliary subroutine \code{init\_realkinds}, which must be called
    once at model startup. \\
  \code{SINGLE\_PRECISION\_REAL} &
    Sets \code{rkx = rk4} (32-bit working precision). The default (flag
    absent) sets \code{rkx = rk8} (64-bit working precision). \\
\end{longtable}

\subsubsection{Kind Parameters}

\begin{longtable}{L{0.9cm} L{1.2cm} L{2.2cm} L{7cm} L{4cm}}
  \toprule
  \textbf{Name} & \textbf{Width} & \textbf{\texttt{F2008} value} & \textbf{Fallback value} & \textbf{Description} \\
  \midrule
  \endhead
  \bottomrule
  \endfoot
  \code{rk4}  & 32-bit  & \code{real32}  & \code{kind(1.0)}                          & Single precision (IEEE-754 \texttt{binary32}) \\
  \code{rk8}  & 64-bit  & \code{real64}  & \code{selected\_real\_kind(2*precision(rk4))} & Double precision (IEEE-754 \texttt{binary64}) \\
  \code{rk16} & 128-bit & \code{real128} & \code{selected\_real\_kind(2*precision(rk8))} & Quad precision; available on supported compilers only \\
  \code{rkx}  & ---     & \code{rk4} or \code{rk8} & same & \textbf{Model working precision} (see below) \\
\end{longtable}

The fallback definitions for \code{rk8} and \code{rk16} follow the
``double-the-precision'' pattern recommended in Metcalf, Reid \& Cohen,
\textit{Fortran 95/2003 Explained} (2004, p.~71): each successive kind is
requested to have at least twice the decimal precision of the previous one,
yielding approximately 7, 15, and 30 significant decimal digits respectively.

\subsubsection{The Working Precision \texttt{rkx}}

\code{rkx} is the single most important symbol in the RegCM type system.
Every physical field, every named constant in \code{mod\_constants}, and every
function argument in the \code{Share/*.inc} include files is typed
\code{real(rkx)}. This means:

\begin{itemize}
  \item Compiling \emph{without} \code{SINGLE\_PRECISION\_REAL} (the default)
        produces a fully double-precision executable --- \code{rkx = rk8}.
  \item Compiling \emph{with} \code{SINGLE\_PRECISION\_REAL} produces a
        single-precision executable --- \code{rkx = rk4} --- at roughly half
        the memory footprint and higher vectorisation throughput, at the cost
        of reduced accuracy.
\end{itemize}

Internal computations inside the thermodynamic kernels
(e.g.\ the Flatau polynomial evaluations in \code{pfesat.inc}) promote
intermediate values to \code{rk8} regardless of \code{rkx}, then cast the
result back. This guards against rounding error in the polynomial evaluation
when the model is built in single precision.

\subsubsection{IEEE Special Values}

Four precision-specific constants and two working-precision aliases are provided:

\begin{longtable}{L{2cm} L{7cm} L{5cm}}
  \toprule
  \textbf{Name} & \textbf{Value / construction} & \textbf{Description} \\
  \midrule
  \endhead
  \bottomrule
  \endfoot
  \code{nan\_r4} &
    \code{transfer(-4194304\_int32, 1.\_real32)} &
    Quiet NaN, 32-bit \\
  \code{inf\_r4} &
    \code{transfer(2139095040\_int32, 1.\_real32)} &
    Positive infinity, 32-bit \\
  \code{nan\_r8} &
    \code{transfer(-2251799813685248\_int64, 1.\_real64)} &
    Quiet NaN, 64-bit \\
  \code{inf\_r8} &
    \code{transfer(9218868437227405312\_int64, 1.\_real64)} &
    Positive infinity, 64-bit \\
  \code{nan}     & alias of \code{nan\_r4} or \code{nan\_r8} & Working-precision NaN \\
  \code{inf}     & alias of \code{inf\_r4} or \code{inf\_r8} & Working-precision positive infinity \\
\end{longtable}

The integer bit patterns used in the \code{transfer} calls are hardwired via
the preprocessor macros \code{\_\_SYSTEM\_NAN\_32\_\_} etc., whose values
were generated by a small self-contained test program (preserved as a comment
at the top of the source file) that calls \code{ieee\_value} and converts the
result to integer via \code{transfer}. This two-step approach allows the
constants to be \code{parameter} (compile-time) rather than requiring runtime
initialisation, while remaining binary-portable across IEEE-754 compliant
platforms.

Under \textbf{NAG Fortran} with \code{F2008}, the standard prohibits
\code{ieee\_value} in a \code{parameter} context; in that case \code{nan} and
\code{inf} are declared as plain (non-constant) module variables and must be
set by calling \code{init\_realkinds()} before any other module is used.

Under the Fortran 95/2003 fallback (no \code{F2008}), the special values are
constructed from octal literals (\code{O'07760000000'} for single-precision
\code{inf}, etc.), which directly encode the IEEE-754 bit patterns and are
accepted as \code{parameter} initialisers by most legacy compilers.

\subsubsection{Generic Interfaces \texttt{is\_nan} and \texttt{is\_inf}}

Two overloaded generic interfaces provide portable special-value detection:

\begin{lstlisting}
interface is_nan
  module procedure is_nan_single   ! real(rk4)
  module procedure is_nan_double   ! real(rk8)
end interface

interface is_inf
  module procedure is_inf_single   ! real(rk4)
  module procedure is_inf_double   ! real(rk8)
end interface
\end{lstlisting}

All four underlying functions are \code{logical elemental}, making them
applicable to arrays of any shape. Their implementations differ by mode:

\begin{longtable}{L{3cm} L{5.5cm} L{5.5cm}}
  \toprule
  \textbf{Function} & \textbf{\texttt{F2008} implementation} & \textbf{Fallback implementation} \\
  \midrule
  \endhead
  \bottomrule
  \endfoot
  \code{is\_nan\_*} &
    \code{ieee\_class(x) == ieee\_quiet\_nan .or.} \newline
    \code{ieee\_class(x) == ieee\_signaling\_nan} &
    \code{(x /= x) .or.} \newline
    \code{((x > 0) .eqv. (x <= 0))} \\
  \code{is\_inf\_*} &
    \code{ieee\_class(x) == ieee\_negative\_inf .or.} \newline
    \code{ieee\_class(x) == ieee\_positive\_inf} &
    \code{x > huge(x) .or. x < -huge(x)} \\
\end{longtable}

The fallback NaN test exploits the IEEE-754 property that NaN is unordered
with respect to any value, including itself (\code{x /= x} is \code{.true.}
for NaN). The second clause \code{(x > 0) .eqv. (x <= 0)} is a supplementary
guard for compilers that may short-circuit or optimise away the self-comparison.
The fallback Inf test relies on \code{huge(x)} returning the largest finite
representable value, so a value strictly greater in magnitude must be infinite.

%------------------------------------------------------------------------------
\subsection{Design Notes}

\textbf{Foundation role.} These two modules sit at the root of the entire
dependency tree. No other \code{Share/} module or physics parameterisation
can be compiled without them, yet they themselves have no RegCM-level
dependencies. Changes to either file propagate to every module at the next
full build.

\textbf{Single entry point for precision.} The \code{rkx} alias in
\code{mod\_realkinds} is the sole control point for model-wide floating-point
precision. Changing precision does not require modifying any physics code;
the compiler resolves all \code{real(rkx)} declarations and all \code{rkx}
literal suffixes through the kind parameter at compile time.

\textbf{Portability layers.} The three-level structure ---
\code{F2008} $\supset$ \code{NAGFOR} exception $\supset$ Fortran~95 fallback
--- ensures that the model can be built with compilers ranging from those
supporting only the Fortran~95 standard to fully conforming Fortran~2008
implementations, including the stricter NAG compiler. The IEEE special-value
approach (precomputed bit patterns as \code{parameter} constants, with a
runtime fallback for NAG) is a pragmatic solution to the long-standing
portability problem of NaN and Inf initialisation in module scope.

\textbf{No \texttt{rk16} in physics.} Although \code{rk16} (quad precision)
is defined, it is not used in any physics module. It is available for
developers who need extended-precision diagnostics or reference computations
during code validation.

%======================================================================
\section{Constants --- \texttt{mod\_constants.F90}}
\label{sec:constants}
%======================================================================

\subsection{Overview}

\code{mod\_constants} is the single authoritative source of all named
physical, mathematical, and numerical constants used throughout the RegCM
model. It is a pure parameter module --- it declares no variables, no
procedures, and no derived types --- and exposes everything with
\code{public} visibility.

The module depends only on \code{mod\_realkinds}, which defines the
precision kinds \code{rk4}, \code{rk8}, and \code{rkx} (the model working
precision). This single dependency means that the precision of nearly every
constant in the model can be controlled by editing \code{mod\_realkinds} alone.

\textbf{Design philosophy:}
\begin{itemize}
  \item All constants are \code{parameter} (compile-time literals) to allow
        constant folding.
  \item Precision-sensitive constants (geometry, trigonometry) are provided
        in all three precisions (\code{\_r4}, \code{\_r8}, \code{rkx}) so
        callers can avoid implicit conversion.
  \item Two compile-time preprocessor flags alter a subset of values:
    \begin{itemize}
      \item \texttt{RCEMIP} --- selects values appropriate for the
            Radiative-Convective Equilibrium Model Intercomparison Project.
      \item \texttt{SINGLE\_PRECISION\_REAL} --- adjusts the tracer minimum
            threshold.
    \end{itemize}
\end{itemize}

\subsection{Dependency}

\begin{lstlisting}
mod_constants
    |-- mod_realkinds   (provides: rk4, rk8, rkx)
\end{lstlisting}

\subsection{Compile-Time Flags}

\begin{longtable}{L{4cm} L{10cm}}
  \toprule
  \textbf{Flag} & \textbf{Effect} \\
  \midrule
  \endhead
  \bottomrule
  \endfoot
  \code{RCEMIP} &
    Uses RCEMIP-prescribed values for \code{egrav}, \code{rgas},
    \code{rwat}, \code{cpd}, \code{cpv}, \code{cvd}, \code{wlhv},
    \code{wlhf}, \code{wlhs}, \code{lrate}, \code{earthrad},
    \code{minqq}, \code{minqc}. \\
  \code{SINGLE\_PRECISION\_REAL} &
    Raises \code{mintr} from $10^{-30}$ to $10^{-20}$ to stay above the
    single-precision underflow threshold. \\
\end{longtable}

%----------------------------------------------------------------------
\subsection{Group 1 --- Numeric Literals}

Plain real literals named to improve readability and avoid magic numbers.
All are \code{real(rkx)} unless noted.

\begin{longtable}{L{3cm} L{3.5cm} L{7cm}}
  \toprule
  \textbf{Name} & \textbf{Value} & \textbf{Description} \\
  \midrule
  \endhead
  \bottomrule
  \endfoot
  \code{d\_zero}   & 0.0       & Zero \\
  \code{d\_one}    & 1.0       & One \\
  \code{d\_two}    & 2.0       & Two \\
  \code{d\_three}  & 3.0       & Three \\
  \code{d\_four}   & 4.0       & Four \\
  \code{d\_five}   & 5.0       & Five \\
  \code{d\_six}    & 6.0       & Six \\
  \code{d\_nine}   & 9.0       & Nine \\
  \code{d\_half}   & 0.5       & One half \\
  \code{d\_rfour}  & 0.25      & One quarter \\
  \code{d\_twelve} & 12.0      & Twelve \\
  \code{d\_10}     & 10.0      & Ten \\
  \code{d\_r10}    & 0.1       & $1/10$ \\
  \code{d\_100}    & 100.0     & One hundred \\
  \code{d\_r100}   & 0.01      & $1/100$ \\
  \code{d\_1000}   & 1000.0    & One thousand \\
  \code{d\_r1000}  & 0.001     & $1/1000$ \\
  \code{onet}      & $1/3$     & One third \\
  \code{twot}      & $2/3$     & Two thirds \\
  \code{fourt}     & $4/3$     & Four thirds \\
  \code{sqrt2}     & $\sqrt{2}$        & Square root of two \\
  \code{rsqrt2}    & $1/\sqrt{2}$      & Reciprocal of $\sqrt{2}$ \\
\end{longtable}

\textbf{Angle literals} (\code{real(rk8)}, operations must be in double precision):

\begin{longtable}{L{3cm} L{3.5cm} L{7cm}}
  \toprule
  \textbf{Name} & \textbf{Value} & \textbf{Description} \\
  \midrule
  \endhead
  \bottomrule
  \endfoot
  \code{deg00}  & 0.0   & $0^\circ$ \\
  \code{deg45}  & 45.0  & $45^\circ$ \\
  \code{deg90}  & 90.0  & $90^\circ$ \\
  \code{deg180} & 180.0 & $180^\circ$ \\
  \code{deg360} & 360.0 & $360^\circ$ \\
\end{longtable}

%----------------------------------------------------------------------
\subsection{Group 2 --- Minimum Threshold Values}

Lower bounds applied to moisture and tracer fields to prevent numerical
underflow, negative values, or division by zero.

\begin{longtable}{L{2.5cm} L{3.5cm} L{3.5cm} L{4.5cm}}
  \toprule
  \textbf{Name} & \textbf{Standard} & \textbf{RCEMIP} & \textbf{Description} \\
  \midrule
  \endhead
  \bottomrule
  \endfoot
  \code{minqq} & $10^{-8}$  & $10^{-14}$ & Generic moisture minimum \\
  \code{minqc} & $10^{-10}$ & $10^{-16}$ & Cloud condensate minimum \\
  \code{minqv} & $100\times\text{minqq}$ & --- & Water vapour minimum \\
  \code{mintr} & $10^{-30}$ (double) & --- & Tracer minimum; $10^{-20}$ single \\
  \code{minqx} & $10^{-16}$ & --- & Extra-strict microphysics minimum \\
\end{longtable}

Missing/sentinel values (all $\approx 10^{\pm20}$; provided in \code{rk4},
\code{rk8} and \code{rkx} variants for NetCDF I/O compatibility):
\code{dlowval}, \code{dhival}, \code{slowval}, \code{shival},
\code{dmissval}, \code{smissval}.

%----------------------------------------------------------------------
\subsection{Group 3 --- Time Conversion}

\begin{longtable}{L{3cm} L{3cm} L{3cm} L{5cm}}
  \toprule
  \textbf{Name} & \textbf{Value} & \textbf{Units} & \textbf{Description} \\
  \midrule
  \endhead
  \bottomrule
  \endfoot
  \code{secpm}  & 60.0     & s min$^{-1}$ & Seconds per minute \\
  \code{secph}  & 3600.0   & s h$^{-1}$   & Seconds per hour \\
  \code{secpd}  & 86400.0  & s d$^{-1}$   & Seconds per day \\
  \code{rsecpd} & 1/86400  & d s$^{-1}$   & Reciprocal seconds per day \\
  \code{minph}  & 60.0     & min h$^{-1}$ & Minutes per hour \\
  \code{minpd}  & 1440.0   & min d$^{-1}$ & Minutes per day \\
  \code{houpd}  & 24.0     & h d$^{-1}$   & Hours per day \\
\end{longtable}

%----------------------------------------------------------------------
\subsection{Group 4 --- Fundamental Physical Constants}

\subsubsection{Universal Constants}

\begin{longtable}{L{3.5cm} L{4.5cm} L{2.5cm} L{3cm}}
  \toprule
  \textbf{Name} & \textbf{Value} & \textbf{Units} & \textbf{Standard} \\
  \midrule
  \endhead
  \bottomrule
  \endfoot
  \code{egrav}        & 9.80665 / 9.79764 (RCEMIP)    & m s$^{-2}$       & 3rd CGPM \\
  \code{speedoflight} & 299\,792\,458.0               & m s$^{-1}$       & SI 2019 (exact) \\
  \code{plankconstant}& $6.62607015\times10^{-34}$    & J s              & SI 2019 \\
  \code{sigm}         & $5.670374419\times10^{-8}$    & W m$^{-2}$ K$^{-4}$ & Stefan-Boltzmann, SI 2019 \\
  \code{boltzk}       & $1.3806490\times10^{-23}$     & J K$^{-1}$       & Boltzmann, SI 2019 \\
  \code{navgdr}       & $6.02214076\times10^{23}$     & mol$^{-1}$       & Avogadro, SI 2019 \\
  \code{rgasmol}      & $\text{navgdr}\times\text{boltzk}\approx8.31446$ & J mol$^{-1}$ K$^{-1}$ & Universal gas constant \\
  \code{m\_euler}     & 0.57721566\ldots               & ---              & Euler-Mascheroni (rk4, rk8, rkx) \\
\end{longtable}

\subsubsection{Molecular Weights}

All in g mol$^{-1}$. Source: CRC Handbook (1997) for air composition;
PubChem for trace gases.

\begin{longtable}{L{3cm} L{3cm} L{4cm} L{4cm}}
  \toprule
  \textbf{Name} & \textbf{Value} & \textbf{Species} & \textbf{Note} \\
  \midrule
  \endhead
  \bottomrule
  \endfoot
  \code{amd}    & 28.96454  & Dry air          & Weighted composition: N$_2$, O$_2$, Ar, CO$_2$, traces \\
  \code{amw}    & 18.01528  & Water (H$_2$O)   & \\
  \code{amo2}   & 31.9988   & O$_2$            & \\
  \code{amo3}   & 47.9982   & O$_3$            & \\
  \code{amco2}  & 44.0095   & CO$_2$           & \\
  \code{amn2o}  & 44.0128   & N$_2$O           & \\
  \code{amch4}  & 16.04246  & CH$_4$           & \\
  \code{amcfc11}& 137.3681  & CFC-11 (CCl$_3$F)& \\
  \code{amcfc12}& 120.9135  & CFC-12 (CCl$_2$F$_2$)& \\
  \code{pdbratio}& 0.0112372& ---              & $^{13}$C/$^{12}$C Pee Dee Belemnite standard \\
\end{longtable}

%----------------------------------------------------------------------
\subsection{Group 5 --- Gas Constants and Specific Heats}

\subsubsection{Gas Constants}

\begin{longtable}{L{2.8cm} L{3.5cm} L{3cm} L{2cm} L{2.5cm}}
  \toprule
  \textbf{Name} & \textbf{Standard} & \textbf{RCEMIP} & \textbf{Units} & \textbf{Description} \\
  \midrule
  \endhead
  \bottomrule
  \endfoot
  \code{rgas} (\code{rdry}) & $(R_\text{mol}/M_d)\times1000\approx287.057$ & 287.04 & J kg$^{-1}$ K$^{-1}$ & Dry air \\
  \code{rwat}               & $(R_\text{mol}/M_w)\times1000\approx461.523$ & 461.50 & J kg$^{-1}$ K$^{-1}$ & Water vapour \\
  \code{rgow}               & $R_d/R_v\approx0.6220$ & --- & --- & Ratio $R_d/R_v$ \\
  \code{rgowi}              & $R_v/R_d\approx1.6078$ & --- & --- & Ratio $R_v/R_d$ \\
\end{longtable}

\subsubsection{Specific Heats for Air}

\begin{longtable}{L{2.5cm} L{3.5cm} L{3cm} L{2cm} L{3.5cm}}
  \toprule
  \textbf{Name} & \textbf{Standard} & \textbf{RCEMIP} & \textbf{Units} & \textbf{Description} \\
  \midrule
  \endhead
  \bottomrule
  \endfoot
  \code{cpd} & $\frac{7}{2}R_d\approx1004.70$ & 1004.64 & J kg$^{-1}$ K$^{-1}$ & Dry air at constant pressure \\
  \code{cpv} & $4R_v\approx1846.09$           & 1846.00 & J kg$^{-1}$ K$^{-1}$ & Water vapour at constant pressure \\
  \code{cvd} & $\frac{5}{2}R_d\approx717.64$  & 717.6   & J kg$^{-1}$ K$^{-1}$ & Dry air at constant volume \\
\end{longtable}

\textit{Note:} The source comment ``Assume dry air is biatomic gas'' reflects
the idealised diatomic derivation $c_p = \frac{7}{2}R$, $c_v = \frac{5}{2}R$.

\subsubsection{Specific Heats for Water Phases}

\begin{longtable}{L{3cm} L{3cm} L{3cm} L{5cm}}
  \toprule
  \textbf{Name} & \textbf{Value} & \textbf{Units} & \textbf{Description} \\
  \midrule
  \endhead
  \bottomrule
  \endfoot
  \code{cpw}      & 4186.95              & J kg$^{-1}$ K$^{-1}$   & Liquid water at 15\,°C \\
  \code{cpw0}     & 4218.0               & J kg$^{-1}$ K$^{-1}$   & Liquid water at 0\,°C \\
  \code{cpi}      & 2106.00              & J kg$^{-1}$ K$^{-1}$   & Ice at 0\,°C (ECMWF value) \\
  \code{ch2o}     & $4.18695\times10^6$  & J m$^{-3}$ K$^{-1}$    & Liquid water volumetric, ISO 15\,°C \\
  \code{shice}    & $0.46\times\text{ch2o}$ & J m$^{-3}$ K$^{-1}$ & Solid ice \\
  \code{csnw}     & $0.20\times\text{ch2o}$ & J m$^{-3}$ K$^{-1}$ & Wind-packed snow \\
  \code{spcpfw}   & 4188.0               & J kg$^{-1}$ K$^{-1}$   & Fresh water \\
  \code{spcpsw}   & 3996.0               & J kg$^{-1}$ K$^{-1}$   & Sea water \\
  \code{spcpice}  & 2117.27              & J kg$^{-1}$ K$^{-1}$   & Fresh ice \\
\end{longtable}

%----------------------------------------------------------------------
\subsection{Group 6 --- Latent Heats and Thermodynamic Compounds}

\subsubsection{Latent Heats at 0\,°C}

\begin{longtable}{L{3cm} L{4cm} L{3.5cm} L{3.5cm} L{2cm}}
  \toprule
  \textbf{Name} & \textbf{Standard} & \textbf{RCEMIP} & \textbf{Description} & \textbf{Units} \\
  \midrule
  \endhead
  \bottomrule
  \endfoot
  \code{wlhv} & 2\,500\,610.0              & $2.501\times10^6$ & Vaporisation   & J kg$^{-1}$ \\
  \code{wlhf} & 333\,560.5                 & $3.337\times10^5$ & Fusion         & J kg$^{-1}$ \\
  \code{wlhs} & $\text{wlhv}+\text{wlhf}$ & $2.834\times10^6$ & Sublimation    & J kg$^{-1}$ \\
\end{longtable}

\subsubsection{Precomputed Utility Compounds}

\begin{longtable}{L{3cm} L{4.5cm} L{2cm} L{5cm}}
  \toprule
  \textbf{Name} & \textbf{Expression} & \textbf{Units} & \textbf{Description} \\
  \midrule
  \endhead
  \bottomrule
  \endfoot
  \code{regrav}   & $1/g$                    & s$^2$ m$^{-1}$  & Reciprocal gravity \\
  \code{rcpd}     & $1/c_{pd}$               & kg K J$^{-1}$   & Reciprocal of $c_{pd}$ \\
  \code{rovcp}    & $R_d/c_{pd}\approx0.2857$& ---             & Poisson exponent $\kappa$ \\
  \code{cpovr}    & $c_{pd}/R_d\approx3.5$   & ---             & $c_p/R_d$ \\
  \code{gdry}     & $-g/c_{pd}$              & K m$^{-1}$      & Dry adiabatic lapse rate (negative) \\
  \code{govcp}    & $g/c_{pd}$               & K m$^{-1}$      & Magnitude of dry lapse rate \\
  \code{wlhvocp}  & $L_v/c_{pd}$             & K               & Vaporisation latent heat / $c_p$ \\
  \code{wlhsocp}  & $L_s/c_{pd}$             & K               & Sublimation latent heat / $c_p$ \\
  \code{wlhfocp}  & $L_f/c_{pd}$             & K               & Fusion latent heat / $c_p$ \\
  \code{hcratio}  & $c_{pv}/c_{pd}$          & ---             & Heat capacity ratio for moist corrections \\
  \code{hcrm1}    & $c_{pv}/c_{pd}-1$        & ---             & Moist correction increment \\
\end{longtable}

%----------------------------------------------------------------------
\subsection{Group 7 --- Temperature and Saturation Pressure Constants}

\subsubsection{Reference Temperatures and Phase Thresholds}

\begin{longtable}{L{3.5cm} L{3cm} L{2cm} L{5.5cm}}
  \toprule
  \textbf{Name} & \textbf{Value} & \textbf{Units} & \textbf{Description} \\
  \midrule
  \endhead
  \bottomrule
  \endfoot
  \code{tzero}           & 273.15  & K & Ice-point (0\,°C) \\
  \code{tiso}            & 216.65  & K & ICAO isothermal layer temperature (11--20\,km) \\
  \code{wattp}           & 273.16  & K & Triple point of water \\
  \code{tboil}           & 373.1339& K & Boiling point at 1\,atm \\
  \code{rtber}           & 268.15  & K & Threshold: mixed-phase / liquid cloud \\
  \code{rtice}           & 250.15  & K & Threshold: pure ice cloud \\
  \code{mpcrt}           & $\approx256.4$& K & Mixed-phase critical temperature \\
  \code{rtwat\_rtice\_r} & $1/(T_w-T_i)$& K$^{-1}$ & Reciprocal of liquid--ice temperature range \\
\end{longtable}

\subsubsection{Tetens Formula Coefficients}

Used by \code{tetens\_pfesat.inc} and \code{tetens\_pfesdt.inc}.
The formula is $e_s = c_{1es}\exp\!\bigl(c_{3x}(T-T_w)/(T-c_{4x})\bigr)$.

\begin{longtable}{L{3cm} L{3cm} L{8cm}}
  \toprule
  \textbf{Name} & \textbf{Value} & \textbf{Description} \\
  \midrule
  \endhead
  \bottomrule
  \endfoot
  \code{c1es}    & 610.78          & Reference saturation pressure at \code{wattp} [Pa] \\
  \code{c3les}   & 17.2693882      & Exponent numerator coefficient -- liquid \\
  \code{c3ies}   & 21.875          & Exponent numerator coefficient -- ice \\
  \code{c4les}   & 35.86\,K        & Exponent denominator offset -- liquid \\
  \code{c4ies}   & 7.66\,K         & Exponent denominator offset -- ice \\
  \code{c5les}   & $c_{3l}(T_0-c_{4l})$ & Precomputed product -- liquid \\
  \code{c5ies}   & $c_{3i}(T_0-c_{4i})$ & Precomputed product -- ice \\
  \code{c5alvcp} & $c_{5l}\,L_v/c_p$    & Compound for equivalent potential temperature -- liquid \\
  \code{c5alscp} & $c_{5i}\,L_s/c_p$    & Compound for equivalent potential temperature -- ice \\
\end{longtable}

\subsubsection{Kain-Fritsch / WSM5 Saturation Constants}

Alternative exponential fits used in convective and microphysics schemes.

\begin{longtable}{L{3cm} L{3cm} L{8cm}}
  \toprule
  \textbf{Name} & \textbf{Value} & \textbf{Description} \\
  \midrule
  \endhead
  \bottomrule
  \endfoot
  \code{aliq} & 613.3   & Reference pressure -- liquid [Pa] \\
  \code{bliq} & 17.502  & Exponent numerator -- liquid \\
  \code{cliq} & 4780.8  & Exponent denominator offset -- liquid [K] \\
  \code{dliq} & 32.19   & Exponent denominator shift -- liquid [K] \\
  \code{aice} & 613.20  & Reference pressure -- ice [Pa] \\
  \code{bice} & 22.452  & Exponent numerator -- ice \\
  \code{cice} & 6133.0  & Exponent denominator offset -- ice [K] \\
  \code{dice} & 0.61    & Exponent denominator shift -- ice [K] \\
\end{longtable}

%----------------------------------------------------------------------
\subsection{Group 8 --- Standard Atmosphere (ICAO 1993)}

\begin{longtable}{L{3cm} L{3.5cm} L{3cm} L{5cm}}
  \toprule
  \textbf{Name} & \textbf{Value} & \textbf{Units} & \textbf{Description} \\
  \midrule
  \endhead
  \bottomrule
  \endfoot
  \code{p00}    & $1.00000\times10^5$ & Pa   & Reference pressure \\
  \code{stdp}   & $1.01325\times10^5$ & Pa   & Standard sea-level pressure \\
  \code{stdpmb} & 1013.250            & hPa  & \code{stdp} in millibar \\
  \code{stdpcb} & 101.325             & cbar & \code{stdp} in centibar \\
  \code{stdt}   & 288.15              & K    & Standard sea-level temperature (15\,°C) \\
  \code{stdrho} & 1.2250              & kg m$^{-3}$ & Standard sea-level air density \\
  \code{lrate}  & 0.00649 / 0.0067 (RCEMIP) & K m$^{-1}$ & Standard tropospheric lapse rate \\
\end{longtable}

%----------------------------------------------------------------------
\subsection{Group 9 --- Geometric and Geodetic Constants}

\subsubsection{Earth Shape and Rotation}

All quantities are provided in \code{rk4}, \code{rk8}, and \code{rkx} variants.

\begin{longtable}{L{3.5cm} L{4.5cm} L{2.5cm} L{3.5cm}}
  \toprule
  \textbf{Name} & \textbf{Standard value} & \textbf{Units} & \textbf{Description} \\
  \midrule
  \endhead
  \bottomrule
  \endfoot
  \code{earthrad}  & $6.371229\times10^6$ & m          & Mean Earth radius \\
  \code{erkm}      & $6371.229$           & km         & Earth radius in kilometres \\
  \code{rearthrad} & $1/\text{earthrad}$  & m$^{-1}$   & Reciprocal Earth radius \\
  \code{eomeg}     & $7.2921159\times10^{-5}$ & rad s$^{-1}$ & Earth angular velocity $\Omega$ \\
  \code{eomeg2}    & $2\Omega$            & rad s$^{-1}$& Coriolis parameter at equator \\
\end{longtable}

\textit{RCEMIP:} \code{earthrad} $= 6.371\times10^6$\,m.

\subsubsection{Trigonometric Constants}

Provided in \code{rk4}, \code{rk8}, and working-precision \code{rkx}.

\begin{longtable}{L{3.5cm} L{4cm} L{6.5cm}}
  \toprule
  \textbf{Name} & \textbf{Value} & \textbf{Description} \\
  \midrule
  \endhead
  \bottomrule
  \endfoot
  \code{mathpi}  & $3.14159265358979\ldots$ & $\pi$ \\
  \code{invpi}   & $1/\pi$                  & Reciprocal of $\pi$ \\
  \code{halfpi}  & $\pi/2$                  & $\pi/2$ \\
  \code{quartpi} & $\pi/4$                  & $\pi/4$ \\
  \code{twopi}   & $2\pi$                   & $2\pi$ \\
  \code{pisqr}   & $\pi^2$                  & $\pi$ squared \\
  \code{degrad}  & $\pi/180$                & Degrees-to-radians conversion \\
  \code{raddeg}  & $180/\pi$                & Radians-to-degrees conversion \\
\end{longtable}

%----------------------------------------------------------------------
\subsection{Group 10 --- Surface, Turbulence, and Vegetation}

\begin{longtable}{L{3.5cm} L{3cm} L{3cm} L{4.5cm}}
  \toprule
  \textbf{Name} & \textbf{Value} & \textbf{Units} & \textbf{Description} \\
  \midrule
  \endhead
  \bottomrule
  \endfoot
  \code{zlnd}    & 0.01    & m            & Aerodynamic roughness length, soil \\
  \code{zoce}    & 0.00023 & m            & Aerodynamic roughness length, ocean \\
  \code{zsno}    & 0.00040 & m            & Aerodynamic roughness length, snow \\
  \code{vonkar}  & 0.4     & ---          & Von Kármán constant \\
  \code{csoilc}  & $4\times10^{-3}$ & --- & Drag coefficient for soil under canopy \\
  \code{wtur}    & 0.1     & m s$^{-1}$   & Minimum turbulent wind speed (stable) \\
  \code{rmax0}   & $2\times10^4$ & s m$^{-1}$ & Maximum stomatal resistance \\
  \code{dewmax}  & 0.1     & mm           & Maximum allowed dew accumulation \\
  \code{trsmx0}  & $2\times10^{-4}$ & kg m$^{-2}$ s$^{-1}$ & Max transpiration rate (saturated soil) \\
  \code{drain}   & $4\times10^{-4}$ & mm s$^{-1}$ & Drainage out of 10\,m soil layer \\
  \code{minwrat} & $10^{-4}$& ---          & Minimum actual/potential soil water ratio \\
\end{longtable}

%----------------------------------------------------------------------
\subsection{Group 11 --- Vapour--Dry Air Ratio Constants}

\begin{longtable}{L{3cm} L{3cm} L{4cm} L{5cm}}
  \toprule
  \textbf{Name} & \textbf{Value} & \textbf{Expression} & \textbf{Description} \\
  \midrule
  \endhead
  \bottomrule
  \endfoot
  \code{ep1}  & $\approx0.6078$ & $M_d/M_w - 1$ & Virtual temperature correction \\
  \code{ep2}  & $\approx0.6220$ & $M_w/M_d$     & Saturation mixing ratio factor \\
  \code{rep2} & $\approx1.6078$ & $M_d/M_w$     & Reciprocal of \code{ep2} \\
\end{longtable}

%----------------------------------------------------------------------
\subsection{Group 12 --- Density Constants}

\begin{longtable}{L{3.5cm} L{3cm} L{3cm} L{4.5cm}}
  \toprule
  \textbf{Name} & \textbf{Value} & \textbf{Units} & \textbf{Description} \\
  \midrule
  \endhead
  \bottomrule
  \endfoot
  \code{rhoh2o}   & 1000.0 & kg m$^{-3}$ & Fresh liquid water \\
  \code{rhosea}   & 1026.0 & kg m$^{-3}$ & Sea water \\
  \code{rhosnow}  & 100.0  & kg m$^{-3}$ & Fresh snow \\
  \code{rhosnowp} & 330.0  & kg m$^{-3}$ & Packed/wind-packed snow \\
  \code{rhoice}   & 917.0  & kg m$^{-3}$ & Pure ice \\
\end{longtable}

%----------------------------------------------------------------------
\subsection{Group 13 --- Microphysics: Terminal Velocity}

Power-law fall speed: $v = A_t\, D^{b_{vt}}$.

\begin{longtable}{L{3cm} L{4cm} L{7cm}}
  \toprule
  \textbf{Name} & \textbf{Value} & \textbf{Description} \\
  \midrule
  \endhead
  \bottomrule
  \endfoot
  \code{avt}  & 841.99667  & Pre-exponential factor \\
  \code{bvt}  & 0.8        & Power-law exponent \\
  \code{g4pb} & 17.837825  & $\Gamma(4+b_{vt})$ \\
  \code{g3pb} & \code{g4pb}$/(3+b_{vt})$ & $\Gamma(3+b_{vt})$ \\
  \code{g5pb} & 1.8273     & $\Gamma(5/2+b_{vt}/2)$ (reflectivity) \\
  \code{vtc}  & \code{avt}$\times$\code{g4pb}$/6$ & Composite terminal velocity constant \\
\end{longtable}

%----------------------------------------------------------------------
\subsection{Group 14 --- Radiative and Emissivity Constants}

\begin{longtable}{L{3.5cm} L{3cm} L{7.5cm}}
  \toprule
  \textbf{Name} & \textbf{Value} & \textbf{Description} \\
  \midrule
  \endhead
  \bottomrule
  \endfoot
  \code{emsw}         & 0.996  & Emissivity of open water \\
  \code{lnd\_sfcemiss}& 0.985  & Land surface longwave emissivity \\
  \code{ocn\_sfcemiss}& 0.996  & Ocean surface longwave emissivity (alias \code{emsw}) \\
  \code{ice\_sfcemiss}& 0.982  & Sea-ice surface longwave emissivity \\
  \code{lowcld}       & 0.01   & Minimum allowed cloud fraction \\
  \code{hicld}        & 0.99   & Maximum allowed cloud fraction \\
\end{longtable}

%----------------------------------------------------------------------
\subsection{Group 15 --- Chemistry and Ocean / Sea-Ice Constants}

\begin{longtable}{L{3.5cm} L{4cm} L{3cm} L{4cm}}
  \toprule
  \textbf{Name} & \textbf{Value} & \textbf{Units} & \textbf{Description} \\
  \midrule
  \endhead
  \bottomrule
  \endfoot
  \code{atmos}   & $2.247\times10^{19}$ & mol cm$^{-3}$ & Atmospheric surface concentration \\
  \code{rtcon}   & $8.314\times10^{-2}$ & L atm mol$^{-1}$ K$^{-1}$ & Henry's law gas constant \\
  \code{dropdif} & $2\times10^{-5}$     & cm$^2$ s$^{-1}$ & Droplet diffusion coefficient \\
  \code{difgas}  & 0.1                  & cm$^2$ s$^{-1}$ & Gas-phase diffusion coefficient \\
  \code{hdmw}    & $1.3889\times10^{-7}$& m$^2$ s$^{-1}$ & Molecular heat diffusion in water \\
  \code{iceminh} & 0.01                 & m               & Minimum sea-ice depth \\
\end{longtable}

%----------------------------------------------------------------------
\subsection{Design Notes}

\textbf{Derived constants.} Wherever possible, compound constants
(\code{rovcp}, \code{gdry}, \code{wlhvocp}, etc.) are expressed as arithmetic
combinations of primitives. Changing a fundamental constant (e.g.\ \code{egrav}
under \code{RCEMIP}) propagates automatically through all derived quantities
at compile time.

\textbf{Triple-precision strategy.} Angle, geometry, and trigonometric constants
are provided in \code{rk4}, \code{rk8}, and \code{rkx} forms to avoid hidden
precision loss in map-projection and solar-orbit code that must remain in double
precision regardless of the model working precision.

\textbf{No procedures.} The module contains no subroutines or functions. All
thermodynamic kernels live in the include files or in other modules, keeping
\code{mod\_constants} a pure compile-time data store with zero runtime cost.

\clearpage

%======================================================================
\section{Include Files --- \texttt{Share/*.inc}}
\label{sec:includes}
%======================================================================

\subsection{Overview}

The \code{Share/} directory contains 20 small, self-contained Fortran functions
that are \emph{textually included} (via \code{include} statements) into multiple
modules across the RegCM codebase. This design allows the same thermodynamic
kernel to be compiled with different precision kinds (\code{rkx}) without code
duplication, and enables aggressive inlining by the compiler
(\code{!DIR\$ ATTRIBUTES FORCEINLINE}).

The files fall into four thematic groups:

\begin{longtable}{L{5cm} L{9cm}}
  \toprule
  \textbf{Group} & \textbf{Files} \\
  \midrule
  \endhead
  \bottomrule
  \endfoot
  Saturation vapour pressure (\code{pfesat}) &
    \code{pfesat.inc}, \code{flatau\_pfesat.inc}, \code{tetens\_pfesat.inc},
    \code{vaisala\_pfesat.inc}, \code{esatliq.inc}, \code{esatice.inc} \\
  Saturation derivatives and specific humidity &
    \code{pfesdt.inc}, \code{flatau\_pfesdt.inc}, \code{tetens\_pfesdt.inc},
    \code{pfqsat.inc}, \code{pfqsdt.inc}, \code{pfwsat.inc}, \code{pfmxsat.inc} \\
  Latent heat and thermodynamic scalars &
    \code{wlh.inc}, \code{cpmf.inc}, \code{clwfromt.inc}, \code{moistlr.inc} \\
  Vertical and surface diagnostics &
    \code{sig2p.inc}, \code{lapserate.inc}, \code{evpt.inc} \\
\end{longtable}

All functions are declared \code{pure elemental} (or \code{pure} for
non-elemental variants), making them safe for OpenACC GPU kernels
(\code{!\$acc routine seq}) and Fortran array contexts.

\subsection{Common Conventions}

\begin{longtable}{L{3cm} L{11cm}}
  \toprule
  \textbf{Symbol} & \textbf{Meaning} \\
  \midrule
  \endhead
  \bottomrule
  \endfoot
  \code{rkx}   & Working-precision real kind (set by the including module) \\
  \code{rk8}   & Double-precision real kind \\
  \code{tzero} & $273.15$\,K --- the ice-point offset \\
  \code{ep2}   & $M_w/M_d \approx 0.622$ \\
  \code{d\_one}, \code{d\_zero} & Double-precision $1.0$ and $0.0$ \\
\end{longtable}

%----------------------------------------------------------------------
\subsection{Saturation Vapour Pressure Functions}

\subsubsection{\texttt{pfesat.inc} --- Flatau 1992 (default, with pressure cap)}

\textbf{Function:} \code{pfesat(t, p) $\to$ es}

\textbf{Purpose:} Computes saturation vapour pressure using 8th-order Flatau
polynomial fits, capped at $0.15\,p$.

\textbf{Arguments:}
\begin{longtable}{L{2cm} L{2cm} L{3.5cm} L{2cm} L{4.5cm}}
  \toprule
  \textbf{Name} & \textbf{Intent} & \textbf{Type} & \textbf{Units} & \textbf{Description} \\
  \midrule
  \endhead
  \bottomrule
  \endfoot
  \code{t} & \code{in} & \code{real(rkx)} & K  & Temperature \\
  \code{p} & \code{in} & \code{real(rkx)} & Pa & Total pressure \\
\end{longtable}

\textbf{Algorithm:} Converts $T$ to Celsius offset
$t_d = \text{clamp}(T - 273.15,\,-75,\,100)$. For $t_d \geq 0$ (liquid),
evaluates the 9-coefficient polynomial $a_0 + t_d(a_1 + t_d(\cdots))$ in Pa.
For $t_d < 0$ (ice), uses the ice-phase polynomial ($c$ coefficients).
Result is capped at $0.15\,p$.

\textbf{Reference:} Flatau, P.J., Walko, R.L., Cotton, W.R. (1992),
\textit{J.\ Appl.\ Meteor.}, 31, 1507--1513.

%----
\subsubsection{\texttt{flatau\_pfesat.inc} --- Flatau 1992 (hPa-based, no pressure cap)}

\textbf{Function:} \code{pfesat(t) $\to$ es}

Same polynomial structure as \code{pfesat.inc} but polynomial constants stored
in hPa; result multiplied by 100 to return Pa. Does not accept a pressure
argument; no $0.15p$ cap applied.

%----
\subsubsection{\texttt{tetens\_pfesat.inc} --- Tetens (1930)}

\textbf{Function:} \code{pfesat(t) $\to$ es}

Classic Tetens exponential formula. Less expensive than the Flatau polynomials.

\[
  e_s =
  \begin{cases}
    c_{1es}\exp\!\left(\dfrac{c_{3les}(T - T_w)}{T - c_{4les}}\right) & T > T_w \\[8pt]
    c_{1es}\exp\!\left(\dfrac{c_{3ies}(T - T_w)}{T - c_{4ies}}\right) & T \leq T_w
  \end{cases}
\]

Constants $c_{1es}$, $c_{3les}$, $c_{4les}$, $c_{3ies}$, $c_{4ies}$ and
$T_w$ (\code{wattp}) are taken from \code{mod\_constants}.

\textbf{Reference:} Tetens, O. (1930), \textit{Zeitschrift f\"ur Geophysik}, 6, 297.

%----
\subsubsection{\texttt{vaisala\_pfesat.inc} --- Wagner \& Pru\ss{} (2002) / IAPWS-95}

\textbf{Function:} \code{pfesat(t) $\to$ es}

High-accuracy formulation valid from $-100$\,°C to $373$\,°C.

\begin{itemize}
  \item For $T > 273.15$\,K (liquid): 6-term reduced-temperature series with
        $T_c = 647.096$\,K, $p_c = 220640$\,hPa.
  \item For $T \leq 273.15$\,K (ice): 2-term series relative to triple point
        $T_n = 273.16$\,K, $p_n = 6.11657$\,hPa.
\end{itemize}

\textbf{Reference:} Wagner, W. \& Pru\ss, A. (2002),
\textit{J.\ Phys.\ Chem.\ Ref.\ Data}, 31, 387--535.

%----
\subsubsection{\texttt{esatliq.inc} --- Ambaum (2020), liquid water}

\textbf{Function:} \code{esatliq(t) $\to$ es}

Exact analytic saturation pressure over liquid water from the Clausius-Clapeyron
equation with temperature-dependent latent heat:
\[
  e_s = e_{s0}\left(\frac{T_0}{T}\right)^{(c_{pl}-c_{pv})/R_v}
        \exp\!\left(\frac{L_{v0}}{R_v T_0} - \frac{L_v(T)}{R_v T}\right)
\]
Parameters: $e_{s0} = 611.655$\,Pa, $T_0 = 273.16$\,K,
$L_{v0} = 2.501\times10^6$\,J\,kg$^{-1}$,
$c_{pl} - c_{pv} = 2320$\,J\,kg$^{-1}$\,K$^{-1}$.

\textbf{Reference:} Ambaum, M.H.P. (2020), \textit{QJRMS},
\href{https://doi.org/10.1002/qj.3899}{doi:10.1002/qj.3899}.

%----
\subsubsection{\texttt{esatice.inc} --- Ambaum (2020), ice}

\textbf{Function:} \code{esatice(t) $\to$ es}

Same Clausius-Clapeyron structure as \code{esatliq.inc} but uses sublimation
latent heat $L_{i0} = 2.835\times10^6$\,J\,kg$^{-1}$ and
$c_{pi} - c_{pv} = 220$\,J\,kg$^{-1}$\,K$^{-1}$.

%----------------------------------------------------------------------
\subsection{Saturation Pressure Derivatives and Specific Humidity}

\subsubsection{\texttt{pfesdt.inc} --- Derivative of $e_s$ w.r.t.\ $T$ (default)}

\textbf{Function:} \code{pfesdt(t) $\to$ esdt}

Returns $\mathrm{d}e_s/\mathrm{d}T$ [Pa\,K$^{-1}$] using Flatau polynomial
fits for the derivative; coefficients in Pa\,K$^{-1}$ ($b$ for liquid,
$d$ for ice).

\subsubsection{\texttt{flatau\_pfesdt.inc} --- Flatau derivative + \texttt{pfqsdt}}

Contains two functions:
\begin{itemize}
  \item \code{pfesdt(t) $\to$ esdt}: derivative of $e_s$ (hPa-based
        coefficients, converted to Pa\,K$^{-1}$).
  \item \code{pfqsdt(t, p[, e][, dedt]) $\to$ qsdt}: derivative of saturation
        specific humidity [K$^{-1}$]:
        \[
          \frac{\mathrm{d}q_s}{\mathrm{d}T}
          = \frac{\mathrm{d}e_s}{\mathrm{d}T}
            \cdot \frac{\varepsilon_2\,p}{(p - 0.378\,e_s)^2}
        \]
        Pre-computed $e_s$ and $\mathrm{d}e_s/\mathrm{d}T$ may be passed
        as optional arguments.
\end{itemize}

\subsubsection{\texttt{tetens\_pfesdt.inc} --- Tetens derivative + \texttt{pfqsdt}}

\begin{itemize}
  \item \code{pfesdt(t) $\to$ esdt}: analytic Tetens derivative,
        $\mathrm{d}e_s/\mathrm{d}T = e_s \cdot c_{3x}(T_0 - c_{4x})/(T-c_{4x})^2$.
  \item \code{pfqsdt(t, p) $\to$ qsdt}: uses internal \code{pfesat} call.
\end{itemize}

\subsubsection{\texttt{pfqsat.inc} --- Saturation specific humidity}

\textbf{Function:} \code{pfqsat(t, p) $\to$ qs}

Saturation specific humidity [kg\,kg$^{-1}$] via embedded Flatau polynomials:
\[
  w_s = \frac{\varepsilon_2\,e_s}{p - e_s}, \quad
  q_s = \frac{w_s}{1 + w_s}
\]

\subsubsection{\texttt{pfqsdt.inc} --- Derivative of $q_s$ (self-contained)}

\textbf{Function:} \code{pfqsdt(t, p) $\to$ qsdt}

Most self-contained variant: embeds all four Flatau polynomial sets (liquid
$e_s$, liquid $\mathrm{d}e_s/\mathrm{d}T$, ice $e_s$, ice
$\mathrm{d}e_s/\mathrm{d}T$) so it has no external dependency on
\code{pfesdt}.

\subsubsection{\texttt{pfwsat.inc} --- Saturation mixing ratio}

\textbf{Function:} \code{pfwsat(t, p) $\to$ ws}

Returns $w_s = \varepsilon_2\,e_s/(p - e_s)$ [kg\,kg$^{-1}$] without the
specific-humidity conversion applied in \code{pfqsat}.

\subsubsection{\texttt{pfmxsat.inc} --- Mixed-phase saturation (subroutine)}

\textbf{Subroutine:} \code{pfmxsat(t, p, ws, Leff)}

Computes saturation mixing ratio and effective latent heat for a mixed
liquid/ice condensate.

\textbf{Arguments:}
\begin{longtable}{L{2cm} L{2cm} L{3.5cm} L{2cm} L{4.5cm}}
  \toprule
  \textbf{Name} & \textbf{Intent} & \textbf{Type} & \textbf{Units} & \textbf{Description} \\
  \midrule
  \endhead
  \bottomrule
  \endfoot
  \code{t}    & \code{in}  & \code{real(rkx)} & K         & Temperature \\
  \code{p}    & \code{in}  & \code{real(rkx)} & Pa        & Pressure \\
  \code{ws}   & \code{out} & \code{real(rkx)} & kg kg$^{-1}$ & Saturation mixing ratio \\
  \code{Leff} & \code{out} & \code{real(rkx)} & J kg$^{-1}$ & Effective latent heat \\
\end{longtable}

\textbf{Phase blending} (Chaboureau \& Bechtold 2002):
\[
  \alpha(T) =
  \begin{cases}
    1 & T \geq T_\text{liq} = 273.15\,\text{K} \\
    \left(\dfrac{T - T_\text{ice}}{T_\text{liq} - T_\text{ice}}\right)^2
      & T_\text{ice} \leq T < T_\text{liq} \\
    0 & T \leq T_\text{ice} = 233.15\,\text{K}
  \end{cases}
\]
\[
  e_{s,\text{eff}} = \alpha\,e_{s,l} + (1-\alpha)\,e_{s,i}, \quad
  L_\text{eff} = \alpha\,L_v^0 + (1-\alpha)\,L_s^0
\]

\textbf{References:} Flatau et al.\ (1992); Chaboureau \& Bechtold (2002),
\textit{J.\ Atmos.\ Sci.}, 59, 2362; Bretherton et al.\ (2004),
\textit{Mon.\ Wea.\ Rev.}, 132.

%----------------------------------------------------------------------
\subsection{Latent Heat and Thermodynamic Scalars}

\subsubsection{\texttt{wlh.inc} --- Latent heat of condensation/sublimation}

\textbf{Function:} \code{wlh(t) $\to$ latent\_heat}

Temperature-dependent latent heat [J\,kg$^{-1}$] (Rogers \& Yau 1989):
\[
  L =
  \begin{cases}
    (2500.79 - 2.36418\,t_c + 1.58927\times10^{-3}\,t_c^2
      - 6.14342\times10^{-5}\,t_c^3)\times10^3 & t_c > 0 \\
    (2834.1 - 0.29\,t_c - 0.004\,t_c^2)\times10^3 & t_c \leq 0
  \end{cases}
\]
where $t_c = T - 273.15$ in \textdegree C.

\textbf{Reference:} Rogers, R.R. \& Yau, M.K. (1989),
\textit{A Short Course in Cloud Physics}, 3rd ed.

\subsubsection{\texttt{cpmf.inc} --- Moist-air heat capacity}

\textbf{Function:} \code{cpmf(qv) $\to$ cp\_moist}

\[
  c_{pm} = c_{pd}(1 + 0.80\,q_v)
\]

Alternative weightings (0.856 and 0.84) are retained as comments, reflecting
the spread in published values of $c_{pv}/c_{pd}$.

\subsubsection{\texttt{clwfromt.inc} --- In-cloud liquid water content}

\textbf{Function:} \code{clwfromt(t) $\to$ clw}

In-cloud liquid water content [g\,m$^{-3}$], clamped to $[0.001,\,0.1]$:
\[
  \text{clw} =
  \begin{cases}
    0.001 & t_c < -50\,\text{\textdegree C} \\
    0.127 + 6.78\times10^{-3}\,t_c + 1.29\times10^{-4}\,t_c^2
         + 8.68\times10^{-7}\,t_c^3 & \text{otherwise}
  \end{cases}
\]

\textbf{Reference:} Lemus, L. et al.\ (1997).

\subsubsection{\texttt{moistlr.inc} --- Moist adiabatic lapse rate}

\textbf{Function:} \code{mlrate(t, q) $\to$ lr}

\[
  \Gamma_m = \frac{g\bigl(1 + L\,q/(R_d\,T)\bigr)}
                  {c_{pd} + L^2\,q/(R_v\,T^2)}
\]

Latent heat $L$ is computed inline using the same polynomial as \code{wlh.inc}
(avoiding a function call in tight loops).

%----------------------------------------------------------------------
\subsection{Vertical and Surface Diagnostic Functions}

\subsubsection{\texttt{sig2p.inc} --- Sigma to pressure coordinate}

\textbf{Function:} \code{sig2p(ps, sigma, ptop) $\to$ p}

\[
  p = \sigma\,(p_s - p_\text{top}) + p_\text{top}
\]

All three arguments and the result must be in Pa (noted explicitly in
the source). This is the standard linear sigma definition used by RegCM's
hydrostatic dynamical core.

\subsubsection{\texttt{lapserate.inc} --- Climatological lapse rate}

\textbf{Function:} \code{lapse\_rate(jday, lat) $\to$ gammac}

Seasonally and latitudinally varying climatological lapse rate [K\,m$^{-1}$]:
\[
  \delta =
  \begin{cases}
    \tfrac{1}{2}\bigl(1 - \cos(2\pi\,d/365)\bigr) & \phi > 30^\circ\text{ (NH)} \\
    \tfrac{1}{2}\bigl(1 + \cos(2\pi\,d/365)\bigr) & \phi < -30^\circ\text{ (SH)} \\
    0.5                                             & |\phi| \leq 30^\circ\text{ (tropics)}
  \end{cases}
\]
\[
  \Gamma_c = 4.5\times10^{-3} + 2.8\times10^{-3}\,\delta \quad[\text{K m}^{-1}]
\]

Range: $4.5\times10^{-3}$ (polar winter) to $7.3\times10^{-3}$\,K\,m$^{-1}$
(summer mid-latitudes).

\subsubsection{\texttt{evpt.inc} --- Potential evapotranspiration (two variants)}

\textbf{Function 1:} \code{evpt(p, t, u, e, es, m, sw, lw) $\to$ ET}
[kg\,m$^{-2}$\,s$^{-1}$]

Penman-Monteith combination equation (Shuttleworth 1993):
\[
  \text{ET} =
  \frac{(\gamma\cdot6.43(1+0.536\,u)\,\Delta e/\lambda) + m\,R_n}
       {(m + \gamma)\cdot86400}
\]
where $\gamma = 0.0016286\,p/\lambda$ is the psychrometric constant,
$\lambda = L_v(T)\times10^{-6}$\,MJ\,kg$^{-1}$, and
$R_n = \max(R_{sw}-R_{lw},\,0)/\lambda$ is net radiation converted to
equivalent evaporation.

\textbf{Function 2:} \code{evpt\_fao(p, t2, u10, e, es, m, sw, lw) $\to$ ET}

ASCE-EWRI standardised reference evapotranspiration (FAO-56 method).
Key differences from \code{evpt}: wind corrected to 2\,m equivalent
($u_2 = \max(u_{10}\times0.7342,\,0.1)$\,m\,s$^{-1}$), radiation in
MJ\,m$^{-2}$\,h$^{-1}$, fixed aerodynamic resistance numerator $37/T_2$.

\textbf{References:} Shuttleworth (1993); ASCE-EWRI (2005).

%----------------------------------------------------------------------
\subsection{Design Notes}

\textbf{Inlining strategy.} The \code{!DIR\$ ATTRIBUTES FORCEINLINE} directive
(Intel Fortran) combined with \code{pure elemental} declarations allows the
compiler to inline these kernels into tight physics loops with zero call
overhead. The \code{!\$acc routine seq} annotation enables the same functions
on GPU devices without modification.

\textbf{Multiple \texttt{pfesat} variants.} The codebase maintains six
alternative implementations of the same \code{pfesat} interface. The including
module selects the desired accuracy/cost trade-off at compile time simply by
choosing the appropriate \code{include} statement.

\textbf{Precision portability.} All internal computations are promoted to
\code{rk8} (double precision) even when \code{rkx} is single precision, to
avoid rounding errors in polynomial evaluations. Results are cast back to
\code{rkx} on return.

\textbf{Pressure cap.} Both \code{pfesat.inc} and related files cap $e_s$ at
$0.15\,p$, preventing the saturation specific humidity from exceeding
$\sim15$\% of the atmospheric column --- a safeguard against numerical
blow-up in very warm and moist conditions.

%==============================================================================
\section{I/O Infrastructure --- \texttt{mod\_stdio.F90} and
         \texttt{mod\_message.F90}}
\label{sec:stdio_message}
%==============================================================================

These two modules form the lowest-level I/O infrastructure of RegCM. They sit
immediately above the kind modules in the dependency hierarchy and are used by
virtually every other module in the model. \code{mod\_stdio} manages Fortran
I/O unit numbers and exposes the three standard streams; \code{mod\_message}
builds on top of it to provide diagnostic output, a portable abort mechanism,
and a lightweight terminal visualisation facility.

\subsection{Dependency Graph}

\begin{lstlisting}
mod_stdio
    |-- mod_intkinds          (ik4)
    |-- iso_fortran_env       (input_unit, output_unit, error_unit)

mod_message
    |-- mod_stdio             (stdout, stderr)
    |-- mod_intkinds          (ik4)
    |-- mod_realkinds         (rkx)
\end{lstlisting}

Both modules use \code{private} as their default accessibility; every public
symbol is explicitly listed.

%------------------------------------------------------------------------------
\subsection{\texttt{mod\_stdio} --- Standard Streams and Unit Management}
\label{subsec:stdio}

\subsubsection{Overview}

\code{mod\_stdio} solves two related problems:

\begin{enumerate}
  \item It exposes the three standard Fortran I/O streams (\code{stdin},
        \code{stdout}, \code{stderr}) as named \code{integer(ik4)} parameters
        derived from the Fortran 2003 intrinsic module \code{iso\_fortran\_env},
        so that all model code refers to them by name rather than by magic
        numbers.
  \item It maintains a registry of Fortran logical unit numbers in the range
        10--99, preventing two parts of the model from accidentally opening
        files on the same unit.
\end{enumerate}

\subsubsection{Public Constants}

\begin{longtable}{L{3cm} L{4cm} L{7cm}}
  \toprule
  \textbf{Name} & \textbf{Value} & \textbf{Description} \\
  \midrule
  \endhead
  \bottomrule
  \endfoot
  \code{stdin}  & \code{input\_unit}  & Standard input stream (Fortran 2003
    \code{iso\_fortran\_env}; typically unit 5 on most platforms). \\
  \code{stdout} & \code{output\_unit} & Standard output stream (typically
    unit 6). \\
  \code{stderr} & \code{error\_unit}  & Standard error stream (typically
    unit 0 or 6 depending on the compiler). \\
\end{longtable}

Binding to \code{iso\_fortran\_env} rather than hardwiring the traditional
values 5, 6, and 0 ensures portability across compilers and operating systems
that may assign different unit numbers to the standard streams.

\subsubsection{Private State}

\begin{longtable}{L{4cm} L{3cm} L{7cm}}
  \toprule
  \textbf{Name} & \textbf{Value} & \textbf{Description} \\
  \midrule
  \endhead
  \bottomrule
  \endfoot
  \code{file\_minunit} & 10 & Lowest unit number managed by the registry. \\
  \code{file\_maxunit} & 99 & Highest unit number managed by the registry. \\
  \code{unit\_tag(10:99)} & \code{logical} array & Per-unit occupancy flag.
    \code{.true.} means the unit has been granted or is known to be open.
    Initialised to \code{.false.}. \\
\end{longtable}

Units 0--9 are left unmanaged: unit 5 (\code{stdin}) and unit 6
(\code{stdout}) are the traditional standard streams, and units 0--9 may be
reserved by the compiler or runtime library.

\subsubsection{Public Procedures}

\paragraph{\texttt{file\_getunit([unitn])} $\to$ \texttt{integer(ik4)}}

Allocates a Fortran logical unit number from the managed pool and returns it.
Returns $-1$ on failure.

\begin{longtable}{L{3cm} L{3cm} L{8cm}}
  \toprule
  \textbf{Argument} & \textbf{Intent} & \textbf{Description} \\
  \midrule
  \endhead
  \bottomrule
  \endfoot
  \code{unitn} & \code{in}, optional & Preferred unit number.
    If present and the unit is free, it is granted. If absent, the routine
    scans from 99 downward for the first available unit. \\
\end{longtable}

\textbf{Algorithm:}
\begin{itemize}
  \item \textbf{Requested unit path} (when \code{unitn} is present): validates
        that \code{unitn} is in $[10, 99]$ and not already tagged. Then calls
        \code{inquire} to check whether some other part of the program has
        opened it outside the registry. If the unit is free, tags it and
        returns it; otherwise returns $-1$.
  \item \textbf{Auto-allocation path} (no \code{unitn}): scans from
        \code{file\_maxunit} down to \code{file\_minunit}. For each unit,
        calls \code{inquire}; if the unit is found open but not tagged (i.e.\
        opened outside the registry), it is tagged retroactively to prevent
        double-use. The first unit that is neither open nor tagged is granted
        and returned.
  \item If no unit is available after the full scan, returns $-1$.
\end{itemize}

The downward scan (from 99 to 10) means higher-numbered units are preferred,
leaving lower-numbered ones available for code that may hard-code unit numbers.

\paragraph{\texttt{file\_freeunit(unitn)}}

Releases a previously allocated unit back to the pool.

\begin{longtable}{L{3cm} L{3cm} L{8cm}}
  \toprule
  \textbf{Argument} & \textbf{Intent} & \textbf{Description} \\
  \midrule
  \endhead
  \bottomrule
  \endfoot
  \code{unitn} & \code{in} & Unit number to release. Must be in the
    managed range. \\
\end{longtable}

If the unit is tagged (\code{unit\_tag(unitn) == .true.}), the routine closes
the unit and clears the tag. If the unit was not tagged, the call is a no-op.
This ensures that \code{close} is always called exactly once per
\code{file\_getunit} call, even if the caller forgot to close the file
explicitly.

\subsubsection{Usage Pattern}

\begin{lstlisting}
use mod_stdio, only : file_getunit, file_freeunit, stdout, stderr

integer(ik4) :: iunit

iunit = file_getunit()
if ( iunit < 0 ) then
  write(stderr,*) 'No available unit'
  stop
end if
open(iunit, file='output.dat', status='replace')
write(iunit, '(f12.6)') result
call file_freeunit(iunit)   ! closes and releases
\end{lstlisting}

%------------------------------------------------------------------------------
\subsection{\texttt{mod\_message} --- Diagnostics, Abort, and Visualisation}
\label{subsec:message}

\subsubsection{Overview}

\code{mod\_message} provides:
\begin{itemize}
  \item A shared character buffer \code{aline} into which callers format
        messages, and a family of output subroutines (\code{say}, \code{note},
        \code{cry}) that write it to the appropriate stream.
  \item An overloaded \code{die} interface for fatal error termination,
        wrapping an external \code{myabort} routine that is linked separately
        (typically to an MPI abort or a simple \code{stop}).
  \item A \code{fatal} subroutine for compiler-macro-style source location
        reporting.
  \item A \code{checkalloc} subroutine for checking \code{allocate} return
        codes in debug builds.
  \item A self-contained ANSI terminal text-mode visualisation subsystem
        (\code{viz\_*}) for quick inline inspection of 2D arrays during
        development.
\end{itemize}

\subsubsection{Private Module State}

\begin{longtable}{L{4cm} L{3cm} L{7cm}}
  \toprule
  \textbf{Name} & \textbf{Type/Value} & \textbf{Description} \\
  \midrule
  \endhead
  \bottomrule
  \endfoot
  \code{aline}    & \code{character(len=512)} &
    Shared message buffer. Callers write into this with \code{write(aline,*)}
    before calling \code{say}, \code{note}, or \code{cry}. Declared
    \code{public} so that calling modules can format it directly. \\
  \code{cline}    & \code{character(len=8)} &
    Scratch buffer for converting line numbers to strings in \code{fatal}
    and \code{checkalloc}. \\
  \code{esc}      & \code{character} &
    ASCII escape character (\code{achar(27)}), used to build ANSI terminal
    escape sequences in the visualisation routines. \\
  \code{iprank}   & \code{integer(ik4)} &
    MPI rank of the current process. Defaults to 0 (serial or rank-0 in
    parallel). Set by \code{setup\_mesg}. Controls whether messages are
    printed: most output is suppressed on non-zero ranks. \\
  \code{maxplot}  & 10 & Number of distinct printable levels in
    \code{viz\_plot}. \\
  \code{plot(0:10)} & \code{character} array &
    Ordered ASCII characters used as density symbols in the text-mode plot:
    \texttt{[` `, `.`, `,`, `:`, `=`, `+`, `o`, `x`, `X`, `\#`, `@`]}. \\
  \code{rows}     & 24 & Default terminal height for \code{viz\_plot}. \\
  \code{cols}     & 48 & Default terminal width for \code{viz\_plot}. \\
  \code{mrows}, \code{mcols} & & Effective dimensions (possibly clamped to
    array size in \code{viz\_plot}). \\
\end{longtable}

\subsubsection{MPI Rank Registration}

\paragraph{\texttt{setup\_mesg(ipid)}}

\begin{longtable}{L{3cm} L{3cm} L{8cm}}
  \toprule
  \textbf{Argument} & \textbf{Intent} & \textbf{Description} \\
  \midrule
  \endhead
  \bottomrule
  \endfoot
  \code{ipid} & \code{in} & MPI process rank. \\
\end{longtable}

Must be called once at model startup, after \code{MPI\_Init}, with the
process rank. Sets the module-level \code{iprank} variable, which is then
consulted by \code{say}, \code{cry}, and \code{die} to suppress duplicate
output from non-zero ranks. In serial runs, the default value of 0 means
all output is emitted.

\subsubsection{Message Output Subroutines}

\paragraph{\texttt{say}}

Writes \code{trim(aline)} to \code{stdout}, but \emph{only} on rank 0.
Intended for informational messages that should appear exactly once in the
output log regardless of MPI decomposition.

\paragraph{\texttt{note}}

Writes a notice banner to \code{stderr} from \emph{every} rank, prefixed
with the rank number. Used for per-process status messages during
initialisation or when each process must report independently.
Overwrites \code{aline} with the banner text.

\paragraph{\texttt{cry}}

Writes an ``important notice'' banner to \code{stderr}, but only on rank 0.
Used for warnings that are more serious than \code{say} but do not abort the
run.

\subsubsection{Error and Abort Routines}

\paragraph{\texttt{die} (generic interface)}

The primary fatal-error mechanism. Writes a diagnostic to \code{stderr} and
calls \code{myabort}. Four overloaded variants are provided, allowing
increasingly detailed error messages:

\begin{longtable}{L{4.5cm} L{9.5cm}}
  \toprule
  \textbf{Signature} & \textbf{Description} \\
  \midrule
  \endhead
  \bottomrule
  \endfoot
  \code{die(msg)} &
    Writes a single message string to \code{stderr} and aborts. \\
  \code{die(msg, msg1)} &
    Writes \texttt{msg : msg1} (two strings). Used throughout the codebase
    as \code{die('ModuleName', 'Descriptive error text')}. \\
  \code{die(msg, msg1, ier1)} &
    Appends an integer error code: \texttt{msg : msg1 : ier1}. \\
  \code{die(msg, msg1, ier1, msg2, ier2)} &
    Two string/integer pairs: \texttt{msg : msg1 : ier1 : msg2 : ier2}.
    Used when two related error codes (e.g.\ two NetCDF status codes) must
    be reported simultaneously. \\
\end{longtable}

All four variants suppress output on non-zero MPI ranks (\code{iprank /= 0})
and call \code{myabort} unconditionally.

\paragraph{\texttt{myabort} (external interface)}

\begin{lstlisting}
interface myabort
  subroutine myabort( )
    implicit none
  end subroutine myabort
end interface myabort
\end{lstlisting}

\code{myabort} is declared as an external interface but not implemented in
this module. It must be provided by the linking application. This design
allows the abort strategy to be chosen at link time:

\begin{itemize}
  \item In \textbf{MPI} builds, \code{myabort} calls \code{MPI\_Abort} to
        cleanly tear down all ranks rather than leaving them hanging.
  \item In \textbf{serial} builds, \code{myabort} contains a bare
        \code{stop} or \code{error stop} statement.
\end{itemize}

This separation keeps \code{mod\_message} free of any MPI dependency,
which is important because it is used in pre-processing tools that do not
link against MPI.

\paragraph{\texttt{fatal(filename, line, str)}}

A higher-level fatal error reporter that includes source file location.

\begin{longtable}{L{3cm} L{3cm} L{8cm}}
  \toprule
  \textbf{Argument} & \textbf{Intent} & \textbf{Description} \\
  \midrule
  \endhead
  \bottomrule
  \endfoot
  \code{filename} & \code{in} & Source file name (typically \code{\_\_FILE\_\_}). \\
  \code{line}     & \code{in} & Source line number (typically \code{\_\_LINE\_\_}). \\
  \code{str}      & \code{in} & Error message string. \\
\end{longtable}

Formats and writes a ``FATAL CALLED'' banner, the file and line location, and
the error message to \code{stderr}, then delegates to \code{die} with the
string representations of the arguments. Intended to be invoked via a
preprocessor macro:

\begin{lstlisting}
#define FATAL(msg) call fatal(__FILE__,__LINE__,msg)
\end{lstlisting}

\paragraph{\texttt{checkalloc(ival, filename, line, arg)}}

Checks the status code \code{ival} returned by a Fortran \code{allocate}
statement and calls \code{die} if it is non-zero.

\begin{longtable}{L{3cm} L{3cm} L{8cm}}
  \toprule
  \textbf{Argument} & \textbf{Intent} & \textbf{Description} \\
  \midrule
  \endhead
  \bottomrule
  \endfoot
  \code{ival}     & \code{in} & \code{allocate} \code{stat=} return value. \\
  \code{filename} & \code{in} & Source file name. \\
  \code{line}     & \code{in} & Source line number. \\
  \code{arg}      & \code{in} & Name of the array being allocated (for
    the error message). \\
\end{longtable}

\textbf{Debug mode only.} When compiled \emph{without} the \code{DEBUG}
preprocessor flag, \code{checkalloc} is declared \code{pure} and compiles
to a no-op \code{return}. When compiled \emph{with} \code{DEBUG}, it writes
a detailed error message and aborts. This means allocation checks carry zero
runtime cost in production builds.

\subsubsection{ANSI Terminal Visualisation (\texttt{viz\_*})}

A self-contained text-mode visualisation subsystem for interactive inspection
of 2D real arrays during development or debugging. It uses ANSI/VT100 escape
sequences and is therefore only useful on terminals that support them.

\paragraph{Attribution}

The \code{viz\_*} routines are adapted from a package by
Axel Kohlmeyer (\texttt{akohlmey@gmail.com}), Copyright \copyright{} 2013,
as noted in the source comments.

\paragraph{\texttt{viz\_init([x, y])}}

Initialises the visualisation subsystem and optionally sets the plot
dimensions.

\begin{longtable}{L{3cm} L{3cm} L{8cm}}
  \toprule
  \textbf{Argument} & \textbf{Intent} & \textbf{Description} \\
  \midrule
  \endhead
  \bottomrule
  \endfoot
  \code{x} & \code{in}, optional & Number of columns (default: 48). \\
  \code{y} & \code{in}, optional & Number of rows (default: 24). \\
\end{longtable}

Hides the terminal cursor to reduce flicker during animation by emitting
the escape sequence \texttt{ESC[?25l}.

\paragraph{\texttt{viz\_clear}}

Clears the terminal screen by emitting \texttt{ESC[2J}.

\paragraph{\texttt{viz\_plot(val [, rmax])}}

Renders a 2D real array as a coloured ASCII density plot on the terminal.

\begin{longtable}{L{3cm} L{3cm} L{8cm}}
  \toprule
  \textbf{Argument} & \textbf{Intent} & \textbf{Description} \\
  \midrule
  \endhead
  \bottomrule
  \endfoot
  \code{val(:,:)} & \code{in} & 2D \code{real(rkx)} array to plot. \\
  \code{rmax}     & \code{in}, optional & Maximum absolute value for
    scaling. If absent, computed from the data. \\
\end{longtable}

\textbf{Algorithm:}
\begin{enumerate}
  \item Clamps the plot area to \code{min(cols, nx)} $\times$
        \code{min(rows, ny)}.
  \item Computes a block-averaging stride \code{dx = nx/mcols},
        \code{dy = ny/mrows} so that the full array is summarised
        into the available character cells.
  \item Finds the maximum absolute value in the averaged data (unless
        \code{rmax} is supplied) and scales to the 11-level character
        palette \code{plot(0:10)}.
  \item Renders rows from top (high index) to bottom (low index),
        colouring negative values blue/cyan and positive values
        red/black using ANSI foreground colour codes:
        \texttt{ESC[34m} (blue), \texttt{ESC[36m} (cyan),
        \texttt{ESC[30m} (black), \texttt{ESC[31m} (red).
  \item Appends a \texttt{|} gutter character at the right edge of each row.
\end{enumerate}

\paragraph{\texttt{viz\_done}}

Restores the terminal to a clean state: positions the cursor at the last row,
resets the foreground colour, re-enables the cursor (\texttt{ESC[?25h}),
and emits a full terminal reset (\texttt{ESC[0m}).

%------------------------------------------------------------------------------
\subsection{Design Notes}
\label{subsec:stdio_message_design}

\textbf{Separation of concerns between the two modules.}
\code{mod\_stdio} is intentionally minimal: it manages units and exposes
stream names, nothing more. This allows it to be used even in the smallest
utility programs without pulling in the diagnostic machinery of
\code{mod\_message}. The split also means \code{mod\_stdio} itself never
calls \code{die}, avoiding a circular dependency.

\textbf{MPI-aware output without an MPI dependency.}
By using the module-level \code{iprank} variable (defaulting to 0), all
\code{mod\_message} routines behave correctly in both serial and parallel
contexts. The MPI rank is injected once via \code{setup\_mesg} rather than
being passed as an argument to every call, keeping call sites clean. No MPI
header or library is referenced inside the module.

\textbf{Link-time abort strategy via \texttt{myabort}.}
Declaring \code{myabort} as an external interface (a Fortran \emph{interface
block} without a body) defers its implementation to the link step.
This is the Fortran equivalent of a C weak symbol or a virtual function:
any object file providing a subroutine named \code{myabort} with no arguments
will satisfy the reference. The RegCM build system provides two implementations
--- one for MPI builds and one for serial builds --- selected by the Makefile
target.

\textbf{Zero-cost allocation checks.}
The \code{pure}/non-\code{pure} duality of \code{checkalloc} controlled by
the \code{DEBUG} flag is a classical Fortran technique for assertions that
disappear entirely in optimised builds. Because \code{pure} procedures cannot
have side effects (no I/O, no \code{stop}), the compiler is guaranteed to
eliminate the no-op body; no branch or call instruction reaches the binary.

\textbf{The \texttt{aline} buffer pattern.}
Rather than requiring each message-printing call to accept a
\code{character(*)} argument (which complicates interfaces and forces callers
to declare local strings), the module exposes the public \code{aline} buffer.
Callers write into it with \code{write(aline,'(...)')} and then call
\code{say} or \code{cry}. This trades thread safety (the buffer is a shared
module variable) for simplicity of the calling interface, which is acceptable
in a model that serialises all diagnostic output through rank 0.

%==============================================================================
\section{Date and Time Arithmetic --- \texttt{mod\_date.F90}}
\label{sec:moddate}
%==============================================================================

\subsection{Overview}

\code{mod\_date} provides the complete calendar and time-management layer for
RegCM. It defines the model's two core derived types --- \code{rcm\_time\_and\_date}
and \code{rcm\_time\_interval} --- together with a rich set of overloaded
operators, conversion utilities, calendar arithmetic, and NetCDF-compatible
time-string parsers.

The module is \code{private} by default; every symbol it exposes is listed
explicitly in a \code{public} statement, making the public interface unambiguous.
It supports three calendar systems (proleptic Gregorian, no-leap, and 360-day),
all selectable at runtime without recompilation.

\textbf{Key capabilities:}
\begin{itemize}
  \item Portable internal representation of date and time as a pair of 64-bit
        integers, independent of calendar system.
  \item Overloaded Fortran operators (\code{=}, \code{+}, \code{-},
        \code{==}, \code{/=}, \code{<}, \code{>}, \code{<=}, \code{>=})
        for natural-looking date and interval arithmetic in physics code.
  \item Conversion to/from multiple external formats: packed 10-digit integer
        (\code{YYYYMMDDHH}), ISO~8601 string, human-readable UTC string, hours,
        and seconds.
  \item Month, week, and day navigation helpers (\code{nextmon}, \code{prevmon},
        \code{monfirst}, \code{monlast}, \code{monmiddle}, \code{nextwk},
        \code{yrfirst}, etc.).
  \item Julian Day computation in the proleptic Gregorian calendar, valid
        for any year, used for the day of the week.
  \item Parsing of CF-Conventions time-unit strings
        (\texttt{hours since YYYY-MM-DD HH:MM:SS}).
  \item CMIP scenario boundary detection (\code{date\_in\_scenario}).
\end{itemize}

\subsection{Dependencies}

\begin{lstlisting}
mod_date
    |-- mod_realkinds   (rk4, rk8, rkx)
    |-- mod_intkinds    (ik4, ik8)
    |-- mod_stdio       (stdout, stderr)
    |-- mod_message     (die)
\end{lstlisting}

\code{mod\_stdio} supplies the unit numbers \code{stdout} and \code{stderr}
used in error reporting. \code{mod\_message} supplies the \code{die} subroutine
that writes a fatal diagnostic and terminates the run.

%------------------------------------------------------------------------------
\subsection{Named Constants}
\label{subsec:date_constants}

\subsubsection{Module-Private Time Conversion Factors}

All are \code{integer(ik8)} to prevent overflow in interval calculations.

\begin{longtable}{L{3cm} L{3cm} L{3cm} L{5cm}}
  \toprule
  \textbf{Name} & \textbf{Value} & \textbf{Units} & \textbf{Description} \\
  \midrule
  \endhead
  \bottomrule
  \endfoot
  \code{i8spw} & 604\,800  & s week$^{-1}$ & Seconds per week \\
  \code{i8spd} & 86\,400   & s d$^{-1}$    & Seconds per day \\
  \code{i8mpd} & 1\,440    & min d$^{-1}$  & Minutes per day \\
  \code{i8hpd} & 24        & h d$^{-1}$    & Hours per day \\
  \code{i8sph} & 3\,600    & s h$^{-1}$    & Seconds per hour \\
  \code{i8spm} & 60        & s min$^{-1}$  & Seconds per minute \\
  \code{i8mph} & 60        & min h$^{-1}$  & Minutes per hour \\
  \code{i8mpy} & 12        & month yr$^{-1}$& Months per year \\
  \code{i8mpc} & 1\,200    & month century$^{-1}$ & Months per century \\
  \code{i8ypc} & 100       & yr century$^{-1}$ & Years per century \\
\end{longtable}

\subsubsection{Public Calendar Selectors}

\begin{longtable}{L{3cm} L{2cm} L{9cm}}
  \toprule
  \textbf{Name} & \textbf{Value} & \textbf{Description} \\
  \midrule
  \endhead
  \bottomrule
  \endfoot
  \code{gregorian} & 1 & Proleptic Gregorian calendar (leap years via the
    400/4/100 rule), with no Julian/Gregorian switch in 1582. \\
  \code{noleap}    & 2 & 365-day calendar with no leap years. \\
  \code{y360}      & 3 & Idealised 360-day calendar (12 months of exactly 30
    days each). \\
\end{longtable}

The calendar names are recognised by three routines, which do not accept
exactly the same strings:

\begin{longtable}{L{3.2cm} L{3.6cm} L{3.6cm} L{3.6cm}}
  \toprule
  \textbf{Calendar} & \code{calendar\_int} & \code{set\_calstring} &
    \code{timeval2date} \\
  \midrule
  \endhead
  \bottomrule
  \endfoot
  \code{gregorian} &
    \texttt{gregorian}, \texttt{standard}, \texttt{proleptic\_gregorian} &
    \texttt{gregorian}, \texttt{standard} &
    any string not listed below \\
  \code{noleap} &
    \texttt{noleap}, \texttt{365\_days}, \texttt{365\_day},
    \texttt{days\_365}, \texttt{day\_365} &
    \texttt{noleap}, \texttt{days\_365}, \texttt{day\_365},
    \texttt{365\_days}, \texttt{365\_day} &
    strings starting with \texttt{noleap}, \texttt{days\_365} or
    \texttt{365\_day} \\
  \code{y360} &
    \texttt{days\_360}, \texttt{day\_360}, \texttt{360\_days},
    \texttt{360\_day} &
    \texttt{days\_360}, \texttt{day\_360}, \texttt{360\_days},
    \texttt{360\_day} &
    strings starting with \texttt{360\_day} \\
  Unknown string &
    \code{gregorian}, silently &
    \code{gregorian}, with a warning on \code{stderr} &
    \code{gregorian}, silently \\
\end{longtable}

\noindent In particular \texttt{proleptic\_gregorian}, the name written by the
model in its own output files, is accepted by \code{set\_calstring} only
through the fallback (with the warning), and \texttt{days\_360},
\texttt{day\_360} and \texttt{day\_365} are read as Gregorian by
\code{timeval2date}.

\subsubsection{Public Time Unit Selectors}

Used in \code{rcm\_time\_interval\%iunit} and the \code{set\_timeunit} routine.

\begin{longtable}{L{3cm} L{2cm} L{9cm}}
  \toprule
  \textbf{Name} & \textbf{Value} & \textbf{Description} \\
  \midrule
  \endhead
  \bottomrule
  \endfoot
  \code{usec} & 1 & Seconds \\
  \code{umin} & 2 & Minutes \\
  \code{uhrs} & 3 & Hours \\
  \code{uday} & 4 & Days \\
  \code{umnt} & 5 & Months \\
  \code{uyrs} & 6 & Years \\
  \code{ucnt} & 7 & Centuries \\
\end{longtable}

\subsubsection{Public Season Selectors}

\begin{longtable}{L{3.5cm} L{2cm} L{8.5cm}}
  \toprule
  \textbf{Name} & \textbf{Value} & \textbf{Description} \\
  \midrule
  \endhead
  \bottomrule
  \endfoot
  \code{isea\_djf} & 1 & December--January--February (boreal winter) \\
  \code{isea\_mam} & 2 & March--April--May (boreal spring) \\
  \code{isea\_jja} & 3 & June--July--August (boreal summer) \\
  \code{isea\_son} & 4 & September--October--November (boreal autumn) \\
\end{longtable}

\subsubsection{Reference and Auxiliary Constants}

\begin{longtable}{L{4cm} L{4cm} L{7cm}}
  \toprule
  \textbf{Name} & \textbf{Value} & \textbf{Description} \\
  \midrule
  \endhead
  \bottomrule
  \endfoot
  \code{reference\_year}  & 1900          & Epoch year for the internal
    days-from-reference counter. \\
  \code{cordex\_refdate}  & 1950010100    & Packed date (1 Jan 1950, 00:00 UTC)
    used as origin for CORDEX time output (\code{curr\_time}). \\
  \code{mlen(1:12)}       & [31,28,31,\ldots] & Standard (non-leap) month
    lengths. \\
  \code{cintstr(1:7)}     & string array  & Human-readable labels for the seven
    time-unit selectors (\texttt{seconds}, \texttt{minutes}, \ldots,
    \texttt{centuries}). \\
  \code{calstr(1:3)}      & string array  & CF-Conventions calendar names
    (\texttt{proleptic\_gregorian}, \texttt{noleap}, \texttt{360\_day}). \\
\end{longtable}

%------------------------------------------------------------------------------
\subsection{Derived Types}
\label{subsec:date_types}

\subsubsection{\texttt{rcm\_time\_and\_date} --- The Primary Timestamp}

\begin{lstlisting}
type rcm_time_and_date
  integer(ik4) :: calendar            = gregorian
  integer(ik8) :: days_from_reference = 0
  integer(ik8) :: second_of_day       = 0
end type rcm_time_and_date
\end{lstlisting}

This is the model's universal timestamp. Every date in RegCM --- simulation
start, current model time, output trigger, boundary-condition timestamp ---
is stored as a value of this type.

\begin{longtable}{L{4cm} L{3cm} L{7cm}}
  \toprule
  \textbf{Field} & \textbf{Type} & \textbf{Description} \\
  \midrule
  \endhead
  \bottomrule
  \endfoot
  \code{calendar}            & \code{integer(ik4)} &
    Calendar system in use (\code{gregorian}, \code{noleap}, or \code{y360}).
    Default: \code{gregorian}. \\
  \code{days\_from\_reference} & \code{integer(ik8)} &
    Signed count of days elapsed since 1 January of \code{reference\_year}
    (1900). Negative values represent dates before 1900. The use of \code{ik8}
    avoids overflow for multi-millennium simulations. \\
  \code{second\_of\_day}      & \code{integer(ik8)} &
    Seconds elapsed since midnight of the current day, in the range
    $[0,\,86399]$. Storing time separately from the day count makes
    sub-daily arithmetic exact and free of floating-point rounding. \\
\end{longtable}

\textbf{Internal representation rationale.} Splitting the timestamp into
integer days and integer seconds rather than using a single real-valued
counter (as in, e.g., Julian Day fractions) avoids accumulated rounding errors
over long simulations with short time steps. The 64-bit day counter alone
could span about $\pm2.5\times10^{16}$ years. The practical range is much
smaller: the decomposed year is an \code{integer(ik4)}, the packed-integer
assignments accept only years with $|y| < 10000$, and the conversion between
the day counter and the calendar date iterates year by year from 1900, so its
cost grows linearly with the distance from the reference year.

\subsubsection{\texttt{rcm\_time\_interval} --- A Duration}

\begin{lstlisting}
type rcm_time_interval
  integer(ik8) :: ival  = 0
  integer(ik4) :: iunit = usec
end type rcm_time_interval
\end{lstlisting}

Represents a signed duration in a specified unit. The unit is carried alongside
the magnitude so that unit-aware arithmetic (e.g.\ adding 1 month vs.\ adding
30 days) is unambiguous.

\begin{longtable}{L{3cm} L{3cm} L{8cm}}
  \toprule
  \textbf{Field} & \textbf{Type} & \textbf{Description} \\
  \midrule
  \endhead
  \bottomrule
  \endfoot
  \code{ival}  & \code{integer(ik8)} &
    Magnitude of the interval in \code{iunit} units. May be negative
    (intervals are signed). \\
  \code{iunit} & \code{integer(ik4)} &
    Unit selector: one of \code{usec}, \code{umin}, \code{uhrs},
    \code{uday}, \code{umnt}, \code{uyrs}, \code{ucnt}.
    Default: \code{usec}. \\
\end{longtable}

\textbf{Note on calendar-dependent units.} Conversions between calendar-independent
units (\code{usec}--\code{uday}) are exact and handled by the module's
\code{isequalit} and comparison functions. Conversions involving \code{umnt},
\code{uyrs}, and \code{ucnt} depend on the calendar; the comparison functions
correctly return \code{.false.} when mixing these with day-based units, and
\code{die} is called in \code{tohours}/\code{toseconds} if such a conversion
is attempted, because the result is calendar-dependent and potentially
ambiguous.

\subsubsection{Auxiliary Input/Output Types}

Three lightweight types are provided to carry a value \emph{together with} a
calendar selector, used as intermediate representations in the overloaded
assignment interface.

\begin{longtable}{L{3cm} L{3cm} L{8cm}}
  \toprule
  \textbf{Type} & \textbf{Fields} & \textbf{Description} \\
  \midrule
  \endhead
  \bottomrule
  \endfoot
  \code{i4wcal}  & \code{i}: \code{integer(ik4)}, \code{cal}: \code{integer(ik4)} &
    Packed 10-digit date integer (32-bit) with calendar. \\
  \code{i8wcal}  & \code{i}: \code{integer(ik8)}, \code{cal}: \code{integer(ik4)} &
    Packed 10-digit date integer (64-bit) with calendar. \\
  \code{c10wcal} & \code{c}: \code{character(len=11)}, \code{cal}: \code{integer(ik4)} &
    Date as 11-character string with calendar. \\
\end{longtable}

\subsubsection{Internal Date/Time Types}

Two further types are used purely inside the module as a human-readable
decomposed representation; they are not public.

\begin{longtable}{L{3cm} L{8cm}}
  \toprule
  \textbf{Type} & \textbf{Fields and purpose} \\
  \midrule
  \endhead
  \bottomrule
  \endfoot
  \code{iadate} &
    \code{calendar} (\code{ik4}), \code{year} (\code{ik4}),
    \code{month} (\code{ik4}), \code{day} (\code{ik4}).
    Decomposed calendar date. Default values: Gregorian, year 1900,
    January 1. \\
  \code{iatime} &
    \code{hour} (\code{ik4}), \code{minute} (\code{ik4}),
    \code{second} (\code{ik4}).
    Time of day in hours, minutes, seconds. Default: all zero. \\
\end{longtable}

%------------------------------------------------------------------------------
\subsection{Overloaded Operators}
\label{subsec:date_operators}

All date arithmetic in model code is expressed through Fortran operator syntax.
The module implements the following generic interfaces.

\subsubsection{Assignment (\texttt{=})}

Converts any of six source representations into an \code{rcm\_time\_and\_date},
and any of four source representations into an \code{rcm\_time\_interval},
by dispatching to the appropriate \code{initfrom*} procedure.

\begin{longtable}{L{5cm} L{9cm}}
  \toprule
  \textbf{Source type} & \textbf{Target and behaviour} \\
  \midrule
  \endhead
  \bottomrule
  \endfoot
  \code{character} (via \code{c10wcal})  &
    Reads a packed date string + calendar into \code{rcm\_time\_and\_date}. \\
  \code{integer(ik4)} &
    Interprets as packed \code{YYYYMMDDHH}; assumes Gregorian calendar. \\
  \code{i4wcal} &
    As above but with explicit calendar. \\
  \code{integer(ik8)} &
    64-bit packed date; assumes Gregorian calendar. \\
  \code{i8wcal} &
    64-bit packed date with explicit calendar. \\
  \code{rcm\_time\_and\_date} &
    Deep copy. \\
  \code{integer(ik4)} $\to$ \code{rcm\_time\_interval} &
    Sets \code{ival} to the integer, \code{iunit} to \code{usec}. \\
  \code{real(rk4)} $\to$ \code{rcm\_time\_interval} &
    Truncates to integer seconds. \\
  \code{real(rk8)} $\to$ \code{rcm\_time\_interval} &
    Truncates to integer seconds. \\
  \code{rcm\_time\_interval} &
    Deep copy of interval. \\
\end{longtable}

\textbf{Packed date format.} The integer encoding is
$\text{date} = \text{YYYY} \times 10^6 + \text{MM} \times 10^4
             + \text{DD} \times 10^2 + \text{HH}$,
which is the legacy RegCM date format used throughout the namelist and
restart files. The \code{split\_i10} / \code{split\_i10\_8} routines decode
this by successive integer division and modulo operations. All the
packed-date assignments check the decoded fields and call \code{die} if the
year is not in $(-10000, 10000)$, the month not in $[1,12]$, the day not in
$[1,31]$, or the hour not in $[0,23]$ (the day is not checked against the
length of the month). The character form is read with format \code{(i11)}
and then treated as an \code{i8wcal}.

\subsubsection{Equality (\texttt{==})}

\begin{longtable}{L{5cm} L{9cm}}
  \toprule
  \textbf{Operands} & \textbf{Behaviour} \\
  \midrule
  \endhead
  \bottomrule
  \endfoot
  \code{rcm\_time\_and\_date} == \code{rcm\_time\_and\_date} &
    Checks calendar compatibility (calls \code{die} if calendars differ), then
    compares \code{days\_from\_reference} and \code{second\_of\_day}. \\
  \code{rcm\_time\_and\_date} == \code{integer(ik4)} &
    Converts the integer to \code{rcm\_time\_and\_date} (in the lhs calendar),
    then delegates to the type-type comparison. \\
  \code{rcm\_time\_interval} == \code{rcm\_time\_interval} &
    Unit-aware comparison. Converts both operands to a common unit before
    comparing; returns \code{.false.} when units are incompatible
    (e.g.\ months vs.\ seconds). \\
\end{longtable}

The interval comparisons (\code{==}, \code{>}, \code{<}) share the private
subroutine \code{common\_unit(x, y, ix, iy, lok)}. It returns the values of
the two intervals in a common unit: unchanged if the units are the same,
otherwise in seconds if both units are between \code{usec} and \code{uday},
or in months if both are between \code{umnt} and \code{ucnt}. The
conversion factors come from \code{unit\_factor}. The flag \code{lok} is
\code{.false.} when one unit is in each group, or when a unit is unknown,
and then every comparison returns \code{.false.}. Converting to the finer
unit keeps the comparison exact: for example 61~seconds is not equal to
1~minute.

\subsubsection{Addition and Subtraction (\texttt{+}, \texttt{-})}

\begin{longtable}{L{5.5cm} L{8.5cm}}
  \toprule
  \textbf{Expression} & \textbf{Result and behaviour} \\
  \midrule
  \endhead
  \bottomrule
  \endfoot
  \code{rcm\_time\_and\_date + rcm\_time\_interval} &
    Returns a new \code{rcm\_time\_and\_date} advanced by the interval.
    Handles all seven unit cases; seconds, minutes, hours, and days update
    \code{second\_of\_day} and \code{days\_from\_reference} by integer
    arithmetic. Month, year, and century increments first decompose the date,
    adjust the year/month, then recompose. A negative number of months
    goes back across the year boundary. For months only, the position
    within the month, $(d-1)/N_\text{month}$, is preserved and mapped onto
    the target month (see Design Notes); year and century increments keep
    the day of the month unchanged (29 February goes to 1 March in a
    common year). \\
  \code{rcm\_time\_and\_date - rcm\_time\_interval} &
    Subtraction. For seconds, minutes, hours and days it is the addition
    of the opposite interval. For months it keeps the day of the month, or
    takes the last day of the target month if that is shorter (31 March
    $-$ 1 month is 28 or 29 February). Unlike the addition it does not
    preserve the position within the month, so \code{x + 1 month - 1
    month} is not always \code{x}. \\
  \code{rcm\_time\_and\_date - rcm\_time\_and\_date} &
    Returns an \code{rcm\_time\_interval} in \code{usec}, equal to the
    signed difference $(x_\text{days} - y_\text{days}) \times 86400 +
    (x_\text{sod} - y_\text{sod})$. Calls \code{die} if calendars differ. \\
  \code{rcm\_time\_interval + rcm\_time\_interval} &
    Returns a new interval. If units match, adds magnitudes directly;
    otherwise converts both to seconds and returns a \code{usec} interval. \\
\end{longtable}

\subsubsection{Ordering (\texttt{>}, \texttt{<}, \texttt{>=}, \texttt{<=}, \texttt{/=})}

The five relational operators are overloaded for \code{rcm\_time\_and\_date},
both between two dates and between a date and a packed integer date
(\code{integer(ik4)}; \code{<} also accepts \code{integer(ik8)}). The integer
is converted to a date in the calendar of the left operand. Only \code{>} and
\code{<} are also defined between two \code{rcm\_time\_interval} values.

Date comparisons first check calendar consistency via \code{check\_cal} and
then compare the pair (\code{days\_from\_reference}, \code{second\_of\_day})
lexicographically. Interval comparisons are unit-aware, using the same
conversion logic as equality, and return \code{.false.} for incompatible
units.

%------------------------------------------------------------------------------
\subsection{Calendar Arithmetic Routines}
\label{subsec:date_calendar}

\subsubsection{Leap Year and Month Length}

\begin{longtable}{L{4.5cm} L{4cm} L{6cm}}
  \toprule
  \textbf{Routine} & \textbf{Signature} & \textbf{Description} \\
  \midrule
  \endhead
  \bottomrule
  \endfoot
  \code{lleap} & \code{(iyear) $\to$ logical} &
    Private. Returns \code{.true.} if \code{iyear} is a Gregorian leap year
    ($\text{mod}(y,400)=0$, or $\text{mod}(y,4)=0$ and
    $\text{mod}(y,100)\ne0$). Pure function. \\
  \code{mdays\_leap} & \code{(iyear, imon) $\to$ integer(ik4)} &
    Private. Returns the number of days in month \code{imon} of year
    \code{iyear} for the Gregorian calendar, accounting for leap February.
    Pure. \\
  \code{ndaypm} & \code{(y, m, cal) $\to$ integer(ik4)} &
    Public. Days in month \code{m} for any calendar: calls
    \code{mdays\_leap} for Gregorian, returns \code{mlen(m)} for
    \code{noleap}, 30 for \code{y360}, and 0 for an unknown calendar.
    Pure. \\
  \code{yeardays} & \code{(y, c) $\to$ integer(ik4)} &
    Public. Total days in year \code{y} under calendar \code{c}:
    365 or 366 (Gregorian, also the default), 365 (\code{noleap}),
    360 (\code{y360}). Pure. \\
\end{longtable}

\subsubsection{Internal Representation Conversion}

The internal representation and the human-readable \code{(year, month, day)}
decomposition are related by two complementary pairs of private routines:

\begin{longtable}{L{5.5cm} L{8.5cm}}
  \toprule
  \textbf{Routine} & \textbf{Description} \\
  \midrule
  \endhead
  \bottomrule
  \endfoot
  \code{date\_to\_days\_from\_reference} &
    Converts an \code{iadate} to the \code{days\_from\_reference} counter by
    iterating over whole years from \code{reference\_year} (1900) using
    \code{yeardays}, then adding the day-of-year computed by \code{idayofyear}.
    Pure. \\
  \code{days\_from\_reference\_to\_date} &
    Inverse: iterates forward or backward year-by-year until the remaining
    day count fits within one year, then calls
    \code{idayofyear\_to\_monthdate} to find month and day. \\
  \code{time\_to\_second\_of\_day} &
    Packs \code{(hour, minute, second)} into \code{second\_of\_day}:
    $\text{sod} = 3600h + 60m + s$. \\
  \code{second\_of\_day\_to\_time} &
    Inverse: recovers \code{(hour, minute, second)} from \code{second\_of\_day}
    by integer division and modulo. \\
  \code{date\_time\_to\_internal} &
    Calls both conversion routines to construct a full
    \code{rcm\_time\_and\_date} from an \code{(iadate, iatime)} pair. \\
  \code{internal\_to\_date\_time} &
    Calls both inverse routines to recover \code{(iadate, iatime)}
    from a full \code{rcm\_time\_and\_date}. \\
\end{longtable}

\subsubsection{Day-of-Year Utilities}

\begin{longtable}{L{4cm} L{9cm}}
  \toprule
  \textbf{Routine} & \textbf{Description} \\
  \midrule
  \endhead
  \bottomrule
  \endfoot
  \code{idayofyear(x)} &
    Returns the day-of-year (1-based) of an \code{iadate} \code{x},
    accumulating month lengths. Pure, calendar-aware. \\
  \code{dayofyear(x)} &
    Public wrapper: converts \code{rcm\_time\_and\_date} to \code{iadate}
    and calls \code{idayofyear}. \\
  \code{idayofyear\_to\_monthdate} &
    Inverse: given a zero-based day-of-year, iterates through months to
    find \code{(month, day)}. Calendar-aware, pure. \\
\end{longtable}

\subsubsection{Month Helper}

The private \code{add\_months(d, n)} adds $n$ months, positive or
negative, to the month and year of an \code{iadate}, with modulo
arithmetic on the zero-based month; the day is left to the caller.

%------------------------------------------------------------------------------
\subsection{Julian Day Computation}
\label{subsec:date_julian}

\begin{longtable}{L{5cm} L{9cm}}
  \toprule
  \textbf{Function} & \textbf{Description} \\
  \midrule
  \endhead
  \bottomrule
  \endfoot
  \code{julianday(iy, im, id) $\to$ real(rk8)} &
    Generic interface with the single procedure \code{ymd\_julianday}.
    Returns the Julian Day of the date (year, month, day) at 00:00~UTC in
    the proleptic Gregorian calendar, so the value always ends in $.5$.
    Pure. \\
  \code{idayofweek(x) $\to$ integer(ik4)} &
    Returns the day of week (1 = Sunday, \ldots, 7 = Saturday) from the
    Julian Day: $\text{dow} = \text{modulo}(JD + 1.5, 7) + 1$. The
    \code{modulo} intrinsic keeps the result in $[1,7]$ also for negative
    Julian Days. Requires the Gregorian calendar; aborts otherwise. \\
\end{longtable}

\code{ymd\_julianday} follows the model calendar: the Gregorian rules are
applied to every year, including those before 1582, and every
(year, month, day) triple maps to a Julian Day. The computation counts
days from 0000-03-01, so that the leap day is the last day of each
counting year:
\begin{enumerate}
  \item Shift the year to start in March: $y' = y - 1$ and $m' = m + 9$ for
        January and February, $y' = y$ and $m' = m - 3$ otherwise, so that
        $m' = 0$ is March and $m' = 11$ is February.
  \item Split $y'$ into 400-year eras, which all have exactly
        146\,097~days: $e = (y' - \text{modulo}(y', 400))/400$ and
        $y_e = y' - 400e \in [0, 399]$. Using \code{modulo} makes this valid
        also for negative years.
  \item Add the days of the whole eras, of the whole years of the era, and
        of the whole months of the year:
        \[
          JD = 1721119.5 + 146097\,e + 365\,y_e
               + \left\lfloor \frac{y_e}{4} \right\rfloor
               - \left\lfloor \frac{y_e}{100} \right\rfloor
               + \left\lfloor \frac{153\,m' + 2}{5} \right\rfloor
               + d - 1
        \]
        where 1721119.5 is the Julian Day of 0000-03-01 00:00~UTC and
        $\lfloor (153m'+2)/5 \rfloor$ is the number of days from 1 March
        to the first day of month $m'$.
\end{enumerate}
All the terms are integer, so the result is exact in \code{real(rk8)}. For
example, 2000-01-01 gives $JD = 2451544.5$ and \code{idayofweek} = 7
(Saturday).

%------------------------------------------------------------------------------
\subsection{Public Conversion and Formatting Functions}
\label{subsec:date_conversion}

\subsubsection{Date/Time to String}

\begin{longtable}{L{4cm} L{4cm} L{6cm}}
  \toprule
  \textbf{Function} & \textbf{Returns} & \textbf{Format} \\
  \midrule
  \endhead
  \bottomrule
  \endfoot
  \code{tochar(x)} &
    \code{character(len=32)} &
    \texttt{YYYY-MM-DD HH:MM:SS UTC}, preceded by a blank (or by
    \texttt{-} for negative years). \\
  \code{toiso8601(x)} &
    \code{character(len=21)} &
    \texttt{YYYY-MM-DDTHH:MM:SSZ} (ISO~8601 extended format), preceded by
    a blank or by \texttt{-} as above. \\
  \code{tochar10(x)} &
    \code{character(len=11)} &
    Zero-padded packed integer \texttt{YYYYMMDDHH}, left-adjusted. \\
\end{longtable}

\subsubsection{Date/Time to Numeric}

\begin{longtable}{L{4cm} L{4cm} L{6cm}}
  \toprule
  \textbf{Function} & \textbf{Returns} & \textbf{Description} \\
  \midrule
  \endhead
  \bottomrule
  \endfoot
  \code{toint10(x)} &
    \code{integer(ik8)} &
    Packed integer $\text{YYYY}\times10^6 + \text{MM}\times10^4
    + \text{DD}\times10^2 + \text{HH}$. \\
  \code{tohours(x)} &
    \code{real(rk8)} &
    Converts an \code{rcm\_time\_interval} to hours. Returns 0 and calls
    \code{die} for calendar-dependent units. \\
  \code{toseconds(x)} &
    \code{integer(ik8)} &
    Converts an \code{rcm\_time\_interval} to seconds. Same restriction.
    Used by the interval addition when the units differ. \\
  \code{yeardayfrac(x)} &
    \code{real(rkx)} &
    Day-of-year as a fractional real (including hour, minute, second
    contributions), 0-based. \\
  \code{yearpoint(x)} &
    \code{real(rk8)} &
    Calendar day used by the orbital routine \code{orb\_decl} (in
    \code{mod\_sun} and in the CLM~4.5 interface). For the Gregorian
    calendar it is the day of the year $n$ (1-based, the time of day is
    ignored) minus $(n + k)/1461$, where $k$ is 1095, 730, 365 or 0 when the
    year is a leap year or precedes one by one, two or three years: the
    correction spreads the leap day over the four-year cycle. For
    \code{noleap} and \code{y360} it is the day of the year. \\
\end{longtable}

\subsubsection{Component Extractors}

Six scalar functions each decode one field from an
\code{rcm\_time\_and\_date}:

\begin{longtable}{L{4.5cm} L{4cm} L{5.5cm}}
  \toprule
  \textbf{Function} & \textbf{Returns} & \textbf{Description} \\
  \midrule
  \endhead
  \bottomrule
  \endfoot
  \code{getyear(x)}   & \code{integer(ik4)} & Calendar year \\
  \code{getmonth(x)}  & \code{integer(ik4)} & Month (1--12) \\
  \code{getday(x)}    & \code{integer(ik4)} & Day of month \\
  \code{gethour(x)}   & \code{integer(ik4)} & Hour of day (0--23) \\
  \code{getminute(x)} & \code{integer(ik4)} & Minute of hour (0--59) \\
  \code{getsecond(x)} & \code{integer(ik4)} & Second of minute (0--59) \\
\end{longtable}

\subsubsection{Decomposition Subroutines (\texttt{split\_idate})}

The generic interface \code{split\_idate} is overloaded across five procedures
that unpack an \code{rcm\_time\_and\_date} (or an integer date) into its
components:

\begin{longtable}{L{6.5cm} L{8.5cm}}
  \toprule
  \textbf{Procedure} & \textbf{Outputs} \\
  \midrule
  \endhead
  \bottomrule
  \endfoot
  \code{split\_i10} &
    From \code{integer(ik4)}: year, month, day, hour. \\
  \code{split\_i10\_8} &
    From \code{integer(ik8)}: year, month, day, hour. \\
  \code{split\_rcm\_date} &
    From \code{rcm\_time\_and\_date}: year, month, day. \\
  \code{split\_rcm\_time\_and\_date} &
    From \code{rcm\_time\_and\_date}: year, month, day, hour. \\
  \code{split\_rcm\_time\_and\_date\_complete} &
    From \code{rcm\_time\_and\_date}: year, month, day, hour, minute, second. \\
\end{longtable}

%------------------------------------------------------------------------------
\subsection{Navigation and Comparison Functions}
\label{subsec:date_navigation}

\subsubsection{Month and Year Navigation}

\begin{longtable}{L{4.5cm} L{4cm} L{5.5cm}}
  \toprule
  \textbf{Function/Subroutine} & \textbf{Returns} & \textbf{Description} \\
  \midrule
  \endhead
  \bottomrule
  \endfoot
  \code{monfirst(x)}  & \code{rcm\_time\_and\_date} &
    First day of the month containing \code{x}, at midnight. \\
  \code{monlast(x)}   & \code{rcm\_time\_and\_date} &
    Last day of the month containing \code{x}, at midnight. \\
  \code{monmiddle(x)} & \code{rcm\_time\_and\_date} &
    Middle of the month: day $\lfloor N/2 \rfloor$ at hour 0 or 12
    depending on whether $N$ is odd. For \code{y360} always day 15. \\
  \code{nextmon(x)}   & \code{rcm\_time\_and\_date} &
    Same date one month later. \\
  \code{prevmon(x)}   & \code{rcm\_time\_and\_date} &
    Same date one month earlier. \\
  \code{yrfirst(x)}   & \code{rcm\_time\_and\_date} &
    January 1 of the year containing \code{x}, at midnight. \\
  \code{nextwk(x)}    & \code{rcm\_time\_and\_date} &
    Same time 7 days later. \\
  \code{prevwk(x)}    & \code{rcm\_time\_and\_date} &
    Same time 7 days earlier. \\
\end{longtable}

\subsubsection{Week Navigation (Gregorian Only)}

\begin{longtable}{L{4.5cm} L{4cm} L{5.5cm}}
  \toprule
  \textbf{Function} & \textbf{Returns} & \textbf{Description} \\
  \midrule
  \endhead
  \bottomrule
  \endfoot
  \code{ifdoweek(x)} & \code{rcm\_time\_and\_date} &
    First day (Sunday) of the week containing \code{x}, at midnight. \\
  \code{ildoweek(x)} & \code{rcm\_time\_and\_date} &
    Last day (Saturday) of the week containing \code{x}, at midnight. \\
  \code{iwkdiff(x,y)}& \code{integer(ik4)} &
    Number of whole weeks in \code{x - y} plus one, i.e.\ the 1-based index
    of the week of \code{x} counted from \code{y}. \\
  \code{lsameweek(x,y)} & \code{logical} &
    Meant to test whether \code{x} and \code{y} fall in the same
    Sunday-to-Saturday week. It returns true if \code{y} lies strictly
    between 00:00 of the Saturday before the week of \code{x} and 00:00 of
    the Saturday of that week. Not used in the current code. \\
\end{longtable}

\subsubsection{Logical Predicates}

\begin{longtable}{L{4.5cm} L{4cm} L{5.5cm}}
  \toprule
  \textbf{Function} & \textbf{Returns} & \textbf{Description} \\
  \midrule
  \endhead
  \bottomrule
  \endfoot
  \code{lmidnight(x)}    & \code{logical} &
    True if \code{second\_of\_day} is zero. \\
  \code{lfdomonth(x)}    & \code{logical} &
    True if \code{x} falls on day 1 of its month. \\
  \code{lfdoyear(x)}     & \code{logical} &
    True if \code{x} falls on January 1 of its year. \\
  \code{lfhomonth(x)}    & \code{logical} &
    True if \code{x} is in the first half of its month
    (i.e.\ \code{x < monmiddle(x)}). \\
  \code{lsamemonth(x,y)} & \code{logical} &
    True if \code{x} and \code{y} have the same month number, whatever
    their year (calendar check). \\
\end{longtable}

\subsubsection{Difference Functions}

\begin{longtable}{L{4.5cm} L{4cm} L{5.5cm}}
  \toprule
  \textbf{Function} & \textbf{Returns} & \textbf{Description} \\
  \midrule
  \endhead
  \bottomrule
  \endfoot
  \code{hourdiff(x,y)}  & \code{real(rk8)} &
    Signed difference in hours: $(x - y)$ computed exactly from the integer
    internal fields. \\
  \code{imondiff(x,y)}  & \code{integer(ik4)} &
    Difference in months, computed as
    $12\,|y_x - y_y| + (m_x - m_y)$: it is the number of months from
    \code{y} to \code{x} when \code{x} is not earlier than \code{y}. \\
  \code{iyeardiff(x,y)} & \code{integer(ik4)} &
    Absolute difference in years. \\
\end{longtable}

\subsubsection{Date and Time Predicates}

The generic interfaces \code{date\_is} and \code{time\_is} provide
tolerance-aware moment testing:

\begin{longtable}{L{5.5cm} L{8.5cm}}
  \toprule
  \textbf{Function} & \textbf{Description} \\
  \midrule
  \endhead
  \bottomrule
  \endfoot
  \code{date\_is(x, y, m, d)} &
    True if the date part of \code{x} equals year \code{y}, month \code{m},
    day \code{d}. \\
  \code{date\_is(x, m, d)} &
    True if the month and day of \code{x} equal \code{m} and \code{d}
    (year-independent). \\
  \code{time\_is(x, h, m, s [,delta])} &
    True if the time of day of \code{x} equals (h, m, s). With optional
    \code{delta} [s], uses a tolerance window
    $|t_x - t_\text{ref}| < \delta$ instead. \\
  \code{time\_is(x, s [,delta])} &
    As above but compares \code{second\_of\_day} directly to \code{s}. \\
\end{longtable}

%------------------------------------------------------------------------------
\subsection{CF-Conventions and External Time Interface}
\label{subsec:date_cf}

\subsubsection{\texttt{timeval2date} --- Numeric Time Value to Date}

The generic interface \code{timeval2date} converts a real-valued ``time since
reference'' (as stored in CF-Conventions NetCDF files) to an
\code{rcm\_time\_and\_date}. It is overloaded for single (\code{rk4}) and
double (\code{rk8}) precision inputs.

\begin{lstlisting}
dd = timeval2date(xval, cunit, ccal)
\end{lstlisting}

\begin{longtable}{L{2.5cm} L{3cm} L{8.5cm}}
  \toprule
  \textbf{Argument} & \textbf{Type} & \textbf{Description} \\
  \midrule
  \endhead
  \bottomrule
  \endfoot
  \code{xval}  & \code{real(rk4/rk8)} &
    Numeric time value (e.g.\ \texttt{hours since 1979-01-01} $\to$
    \code{xval} in hours). May be negative. \\
  \code{cunit} & \code{character(*)} &
    CF time-unit string, e.g.\ \texttt{hours since 1979-01-01 00:00:00}. \\
  \code{ccal}  & \code{character(*)} &
    CF calendar name (\texttt{noleap}, \texttt{360\_day}, or
    \texttt{gregorian}/\texttt{standard}). \\
\end{longtable}

The routine caches \code{cunit}, \code{ccal} and the reference date in
\code{save} variables (separately for the two precisions), so that the
string parse (via \code{parse\_timestring}) is only performed on the first
call or when the unit or calendar string changes. The value is truncated
towards zero to a whole number of units, and the fraction is discarded
except for days: a value of 15.5 days is interpreted as 15 days plus
12 hours, consistent with common NetCDF practice (the hours are truncated as
well). The calendar strings recognised are listed in
Section~\ref{subsec:date_constants}.

\subsubsection{\texttt{timeval2ym} --- Numeric Time Value to Year/Month}

\begin{lstlisting}
call timeval2ym(xval, cunit, year, month)
\end{lstlisting}

Lighter-weight variant that parses only the year and month from a
\texttt{months since \ldots} unit string. Uses the same caching strategy.

\subsubsection{\texttt{parse\_timestring} --- CF Time-Unit String Parser}

Private function that parses a CF-Conventions time-unit string of the form
\texttt{UNIT since YYYY-MM-DD HH:MM:SS}. The unit is recognised by the first
word (\texttt{seconds}, \texttt{minutes}, \texttt{hours}, \texttt{days},
\texttt{months} or \texttt{years}), and the date is read from the fixed
position after \texttt{UNIT since }. Months, days, hours and minutes may
have one or two digits; the date and the time may be separated by a blank or
by \texttt{T}. The string may stop after any field: the missing fields keep
their defaults (1900-01-01 00:00:00). Returns:
\begin{itemize}
  \item The unit selector (\code{aunit}): one of \code{usec}--\code{uyrs},
        or $-1$ for an unrecognised unit.
  \item The reference date components: year, month, day, hour, minute, second.
  \item An integer result code: 0 on success, $-1$ if a date field cannot be
        read. An unrecognised unit is reported on \code{stderr} but returns
        0, with \code{aunit} $= -1$.
\end{itemize}

%------------------------------------------------------------------------------
\subsection{Model Interface Routines}
\label{subsec:date_model}

\subsubsection{Reference and Current Time}

\begin{longtable}{L{5.5cm} L{8.5cm}}
  \toprule
  \textbf{Routine} & \textbf{Description} \\
  \midrule
  \endhead
  \bottomrule
  \endfoot
  \code{ref\_date(iy, im, id, isec)} &
    Returns the fixed CORDEX reference date: 1 January 1950, 00:00:00.
    Hard-coded; the year 1950 is the WCRP CORDEX standard epoch. \\
  \code{curr\_date(x, iy, im, id, isec [, offset])} &
    Decomposes an \code{rcm\_time\_and\_date} into (year, month, day,
    second-of-day). The optional \code{offset} argument (in seconds) shifts
    the timestamp before decomposition, used to advance to the next time
    step for output labelling. \\
  \code{curr\_time(x, iday, isec)} &
    Returns days and seconds elapsed since the CORDEX reference date
    (1950-01-01). Uses a \code{save}d copy of the reference date to avoid
    recomputing it on every call; reinitialises if the calendar changes. \\
\end{longtable}

\subsubsection{Calendar String/Integer Conversion}

\begin{longtable}{L{5cm} L{9cm}}
  \toprule
  \textbf{Function} & \textbf{Description} \\
  \midrule
  \endhead
  \bottomrule
  \endfoot
  \code{calendar\_str(x)} &
    Returns the CF-Conventions calendar string
    (\texttt{proleptic\_gregorian}, \texttt{noleap}, or \texttt{360\_day})
    for the calendar stored in \code{x}. \\
  \code{calendar\_int(calstr)} &
    Pure function: maps a CF calendar string to the corresponding integer
    constant (\code{gregorian}, \code{noleap}, or \code{y360}).
    Unknown strings silently default to \code{gregorian}. \\
\end{longtable}

\subsubsection{Calendar Setting}

The generic interface \code{setcal} changes the calendar of an existing
\code{rcm\_time\_and\_date}, recomputing \code{days\_from\_reference}
for the new calendar:

\begin{longtable}{L{5cm} L{9cm}}
  \toprule
  \textbf{Procedure} & \textbf{Description} \\
  \midrule
  \endhead
  \bottomrule
  \endfoot
  \code{set\_calint(x, c)} &
    Sets the calendar to the integer constant \code{c} (an unknown value
    selects \code{gregorian}). If the calendar changes, decomposes the
    current date and recomposes the same (year, month, day) under the new
    system, so the day counter changes. Pure. \\
  \code{set\_calstring(x, c)} &
    Maps the string \code{c} to a calendar (see the table in
    Section~\ref{subsec:date_constants}) and delegates to
    \code{set\_calint}. \\
  \code{set\_caltype(x, y)} &
    Copies the calendar from \code{y} into \code{x}, calling
    \code{set\_calint} only if they differ. \\
  \code{set\_timeunit(x, u)} &
    Sets the unit of an \code{rcm\_time\_interval}; validates against the
    seven legal unit selectors and defaults to \code{uhrs} with a warning
    for unknown values. \\
\end{longtable}

\subsubsection{Diagnostic Utilities}

\begin{longtable}{L{5.5cm} L{8.5cm}}
  \toprule
  \textbf{Routine} & \textbf{Description} \\
  \midrule
  \endhead
  \bottomrule
  \endfoot
  \code{print\_rcm\_time\_and\_date(x)} &
    Writes the formatted timestamp to \code{stdout} via \code{tochar}. \\
  \code{print\_rcm\_time\_interval(x)} &
    Writes the interval magnitude and unit name to \code{stdout}. \\
  \code{season\_index(x)} &
    Returns the meteorological season index
    (\code{isea\_djf}, \ldots, \code{isea\_son}). \\
  \code{date\_in\_scenario(x, cmip\_cycle [, l2006])} &
    Returns \code{.true.} if \code{x} falls in the CMIP scenario period:
    after 2000 for CMIP3, after 2005 (or December 2005, controlled by
    \code{l2006}) for CMIP5. \\
\end{longtable}

%------------------------------------------------------------------------------
\subsection{Design Notes}
\label{subsec:date_design}

\textbf{Integer-only internal state.} All time quantities stored in
\code{rcm\_time\_and\_date} are exact integers. This means that date
comparisons (\code{==}, \code{<}, etc.) are exact and never suffer from
floating-point epsilon issues, which would be a serious problem in
output-trigger logic run at every model time step.

\textbf{Calendar enforcement.} Every multi-date operation (\code{-},
comparison operators, \code{lsamemonth}, etc.) calls the private
\code{check\_cal} guard, which terminates the run if the two operands use
different calendars. This prevents silent incorrect results from mixing,
say, a Gregorian simulation date with a no-leap boundary date.

\textbf{Month arithmetic and day-of-month preservation.} Adding or subtracting
months is inherently ambiguous (what is 31 January + 1 month?). The module
resolves this by computing a fractional position $\alpha$ within the source
month and mapping it proportionally onto the target month, so that the
``middle of January'' becomes the ``middle of February'' rather than day 28
or day 1. The new day is $\text{nint}(N_\text{target}\,\alpha) + 1$ with
$\alpha = (d - 1)/N_\text{source}$; the source comments call it ``a trick'',
whose purpose is to make day 14 the middle of February in the Gregorian
calendar. Only the addition of months uses it (see the table of the
\code{+} and \code{-} operators).

\textbf{Corrections.} The following were corrected in October 2026, and
are checked by \code{Testing/share/test\_date.F90}: the subtraction of
seconds added them, and that of minutes and hours gave a negative second
of the day when the result was on a previous day from a time other than
midnight; the subtraction of months could give a day beyond the end of
the month (31 March $-$ 1 month was 3 March); the addition of a negative
number of months gave month 0 instead of going back to December. No code
outside the module subtracted seconds, minutes or hours, and
\code{prevmon} is applied to dates up to the middle of the month, so the
model results do not change. The unused recursive day helpers
(\code{add\_days\_leap}, \code{sub\_days\_leap}, \code{add\_days\_noleap},
\code{sub\_days\_noleap}, the last two inconsistent with the first two) and
\code{adjustpm}, \code{adjustmp} were removed.

\textbf{CF-Conventions compatibility.} The \code{timeval2date},
\code{timeval2ym}, and \code{parse\_timestring} routines ensure that any
time axis read from a CF-compliant NetCDF file can be decoded directly into
the internal representation, regardless of the reference date, unit, or
calendar stored in the file's \texttt{time:units} and
\texttt{time:calendar} attributes.

\textbf{CORDEX time output.} The \code{curr\_time} subroutine expresses the
model time as (days, seconds) from 1950-01-01 00:00:00, which is the standard
CORDEX time reference used in all output NetCDF files. The reference date is
initialised once using a \code{save}d variable to avoid repeated construction.

%==============================================================================
\section{Global Parameters and Namelist Interface --- \texttt{mod\_dynparam.F90}}
\label{sec:dynparam}
%==============================================================================

\subsection{Overview}

\code{mod\_dynparam} is the central configuration registry of RegCM. It
declares --- as public module-level variables --- every parameter that
controls a simulation: grid dimensions, map projection, physics switches,
dynamical core selection, I/O frequencies, boundary condition settings,
ensemble perturbation fractions, and CMIP/PMIP dataset identifiers. It also
provides the single entry point for reading the model namelist file
(\code{initparam}) and exposes helper subroutines for nested-domain and
global-window configuration.

Because all parameters are plain module variables (not derived-type members),
any module that \code{use}s \code{mod\_dynparam} gains read/write access to
the entire global configuration without argument passing. This is the
traditional RegCM design pattern: a single \code{initparam} call early in
the main program populates the module, and all subsequent modules read from
it directly.

\textbf{Key responsibilities:}
\begin{itemize}
  \item Hard-coded array-size ceilings (\code{maxvert}, \code{maxtracer})
        that are used as compile-time bounds in CUDA/OpenACC kernel launches.
  \item Grid-dimension variables (\code{iy}, \code{jx}, \code{kz}) and the
        full set of derived convenience dimensions (\code{iym1}, \code{kzp1},
        etc.).
  \item Staggered-grid index ranges for cross and dot points, MPI
        subdomain decomposition bounds, ghost-cell stencil extents, and
        sub-grid index ranges.
  \item Calendar, date, and simulation-period state (\code{globidate1},
        \code{globidate2}, \code{dayspy}, \code{dpd}, etc.).
  \item NetCDF I/O mode selection at compile time via preprocessor flags.
  \item Random number generator seeding (CPU and optional GPU/CUDA paths).
\end{itemize}

\subsection{Dependencies}

\begin{lstlisting}
mod_dynparam
    |-- mod_stdio       (stderr, stdout)
    |-- mod_intkinds    (ik4, ik8)
    |-- mod_realkinds   (rkx, rk8)
    |-- mod_constants   (erkm, degrad, twopi, d_two, d_r1000, ...)
    |-- mod_date        (rcm_time_and_date, i8wcal, gregorian, noleap, y360)
    |-- netcdf          [conditional: absent when PNETCDF is defined]
    |-- mpi_f08         (mpi_comm type only)
\end{lstlisting}

The \code{netcdf} module is omitted when the build uses PnetCDF
(\code{PNETCDF} macro), in which case NetCDF constants such as
\code{nf90\_clobber} are provided through the PnetCDF interface instead.
Only the \code{mpi\_comm} derived type is imported from \code{mpi\_f08};
no MPI communication calls are made inside this module.

\subsection{Compile-Time Preprocessor Flags}

\begin{longtable}{L{5cm} L{9cm}}
  \toprule
  \textbf{Flag} & \textbf{Effect} \\
  \midrule
  \endhead
  \bottomrule
  \endfoot
  \code{PNETCDF} &
    Excludes the \code{use netcdf} statement; \code{iomode} is set to 0
    (PnetCDF sets its own mode). \\
  \code{NETCDF4\_HDF5} &
    Sets \code{iomode = ior(nf90\_clobber, nf90\_netcdf4)} for HDF5-backed
    NetCDF4 output. \\
  \code{NETCDF4\_HDF5 + NETCDF4\_COMPRESS} &
    Additionally reads the \code{ncfilters} namelist to configure HDF5
    compression filter parameters. \\
  \code{NETCDF\_CDF5} &
    Sets \code{iomode = ior(nf90\_clobber, nf90\_cdf5)} for 64-bit data
    CDF-5 format. \\
  \code{CLM45} &
    Enables the CLM 4.5 land-surface scheme; adds the \code{scalez},
    \code{num\_soil\_layers} (10 instead of 3), and
    \code{clm\_regcm} namelist group, and the four CLM feature flags. \\
  \code{OPENACC} &
    Exposes \code{init\_gpu\_random\_number\_generator}, which seeds the
    cuRAND XORWOW generator used by OpenACC device kernels. \\
  \code{NAGFOR} &
    Triggers a call to \code{init\_realkinds} at the start of
    \code{initparam} (see \code{mod\_realkinds} documentation). \\
  \code{DEBUG} &
    Enables \code{checkalloc} checks in \code{mod\_message}
    (used transitively). \\
\end{longtable}

%------------------------------------------------------------------------------
\subsection{Hard-Coded Compile-Time Ceilings}
\label{subsec:dynparam_ceilings}

\begin{longtable}{L{3.5cm} L{2.5cm} L{8cm}}
  \toprule
  \textbf{Name} & \textbf{Value} & \textbf{Description} \\
  \midrule
  \endhead
  \bottomrule
  \endfoot
  \code{maxvert}   & 256 &
    Maximum number of vertical levels. Fixed to permit static allocation in
    CUDA/OpenACC kernels, where dynamic allocation is expensive. Comment
    notes that ECMWF uses 137 levels and MIROC-ESM uses 80 levels; 256
    provides ample headroom. Must be $\geq$ \code{kz}; \code{initparam}
    enforces this at runtime. \\
  \code{maxtracer} & 64 &
    Maximum number of chemical tracer species. The MINE chemistry
    configuration used approximately 56 tracers as of September 2026. \\
\end{longtable}

%------------------------------------------------------------------------------
\subsection{Grid and Projection Variables}
\label{subsec:dynparam_grid}

All variables in this section are module-level \code{public} scalars; they
are set by reading the \code{dimparam} and \code{geoparam} namelists inside
\code{initparam}.

\subsubsection{Primary Grid Dimensions}

\begin{longtable}{L{3cm} L{3cm} L{8cm}}
  \toprule
  \textbf{Name} & \textbf{Type} & \textbf{Description} \\
  \midrule
  \endhead
  \bottomrule
  \endfoot
  \code{iy}  & \code{integer(ik4)} &
    Number of grid points in the Y (latitude) direction, including boundary
    zones. \\
  \code{jx}  & \code{integer(ik4)} &
    Number of grid points in the X (longitude) direction. \\
  \code{kz}  & \code{integer(ik4)} &
    Number of vertical levels. Must satisfy \code{kz} $\leq$ \code{maxvert}. \\
  \code{nsg} & \code{integer(ik4)} &
    Sub-grid refinement factor for CLM or statistical downscaling.
    Default: 1 (no sub-gridding). \\
  \code{ds}  & \code{real(rkx)} &
    Horizontal grid spacing in km. For ROTLLR projection a negative value
    (degrees) is accepted and converted automatically. \\
  \code{dsmax}, \code{dsmin} & \code{real(rkx)} &
    Maximum and minimum normalised sigma level spacing. Defaults
    0.05 and 0.01 respectively; used to bound the sigma-level
    distribution. \\
  \code{ptop} & \code{real(rkx)} &
    Pressure at the model top in cbar (centibars). Set to 0 for the
    non-hydrostatic Moloch core (\code{idynamic}~=~3). \\
\end{longtable}

\subsubsection{Derived Convenience Dimensions}

Computed from \code{iy}, \code{jx}, \code{kz}, and \code{nsg} at the end
of \code{initparam}. Named to avoid repeated off-by-one arithmetic
throughout the model.

\begin{longtable}{L{3.5cm} L{8.5cm}}
  \toprule
  \textbf{Name(s)} & \textbf{Description} \\
  \midrule
  \endhead
  \bottomrule
  \endfoot
  \code{iym1}, \code{iym2}, \code{iym3} &
    $\text{iy}-1$, $\text{iy}-2$, $\text{iy}-3$. \\
  \code{jxm1}, \code{jxm2}, \code{jxm3} &
    $\text{jx}-1$, $\text{jx}-2$, $\text{jx}-3$. \\
  \code{kzm1}, \code{kzm2} &
    $\text{kz}-1$, $\text{kz}-2$ (below top layer indices). \\
  \code{kzp1}, \code{kzp2}, \code{kzp3}, \code{kzp4} &
    $\text{kz}+1$ through $\text{kz}+4$ (interface and ghost layer indices). \\
  \code{iysg}, \code{jxsg} &
    Sub-grid dimensions: $\text{iy}\times\text{nsg}$, $\text{jx}\times\text{nsg}$. \\
  \code{iym1sg}\ldots\code{iym3sg}, \code{jxm1sg}\ldots\code{jxm3sg} &
    Sub-grid analogues of the $-1$, $-2$, $-3$ offsets. \\
  \code{nnsg} &
    $\text{nsg}^2$; number of sub-grid cells per coarse cell. \\
\end{longtable}

\subsubsection{Map Projection Parameters}

Read from the \code{geoparam} namelist.

\begin{longtable}{L{3cm} L{3cm} L{8cm}}
  \toprule
  \textbf{Name} & \textbf{Type} & \textbf{Description} \\
  \midrule
  \endhead
  \bottomrule
  \endfoot
  \code{iproj} & \code{character(6)} &
    Projection identifier. Accepted values: \texttt{LAMCON} (Lambert
    Conformal), \texttt{POLSTR} (Polar Stereographic), \texttt{NORMER}
    (Normal Mercator), \texttt{ROTMER} (Rotated Mercator),
    \texttt{ROTLLR} (Rotated Lat-Lon, degree spacing). \\
  \code{clat}, \code{clon} & \code{real(rkx)} &
    Central latitude and longitude of the projection [degrees].
    North/East positive. \\
  \code{plat}, \code{plon} & \code{real(rkx)} &
    Pole latitude/longitude for rotated Mercator; equals \code{clat}/\code{clon}
    for other projections. \\
  \code{cntri}, \code{cntrj} & \code{real(rkx)} &
    Grid index (I, J) of the projection centre. Defaults to
    $\text{iy}/2$, $\text{jx}/2$ if not specified (value $< 0$). \\
  \code{truelatl}, \code{truelath} & \code{real(rkx)} &
    Lower and upper true latitudes for Lambert Conformal. If both are zero
    on entry, \code{initparam} warns and sets both to \code{clat}. \\
  \code{xcone} & \code{real(rkx)} &
    Lambert/Polar cone factor (computed elsewhere from the true latitudes). \\
  \code{i\_band} & \code{integer(ik4)} &
    Tropical band mode flag (1 = active). In band mode \code{jcross2 = jx}
    (periodic zonal boundary). Automatically set to 1 when \code{i\_crm = 1}. \\
  \code{i\_crm} & \code{integer(ik4)} &
    Cloud-Resolving Model mode flag. Activates periodic boundaries in both
    directions. \\
\end{longtable}

%------------------------------------------------------------------------------
\subsection{Staggered-Grid Index Ranges}
\label{subsec:dynparam_stagger}

RegCM uses an Arakawa B-grid with \emph{cross} (scalar/thermodynamic) and
\emph{dot} (vorticity/wind) points. The following index-range variables
define the active portions of each grid type in the global (non-decomposed)
domain.

\subsubsection{Cross-Point Ranges}

\begin{longtable}{L{4cm} L{8cm}}
  \toprule
  \textbf{Name(s)} & \textbf{Description} \\
  \midrule
  \endhead
  \bottomrule
  \endfoot
  \code{ice1}, \code{ice2} & External (boundary-inclusive) I range: 1 to $\text{iym1}$
    (or \code{iy} in CRM mode). \\
  \code{jce1}, \code{jce2} & External J range: 1 to $\text{jxm1}$
    (or \code{jx} in band mode). \\
  \code{ici1}, \code{ici2} & Internal I range (first and last boundary row excluded). \\
  \code{jci1}, \code{jci2} & Internal J range. \\
  \code{icii1}, \code{icii2} & Doubly-internal I range (two boundary rows excluded). \\
  \code{jcii1}, \code{jcii2} & Doubly-internal J range. \\
  \code{nicross}, \code{njcross} & Number of cross points in I and J. \\
\end{longtable}

\subsubsection{Dot-Point Ranges}

\begin{longtable}{L{4cm} L{8cm}}
  \toprule
  \textbf{Name(s)} & \textbf{Description} \\
  \midrule
  \endhead
  \bottomrule
  \endfoot
  \code{ide1}, \code{ide2} & External I range for dot points: 1 to \code{iy}. \\
  \code{jde1}, \code{jde2} & External J range: 1 to \code{jx}. \\
  \code{idi1}, \code{idi2} & Internal I range. \\
  \code{jdi1}, \code{jdi2} & Internal J range. \\
  \code{idii1}, \code{idii2} & Doubly-internal I. \\
  \code{jdii1}, \code{jdii2} & Doubly-internal J. \\
  \code{nidot}, \code{njdot} & Number of dot points in I and J. \\
\end{longtable}

\subsubsection{Output, Sub-Grid, and Ghost-Cell Ranges}

\begin{longtable}{L{4cm} L{8cm}}
  \toprule
  \textbf{Name(s)} & \textbf{Description} \\
  \midrule
  \endhead
  \bottomrule
  \endfoot
  \code{iout1/2}, \code{jout1/2} &
    Output domain I/J range (interior only, excluding boundary sponge rows). \\
  \code{ioutsg1/2}, \code{joutsg1/2} &
    Sub-grid output domain range. \\
  \code{ide1sg/2sg}, \code{jde1sg/2sg} &
    Sub-grid dot-point external range. \\
  \code{ice1ga/2ga}\ldots & Ghost-layer A extents for cross and dot points
    (used in halo exchange). Analogous \code{gb} and \code{gc} sets for
    deeper halo depths. \\
  \code{ice1sl/2sl}\ldots & Semi-Lagrangian stencil extents. \\
\end{longtable}

%------------------------------------------------------------------------------
\subsection{MPI Decomposition Variables}
\label{subsec:dynparam_mpi}

\begin{longtable}{L{4cm} L{3cm} L{7cm}}
  \toprule
  \textbf{Name} & \textbf{Type} & \textbf{Description} \\
  \midrule
  \endhead
  \bottomrule
  \endfoot
  \code{mycomm}   & \code{type(mpi\_comm)} &
    MPI communicator for the model domain. \\
  \code{nproc}    & \code{integer(ik4)} &
    Total number of MPI processes. \\
  \code{nprocshm} & \code{integer(ik4)} &
    Number of processes sharing memory (shared-memory window). \\
  \code{myid}     & \code{integer(ik4)} &
    Rank of the current process in \code{mycomm}. \\
  \code{myidshm}  & \code{integer(ik4)} &
    Rank within the shared-memory group. \\
  \code{njxcpus}, \code{niycpus} & \code{integer(ik4)} &
    Requested decomposition: number of subdomains in J and I. $-1$ means
    auto-select (default). \\
  \code{iyp}, \code{jxp} & \code{integer(ik4)} &
    Local subdomain size in I and J for the current rank. \\
  \code{iypsg}, \code{jxpsg} & \code{integer(ik4)} &
    Local sub-grid subdomain size. \\
\end{longtable}

%------------------------------------------------------------------------------
\subsection{Dynamical Core and Physics Configuration}
\label{subsec:dynparam_core}

\subsubsection{Dynamical Core}

\begin{longtable}{L{3.5cm} L{3cm} L{7.5cm}}
  \toprule
  \textbf{Name} & \textbf{Value(s)} & \textbf{Description} \\
  \midrule
  \endhead
  \bottomrule
  \endfoot
  \code{idynamic} & 1, 2, or 3 &
    Dynamical core selector. 1 = hydrostatic RegCM (default), 2 = non-hydrostatic
    MM5-style, 3 = non-hydrostatic Moloch. Validated in \code{initparam};
    values outside [1,3] are fatal. \\
  \code{base\_state\_pressure} & \code{real(rkx)} &
    Base-state reference pressure [Pa] for \code{idynamic}~=~2.
    Default: 101\,325~Pa. \\
  \code{logp\_lrate} & \code{real(rkx)} &
    Log-pressure lapse rate $\mathrm{d}T/\mathrm{d}(\ln p)$ [K\,ln(Pa)$^{-1}$]
    for \code{idynamic}~=~2. Default: 47.70. \\
\end{longtable}

\subsubsection{Moloch Non-Hydrostatic Parameters}

Read from the \code{molochparam} namelist only when \code{idynamic}~=~3.

\begin{longtable}{L{4cm} L{3cm} L{7cm}}
  \toprule
  \textbf{Name} & \textbf{Default} & \textbf{Description} \\
  \midrule
  \endhead
  \bottomrule
  \endfoot
  \code{mo\_ztop}           & 30\,000~m & Model top height for Moloch. \\
  \code{mo\_h}              & 8\,000~m  & Reference scale height. \\
  \code{mo\_a0}             & 0.0       & Orography perturbation amplitude. \\
  \code{mo\_divdamp}        & \code{.true.} & Divergence damping switch. \\
  \code{mo\_spectral\_nudge}& \code{.true.} & Spectral nudging towards driving fields. \\
  \code{mo\_divfilter}      & \code{.true.} & Divergence filter switch. \\
  \code{mo\_top\_nudge}     & \code{.true.} & Upper-boundary nudging. \\
  \code{mo\_nadv}           & 2  & Number of advection sub-steps. \\
  \code{mo\_nsound}         & 5  & Number of acoustic sub-steps per advection step. \\
\end{longtable}

%------------------------------------------------------------------------------
\subsection{Boundary Conditions and Sponge Zone}
\label{subsec:dynparam_bdy}

\begin{longtable}{L{4.5cm} L{3cm} L{6.5cm}}
  \toprule
  \textbf{Name} & \textbf{Default} & \textbf{Description} \\
  \midrule
  \endhead
  \bottomrule
  \endfoot
  \code{nspgx}            & 6 (scaled) &
    Sponge zone depth in J (cross points). Scaled automatically to the grid
    spacing: $\max(\lfloor 6 \times 50\,\text{km}/ds \rfloor,\,6)$, then
    overridden by the \code{boundaryparam} namelist. For Moloch with Lehmann
    BCs, rounded to the next power-of-2 plus 2. \\
  \code{nspgd}            & 6 (scaled) &
    Sponge zone depth in I (dot points). Same scaling as \code{nspgx}. \\
  \code{bdy\_use\_lehmann}& \code{.false.} &
    Use the Lehmann spectral lateral boundary condition instead of the
    standard relaxation scheme. \\
  \code{high\_nudge}      & 3.0  & Nudging coefficient for the outermost sponge ring. \\
  \code{medium\_nudge}    & 2.0  & Nudging coefficient for intermediate ring. \\
  \code{low\_nudge}       & 1.0  & Nudging coefficient for innermost ring. \\
  \code{bdy\_nm}          & $-1$ & Normal-mode boundary coefficient ($-1$ = auto). \\
  \code{bdy\_dm}          & $-1$ & Diffusive boundary coefficient ($-1$ = auto). \\
\end{longtable}

%------------------------------------------------------------------------------
\subsection{Calendar and Simulation Period}
\label{subsec:dynparam_calendar}

\begin{longtable}{L{4.5cm} L{3cm} L{6.5cm}}
  \toprule
  \textbf{Name} & \textbf{Type} & \textbf{Description} \\
  \midrule
  \endhead
  \bottomrule
  \endfoot
  \code{calendar}        & \code{character(24)} &
    CF-Conventions calendar name from the namelist. \\
  \code{ical}            & \code{integer(ik4)} &
    Integer calendar selector (\code{gregorian}, \code{noleap}, or
    \code{y360} from \code{mod\_date}). \\
  \code{globidate1}      & \code{rcm\_time\_and\_date} &
    Global simulation start date (from \code{gdate1} in namelist). \\
  \code{globidate2}      & \code{rcm\_time\_and\_date} &
    Global simulation end date (from \code{gdate2} in namelist). \\
  \code{dayspy}          & \code{real(rkx)} &
    Days per year: 365.2422 (Gregorian), 365.0 (no-leap), 360.0 (360-day). \\
  \code{dpd}             & \code{real(rkx)} &
    Degrees of longitude per day: $360/\text{dayspy}$. \\
  \code{vernal\_equinox} & \code{real(rkx)} &
    Julian day of the vernal equinox, computed as
    $(80.50/\text{dayspy})\times365$ (fractional day-of-year). \\
  \code{half\_dayspy}    & \code{real(rkx)} & $\text{dayspy}/2$. \\
  \code{sixteenth\_dayspy} & \code{real(rkx)} & $\text{dayspy}/16$. \\
\end{longtable}

%------------------------------------------------------------------------------
\subsection{Input Data Source Variables}
\label{subsec:dynparam_data}

\subsubsection{Standard Analysis and Surface Datasets}

\begin{longtable}{L{4cm} L{3cm} L{7cm}}
  \toprule
  \textbf{Name} & \textbf{Default} & \textbf{Description} \\
  \midrule
  \endhead
  \bottomrule
  \endfoot
  \code{dattyp}  & \texttt{UNKNW} &
    5-character code for the global analysis dataset used in pre-processing.
    Examples: \texttt{ERA5}, \texttt{CMIP6}, \texttt{CMIP7}, \texttt{PMIP4}. \\
  \code{ssttyp}  & ---           &
    5-character code for the sea-surface temperature dataset.
    Examples: \texttt{OI2ST}, \texttt{CMIP6}, \texttt{ERSST}. \\
  \code{chemtyp} & \texttt{MZCLM}&
    5-character code for the chemistry boundary conditions.
    \texttt{MZ6HR} (MOZART 6-hourly) or \texttt{MZCLM} (MOZART climatology). \\
  \code{tersrc}  & \texttt{GMTED}& Terrain data source. \\
  \code{texsrc}  & \texttt{GLZB} & Soil texture data source. \\
  \code{smsrc}   & \texttt{ESACCI}& Soil moisture data source. \\
\end{longtable}

\subsubsection{CMIP6, CMIP7, and PMIP4 Identifiers}

Read from the \code{cmip6param}, \code{cmip7param}, or \code{pmip4param}
namelists when \code{dattyp} or \code{ssttyp} matches the corresponding
protocol identifier.

\begin{longtable}{L{4.5cm} L{6.5cm} L{3cm}}
  \toprule
  \textbf{Name} & \textbf{Description} & \textbf{Default} \\
  \midrule
  \endhead
  \bottomrule
  \endfoot
  \code{cmip6\_inp}        & OPeNDAP or local root URL/path for CMIP6 data &
    \texttt{https://esgf3.dkrz.de/\ldots} \\
  \code{cmip6\_model}      & Model identifier & \texttt{MPI-ESM1-2-HR} \\
  \code{cmip6\_variant}    & Variant label & \texttt{r1i1p1f1} \\
  \code{cmip6\_experiment} & Experiment identifier & \texttt{ssp585} \\
  \code{cmip6\_grid}       & Grid label & \texttt{gn} \\
  \code{cmip7\_inp}        & Root path for CMIP7 data & \texttt{cmip7} \\
  \code{cmip7\_model}      & Model identifier & \texttt{EC-Earth3-ESM-1-1} \\
  \code{cmip7\_experiment} & Experiment & \texttt{HL-ext} \\
  \code{cmip7\_region}     & Region tag & \texttt{glb} \\
  \code{cmip7\_atmo\_grid\_label} & Atmosphere grid label & \texttt{g114} \\
  \code{cmip7\_ocn\_grid\_label}  & Ocean grid label & \texttt{g102} \\
  \code{pmip4\_inp}        & Root URL for PMIP4 data & \texttt{https://esgf3.dkrz.de/\ldots} \\
  \code{pmip4\_model}      & Model identifier & \texttt{MPI-ESM1-2-LR} \\
  \code{pmip4\_experiment} & Experiment & \texttt{lgm} \\
  \code{pmip4\_variant}    & Variant label & \texttt{r1i1p1f2} \\
  \code{pmip4\_grid}       & Grid label & \texttt{gn} \\
\end{longtable}

%------------------------------------------------------------------------------
\subsection{I/O Control and Output Frequencies}
\label{subsec:dynparam_io}

\subsubsection{Output Switches}

\begin{longtable}{L{3cm} L{9cm}}
  \toprule
  \textbf{Name} & \textbf{Description} \\
  \midrule
  \endhead
  \bottomrule
  \endfoot
  \code{ifsave}   & Write restart files. \\
  \code{ifatm}    & Write atmospheric output. \\
  \code{ifshf}    & Write surface heat flux output. \\
  \code{ifrad}    & Write radiation output. \\
  \code{ifsrf}    & Write surface variable output. \\
  \code{ifsub}    & Write sub-grid output. \\
  \code{ifsts}    & Write surface temperature/soil output. \\
  \code{iflak}    & Write lake model output. \\
  \code{ifopt}    & Write optical depth / aerosol output. \\
  \code{ifchem}   & Write chemistry output. \\
  \code{ifcordex} & Write CORDEX-format output (additional variable set). \\
\end{longtable}

\subsubsection{Output Frequencies}

All frequencies are in hours unless the name implies a different unit.

\begin{longtable}{L{3.5cm} L{9cm}}
  \toprule
  \textbf{Name} & \textbf{Description} \\
  \midrule
  \endhead
  \bottomrule
  \endfoot
  \code{outnwf}  & Nominal output frequency [h]. \\
  \code{savfrq}  & Restart save frequency [h]. \\
  \code{atmfrq}  & Atmospheric output frequency [h]. \\
  \code{radfrq}  & Radiation output frequency [h]. \\
  \code{lakfrq}  & Lake model output frequency [h]. \\
  \code{subfrq}  & Sub-grid output frequency [h]. \\
  \code{srffrq}  & Surface output frequency [h]. \\
  \code{chemfrq} & Chemistry output frequency [h]. \\
  \code{optfrq}  & Optical/aerosol output frequency [h]. \\
  \code{ibdyfrq} & Boundary condition update frequency [h]. Default: 6. \\
\end{longtable}

\subsubsection{NetCDF I/O Mode}

\code{iomode} is an \code{integer(ik4)} whose value is set at compile time:

\begin{longtable}{L{5cm} L{9cm}}
  \toprule
  \textbf{Condition} & \textbf{\texttt{iomode} value} \\
  \midrule
  \endhead
  \bottomrule
  \endfoot
  \code{PNETCDF}              & 0 (PnetCDF manages mode internally) \\
  \code{NETCDF4\_HDF5}        & \code{ior(nf90\_clobber, nf90\_netcdf4)} \\
  \code{NETCDF\_CDF5}         & \code{ior(nf90\_clobber, nf90\_cdf5)} \\
  neither                     & \code{nf90\_clobber} (classic NetCDF-3) \\
\end{longtable}

The HDF5 compression filter is configured via \code{ncfilter} (default:
gzip, code 1), \code{ncfilter\_nparams}, and \code{ncfilter\_params}
(up to 8 integer parameters, first entry defaults to compression level 1).

%------------------------------------------------------------------------------
\subsection{Terrain and Surface Processing}
\label{subsec:dynparam_terrain}

\begin{longtable}{L{4cm} L{3cm} L{7cm}}
  \toprule
  \textbf{Name} & \textbf{Default} & \textbf{Description} \\
  \midrule
  \endhead
  \bottomrule
  \endfoot
  \code{domname}    & ---          & 64-character domain name (output file prefix). \\
  \code{prestr}     & ---          & 64-character pre-processing string prefix. \\
  \code{dirter}     & \texttt{../../Input} & Terrain input directory. \\
  \code{inpter}     & \texttt{../DATA}     & Terrain auxiliary data directory. \\
  \code{dirglob}    & \texttt{../../Input} & Global analysis input directory. \\
  \code{inpglob}    & \texttt{../DATA}     & Global auxiliary data directory. \\
  \code{dirout}     & ---          & Model output directory. \\
  \code{moist\_filename} & \texttt{moist.nc} & Initial soil moisture file name. \\
  \code{smthbdy}    & \code{.false.} & Apply extra smoothing at lateral boundaries. \\
  \code{nsmthbdy}   & 5 (min)      & Number of boundary smoothing passes ($\geq5$
    when \code{smthbdy} is true). \\
  \code{lresamp}    & \code{.false.} & Perform a first resampling pass before
    final terrain interpolation. \\
  \code{ismthlev}   & 1 & Terrain smoothness level. \\
  \code{h2opct}     & 50.0 & Minimum land-fraction percentage to classify a cell
    as water [$\%$]. \\
  \code{h2ohgt}     & \code{.false.} & Allow water grid cells to carry elevation. \\
  \code{roidem}     & $0.5\sqrt{2}$ & Interpolation search radius [grid spacings]
    for terrain resampling. \\
  \code{lakedpth}   & \code{.false.} & Create lake bathymetry for the Hostetler
    lake model. \\
  \code{lsmoist}    & \code{.false.} & Generate initial soil moisture dataset. \\
  \code{fudge\_lnd}, \code{fudge\_lnd\_s} & \code{.false.} &
    Override land-use classification for grid / sub-grid. \\
  \code{fudge\_tex}, \code{fudge\_tex\_s} & \code{.false.} &
    Override soil texture classification for grid / sub-grid. \\
  \code{fudge\_lak}, \code{fudge\_lak\_s} & \code{.false.} &
    Override lake mask for grid / sub-grid. \\
  \code{nveg}       & 22 & Number of land-use vegetation categories. \\
\end{longtable}

%------------------------------------------------------------------------------
\subsection{Fixed-Size Domain Parameters}
\label{subsec:dynparam_fixed}

\begin{longtable}{L{4cm} L{3cm} L{7cm}}
  \toprule
  \textbf{Name} & \textbf{Value} & \textbf{Description} \\
  \midrule
  \endhead
  \bottomrule
  \endfoot
  \code{mpy}       & 12  & Months per year. \\
  \code{ntex}      & 17  & Number of soil texture categories (STATSGO/USDA). \\
  \code{nats}      & 12  & Number of soil classes ($\text{ntex}-5$). \\
  \code{ndpmax}    & 200 & Maximum number of depth levels in the lake model. \\
  \code{nspi}      & 19  & Number of spectral bins in the solar radiation scheme. \\
  \code{npgwlev}   & 17  & Number of pressure levels for pseudo-global-warming
    perturbations. \\
  \code{num\_soil\_layers} & 3 (10 with CLM45) &
    Number of soil layers. \\
  \code{scalez}    & 0.025\,m (CLM45 only) &
    Soil layer thickness discretisation for CLM 4.5. \\
  \code{pthsep}    & \texttt{/} & OS path separator character. \\
\end{longtable}

%------------------------------------------------------------------------------
\subsection{Ensemble and Perturbation Parameters}
\label{subsec:dynparam_ensemble}

\begin{longtable}{L{4.5cm} L{3cm} L{6.5cm}}
  \toprule
  \textbf{Name} & \textbf{Default} & \textbf{Description} \\
  \midrule
  \endhead
  \bottomrule
  \endfoot
  \code{ensemble\_run}      & \code{.false.} &
    Enable ensemble mode (reads \code{perturbparam} namelist). \\
  \code{lperturb\_topo}     & \code{.false.} & Perturb topography. \\
  \code{lperturb\_ts}       & \code{.false.} & Perturb surface temperature IC. \\
  \code{lperturb\_ps}       & \code{.false.} & Perturb surface pressure IC. \\
  \code{lperturb\_t}        & \code{.false.} & Perturb atmospheric temperature IC. \\
  \code{lperturb\_q}        & \code{.false.} & Perturb specific humidity IC. \\
  \code{lperturb\_u}        & \code{.false.} & Perturb zonal wind IC. \\
  \code{lperturb\_v}        & \code{.false.} & Perturb meridional wind IC. \\
  \code{perturb\_frac\_*}   & $10^{-3}$      &
    Fractional perturbation amplitude for each field (all default to
    \code{d\_r1000} $=0.001$). \\
\end{longtable}

%------------------------------------------------------------------------------
\subsection{Public Procedures}
\label{subsec:dynparam_procs}

\subsubsection{\texttt{initparam(filename, ierr)}}

The primary initialisation routine. Reads the model namelist file
sequentially and populates all module-level variables.

\begin{longtable}{L{3cm} L{3cm} L{8cm}}
  \toprule
  \textbf{Argument} & \textbf{Intent} & \textbf{Description} \\
  \midrule
  \endhead
  \bottomrule
  \endfoot
  \code{filename} & \code{in} &
    Path to the RegCM namelist file (typically \texttt{regcm.in}). \\
  \code{ierr}     & \code{out} &
    Error code. 0 on success; positive values indicate which namelist
    group failed (see table below). \\
\end{longtable}

\textbf{Namelist groups read and corresponding error codes:}

\begin{longtable}{L{3cm} L{2cm} L{9cm}}
  \toprule
  \textbf{Namelist group} & \textbf{\texttt{ierr}} & \textbf{Contents} \\
  \midrule
  \endhead
  \bottomrule
  \endfoot
  \code{dimparam}      & 2 & \code{iy}, \code{jx}, \code{kz}, \code{dsmax},
    \code{dsmin}, \code{nsg}, \code{njxcpus}, \code{niycpus}. \\
  \code{coreparam}     & 3 & \code{idynamic}. \\
  \code{referenceatm}  & --- & \code{base\_state\_pressure}, \code{logp\_lrate}
    (only when \code{idynamic}~=~2). \\
  \code{molochparam}   & --- & Moloch NH settings
    (only when \code{idynamic}~=~3). \\
  \code{ncfilters}     & --- & NetCDF compression parameters
    (only with \code{NETCDF4\_HDF5+NETCDF4\_COMPRESS}). \\
  \code{geoparam}      & 4 & Projection, \code{ds}, \code{ptop}, center
    coordinates, true latitudes, band/CRM flags. \\
  \code{terrainparam}  & 5 & Terrain processing switches and directory paths. \\
  \code{debugparam}    & --- & \code{debug\_level}, \code{dbgfrq}
    (optional; failure ignored). \\
  \code{boundaryparam} & --- & Sponge zone and nudging (optional). \\
  \code{globdatparam}  & 6 & Dataset types, dates, directories,
    \code{ibdyfrq}, calendar, \code{ensemble\_run}. \\
  \code{cmip6param}    & 7 & CMIP6 model/experiment identifiers
    (only when \code{dattyp} or \code{ssttyp} is \texttt{CMIP6}). \\
  \code{cmip7param}    & 7 & CMIP7 identifiers (conditional). \\
  \code{pmip4param}    & 7 & PMIP4 identifiers (conditional). \\
  \code{perturbparam}  & 8 & Perturbation flags and fractions
    (only when \code{ensemble\_run}~=~\code{.true.}). \\
  \code{clm\_regcm}    & --- & CLM 4.5 feature flags
    (only with \code{CLM45}; failure ignored). \\
\end{longtable}

After all namelists are read, \code{initparam} derives the convenience
dimensions (\code{iym1}, \code{kzp1}, etc.), the staggered-grid index
ranges, the sub-grid dimensions, the calendar scalars (\code{dayspy},
\code{dpd}, \code{vernal\_equinox}), and the global date variables
(\code{globidate1}, \code{globidate2}). It then calls
\code{init\_random\_number\_generator}.

\subsubsection{\texttt{init\_random\_number\_generator}}

Seeds the Fortran intrinsic random number generator with a reproducible but
rank-dependent seed:
\[
  \text{seed}_i = \texttt{regcm\_magick} + \texttt{myid} \times jx \times iy
                + 37 \times (i-1)
\]
where \code{regcm\_magick} = 19\,590\,209 (a fixed salt constant). The
rank-dependent term ensures that different MPI processes draw from
uncorrelated sequences, which is required for ensemble perturbations.

\subsubsection{\texttt{init\_gpu\_random\_number\_generator} (OpenACC only)}

Seeds the cuRAND XORWOW pseudo-random generator on the GPU device using the
same seed formula (adapted to \code{ik8} arithmetic). Uses the \code{curand}
and \code{cutensorex} modules; only compiled and exposed when the
\code{OPENACC} preprocessor flag is set.

\subsubsection{\texttt{init\_fnestparam(filename, coarse\_outdir, coarse\_domname)}}

Reads the \code{fnestparam} namelist group from \code{filename} to retrieve
the output directory and domain name of the coarse-resolution parent domain
in a one-way nested simulation. Silent on failure (defaults to empty strings).

\subsubsection{\texttt{init\_globwindow(filename, lat0, lon0, lat1, lon1)}}

Reads the \code{globwindow} namelist group to define a geographic bounding
box for extracting a subset of global input data. Defaults to
$(0, 0, 0, 0)$ (global) on failure.

%------------------------------------------------------------------------------
\subsection{Design Notes}
\label{subsec:dynparam_design}

\textbf{Global mutable state.} The deliberate use of public module-level
variables rather than a configuration derived type makes \code{mod\_dynparam}
a global mutable singleton. This is idiomatic Fortran~90/95 practice and
avoids the need to thread a configuration object through every subroutine
call. It does mean that the entire configuration is implicitly shared across
all threads in an OpenMP context; the model serialises configuration
access by ensuring \code{initparam} completes before any parallel region.

\textbf{Computed index ranges.} The large number of named index bounds
(\code{ice1}--\code{jcii2}, ghost-cell ranges, etc.) is a deliberate
performance optimisation: it prevents repeated subtraction and boundary
checks inside inner loops, and it makes loop bounds self-documenting.
The naming convention --- \code{i}~=~I direction (latitude), \code{j}~=~J
direction (longitude), \code{c}~=~cross, \code{d}~=~dot, \code{e}~=~external,
\code{i} suffix~=~internal, \code{ii}~=~doubly internal, \code{g\{a,b,c\}}~=~ghost
layers A/B/C, \code{sl}~=~semi-Lagrangian stencil --- is consistent
throughout the codebase.

\textbf{Sponge zone scaling.} Rather than requiring users to specify the
sponge width in grid points for every resolution, \code{initparam} scales
a reference width (6 points at 50\,km) to the actual grid spacing. This
ensures that the physical width of the sponge zone (in km) is approximately
constant across resolutions. The user can override the scaled value via the
\code{boundaryparam} namelist.

\textbf{NetCDF format selection at compile time.} The \code{iomode} variable
is a compile-time constant (though declared as a variable to allow
initialisation from NetCDF constants). This means the output format cannot
be changed at runtime; switching between NetCDF-3 classic, CDF-5, and
NetCDF4/HDF5 requires a rebuild. This is consistent with the requirement that
all output files within a simulation use the same format.

\textbf{Dependency on \texttt{mpi\_f08}.} Only \code{type(mpi\_comm)} is
imported; no MPI procedures are called. The communicator variable
\code{mycomm} is set by the parallel initialisation code in the main program,
not by \code{initparam}. This keeps \code{mod\_dynparam} free of MPI
call semantics while still being able to store and pass the communicator
handle.

%==============================================================================
\section{Spherical Geometry --- \texttt{mod\_earth.F90}}
\label{sec:earth}
%==============================================================================

\subsection{Overview}

\code{mod\_earth} provides the spherical-geometry primitives used throughout
the RegCM pre-processing and interpolation pipeline. It supplies three
categories of functionality:

\begin{itemize}
  \item \textbf{Great-circle distance functions} --- two complementary
        implementations (\code{gcdist\_simple} and \code{gcdist}) that
        compute arc lengths on a sphere, both overloaded for single and
        double precision and annotated for OpenACC GPU execution.
  \item \textbf{Geodetic-to-Cartesian coordinate conversion} (\code{ll2xyz})
        --- five overloaded variants covering scalar pairs, 3-element array
        output, 1-D point clouds, structured 1-D latitude/longitude vectors,
        and full 2-D curvilinear grids.
  \item \textbf{Global-to-regional window extraction} (\code{get\_window},
        \code{global\_domain}) --- a derived type and overloaded subroutine
        that compute the index range on a global regular grid that fully
        encloses a (potentially curvilinear) regional model domain, handling
        all longitude-convention cases (0--360 vs.\ $-180$--180), dateline
        crossing, Greenwich meridian crossing, band mode, and polar domains.
\end{itemize}

The module is \code{private} by default; all public symbols are explicitly
listed.

\subsection{Dependencies}

\begin{lstlisting}
mod_earth
    |-- mod_realkinds   (rk4, rk8, rkx)
    |-- mod_intkinds    (ik4)
    |-- mod_constants   (degrad, erkm, mathpi, d_two)
\end{lstlisting}

\subsection{Private Constants}

\begin{longtable}{L{3cm} L{3cm} L{8cm}}
  \toprule
  \textbf{Name} & \textbf{Value} & \textbf{Description} \\
  \midrule
  \endhead
  \bottomrule
  \endfoot
  \code{mindist} & $10^{-10}$ (\code{rk8}) &
    Minimum distance floor applied in \code{gcdist} to prevent a zero
    return value when two points are numerically coincident. \\
\end{longtable}

%------------------------------------------------------------------------------
\subsection{Derived Type: \texttt{global\_domain}}
\label{subsec:earth_globaldom}

\begin{lstlisting}
type global_domain
  integer(ik4) :: global_ni
  integer(ik4) :: global_nj
  integer(ik4) :: ntiles
  integer(ik4), dimension(2) :: ni
  integer(ik4), dimension(2) :: igstart
  integer(ik4), dimension(2) :: igstop
  integer(ik4) :: nj
  integer(ik4) :: jgstart
  integer(ik4) :: jgstop
end type global_domain
\end{lstlisting}

\code{global\_domain} is the output descriptor produced by \code{get\_window}.
It describes how the global regular input grid must be sliced to cover the
regional model domain. Because the domain may straddle the dateline or the
Greenwich meridian, the I (longitude) dimension is split into up to two
contiguous \emph{tiles}; the J (latitude) dimension is always a single
contiguous range.

\begin{longtable}{L{4cm} L{3cm} L{7cm}}
  \toprule
  \textbf{Field} & \textbf{Type} & \textbf{Description} \\
  \midrule
  \endhead
  \bottomrule
  \endfoot
  \code{global\_ni} & \code{integer(ik4)} &
    Total number of longitude points in the global source grid. \\
  \code{global\_nj} & \code{integer(ik4)} &
    Total number of latitude points in the global source grid. \\
  \code{ntiles}     & \code{integer(ik4)} &
    Number of I-direction tiles required: 1 (domain fits in a single
    contiguous longitude block) or 2 (domain wraps around a longitude
    boundary). \\
  \code{igstart(2)} & \code{integer(ik4)} &
    Start longitude index of each tile in the global grid. The second
    element is 0 when \code{ntiles}~=~1. \\
  \code{igstop(2)}  & \code{integer(ik4)} &
    End longitude index of each tile. \\
  \code{ni(2)}      & \code{integer(ik4)} &
    Width of each tile: \code{igstop(i) - igstart(i) + 1}. \\
  \code{jgstart}    & \code{integer(ik4)} &
    Start latitude index in the global grid (clamped to $[1, \text{global\_nj}]$). \\
  \code{jgstop}     & \code{integer(ik4)} &
    End latitude index (clamped). \\
  \code{nj}         & \code{integer(ik4)} &
    Height of the window: \code{jgstop - jgstart + 1}. \\
\end{longtable}

%------------------------------------------------------------------------------
\subsection{Great-Circle Distance Functions}
\label{subsec:earth_gcdist}

\subsubsection{\texttt{gcdist\_simple} --- Arc Length in Kilometres}

\textbf{Generic interface overloaded for \code{rk4} and \code{rk8}.}

\begin{lstlisting}
pure real(rk8) function gcdist_simple_r8(lat1,lon1,lat2,lon2)
pure real(rk4) function gcdist_simple_r4(lat1,lon1,lat2,lon2)
\end{lstlisting}

\begin{longtable}{L{2.5cm} L{2.5cm} L{3cm} L{6cm}}
  \toprule
  \textbf{Argument} & \textbf{Intent} & \textbf{Type} & \textbf{Description} \\
  \midrule
  \endhead
  \bottomrule
  \endfoot
  \code{lat1}, \code{lon1} & \code{in} & \code{real(rk4/rk8)} &
    Geodetic latitude and longitude of point 1 [degrees]. \\
  \code{lat2}, \code{lon2} & \code{in} & \code{real(rk4/rk8)} &
    Geodetic latitude and longitude of point 2 [degrees]. \\
\end{longtable}

\textbf{Returns:} great-circle arc length in kilometres (\code{rk8} or
\code{rk4} respectively).

\textbf{Algorithm:} Spherical law of cosines. All intermediate computation
is performed in \code{rk8} regardless of the input precision:
\[
  c = \sin\phi_1\sin\phi_2 + \cos\phi_1\cos\phi_2\cos(\lambda_1-\lambda_2)
\]
\[
  d = R_\text{km}\arccos\!\bigl(\max(-1,\,\min(1,\,c))\bigr)
\]
The \code{max}/\code{min} clamp prevents \code{acos} domain errors from
floating-point rounding when the two points are coincident or antipodal.

\textbf{GPU availability:} both variants are annotated \code{!\$acc routine seq},
making them callable from OpenACC device kernels.

\subsubsection{\texttt{gcdist} --- Arc Length in Grid Spacings}

\begin{lstlisting}
pure real(rkx) function gcdist(ds, lat1, lon1, lat2, lon2)
\end{lstlisting}

\begin{longtable}{L{2.5cm} L{2.5cm} L{3cm} L{6cm}}
  \toprule
  \textbf{Argument} & \textbf{Intent} & \textbf{Type} & \textbf{Description} \\
  \midrule
  \endhead
  \bottomrule
  \endfoot
  \code{ds}           & \code{in} & \code{real(rkx)} &
    Grid spacing [km] used to normalise the result. \\
  \code{lat1}, \code{lon1} & \code{in} & \code{real(rkx)} &
    Latitude and longitude of point 1 [degrees]. \\
  \code{lat2}, \code{lon2} & \code{in} & \code{real(rkx)} &
    Latitude and longitude of point 2 [degrees]. \\
\end{longtable}

\textbf{Returns:} great-circle distance normalised by \code{ds}
[dimensionless grid spacings], floored at \code{mindist} ($10^{-10}$).

\textbf{Algorithm:} Vincenty formula (numerically stable for all
separations, including near-antipodal points):
\[
  y = \sqrt{(\cos\phi_2\sin\Delta\lambda)^2 +
            (\cos\phi_1\sin\phi_2 - \sin\phi_1\cos\phi_2\cos\Delta\lambda)^2}
\]
\[
  x = \sin\phi_1\sin\phi_2 + \cos\phi_1\cos\phi_2\cos\Delta\lambda
\]
\[
  d = \max\!\left(\frac{R_\text{km}\,\text{atan2}(y,x)}{ds},\;
                  \text{mindist}\right)
\]
Using \code{atan2} rather than \code{acos} avoids the clamping requirement
and gives better accuracy near zero separation. This function is the one
used in the kd-tree search (\code{mod\_kdinterp}) to convert geographic
distances to grid-point units.

\textbf{GPU availability:} annotated \code{!\$acc routine seq}.

%------------------------------------------------------------------------------
\subsection{Geodetic-to-Cartesian Conversion: \texttt{ll2xyz}}
\label{subsec:earth_ll2xyz}

The generic interface \code{ll2xyz} converts geographic coordinates
(latitude, longitude in degrees) to unit-sphere Cartesian coordinates
$(x, y, z)$ using the convention:
\[
  x = \cos\phi\,\sin\lambda, \quad
  y = \sin\phi, \quad
  z = \cos\phi\,\cos\lambda
\]
where $\phi$ is latitude and $\lambda$ is longitude, both in radians.
All variants clamp latitude to $[-89.999°, +89.999°]$ to avoid exact pole
singularities, and promote all computation to \code{rk8} regardless of
input precision. Output is always \code{real(rk8)}.

Five overloaded procedures are provided:

\begin{longtable}{L{4.5cm} L{9.5cm}}
  \toprule
  \textbf{Procedure} & \textbf{Signature and description} \\
  \midrule
  \endhead
  \bottomrule
  \endfoot
  \code{ll2xyz\_values} &
    \code{(lat, lon, x, y, z)}: scalar inputs; three scalar \code{rk8}
    outputs. The most explicit form. \\
  \code{ll2xyz\_array} &
    \code{(lat, lon, x(3))}: scalar inputs; output as a 3-element vector
    \code{x(1:3)} = $(x,y,z)$. Suitable for passing to kd-tree point
    insertion routines. \\
  \code{ll2xyz\_1d} &
    \code{(ni, lat(ni), lon(ni), x(3,ni))}: 1-D arrays of $N$ points;
    output is a $3\times N$ matrix. Uses \code{do concurrent} for
    compiler-assisted parallelism. \\
  \code{ll2xyz\_arrays} &
    \code{(lat(:), lon(:), x(3,:))}: separate 1-D latitude and longitude
    vectors forming a structured grid; output is $3\times(n_\text{lat}
    \times n_\text{lon})$, packed in row-major order $n = (j-1)n_\text{lon}+i$.
    Returns silently if the output array has the wrong shape. Uses
    \code{do concurrent}. \\
  \code{ll2xyz\_grid} &
    \code{(lat(:,:), lon(:,:), x(3,:))}: 2-D curvilinear grid; output is
    $3\times(n_i\times n_j)$, same packing. Validates shape consistency
    between \code{lat}, \code{lon}, and \code{x}; returns silently on
    mismatch. Uses \code{do concurrent}. \\
\end{longtable}

%------------------------------------------------------------------------------
\subsection{\texttt{longitude\_circle} --- Zonal Circumference}
\label{subsec:earth_loncircle}

\begin{lstlisting}
pure real(rk8) function longitude_circle(lat) result(er)
  !$acc routine seq
  real(rk8), intent(in) :: lat
\end{lstlisting}

Returns the circumference of the circle of latitude \code{lat} [degrees]
in kilometres:
\[
  C(\phi) = 2\pi\,R_\text{km}\cos\phi
\]
Used in pre-processing to determine the zonal extent of the domain relative
to the full equatorial circumference, e.g.\ to detect whether a band-mode
domain effectively spans the globe. Annotated \code{!\$acc routine seq}.

%------------------------------------------------------------------------------
\subsection{\texttt{get\_window} --- Global Grid Window Extraction}
\label{subsec:earth_getwindow}

\textbf{Generic interface overloaded for \code{rk4} and \code{rk8} global
grid coordinate arrays.}

\begin{lstlisting}
subroutine get_window_r8(glat, glon, xlat, xlon, i_band, domain)
subroutine get_window_r4(glat, glon, xlat, xlon, i_band, domain)
\end{lstlisting}

\begin{longtable}{L{3cm} L{2.5cm} L{3.5cm} L{5cm}}
  \toprule
  \textbf{Argument} & \textbf{Intent} & \textbf{Type} & \textbf{Description} \\
  \midrule
  \endhead
  \bottomrule
  \endfoot
  \code{glat(:)} & \code{in} & \code{real(rk8/rk4)} &
    Latitudes of the global regular grid, monotone, ordered South-to-North
    or North-to-South. \\
  \code{glon(:)} & \code{in} & \code{real(rk8/rk4)} &
    Longitudes of the global regular grid, either $0$--$360$ or
    $-180$--$180$. \\
  \code{xlat(:,:)} & \code{in} & \code{real(rkx)} &
    2-D latitude array of the regional model domain, ordered
    South-to-North in J. \\
  \code{xlon(:,:)} & \code{in} & \code{real(rkx)} &
    2-D longitude array of the regional domain in the range
    $-180$--$180$. \\
  \code{i\_band} & \code{in} & \code{integer(ik4)} &
    Tropical band mode flag from \code{mod\_dynparam}. When 1, the
    full zonal extent of the global grid is always selected. \\
  \code{domain} & \code{out} & \code{type(global\_domain)} &
    Populated window descriptor (see above). \\
\end{longtable}

\subsubsection{Algorithm}

The subroutine handles five distinct longitude-convention cases, determined
by inspecting \code{glon} and \code{xlon}:

\begin{longtable}{L{4cm} L{10cm}}
  \toprule
  \textbf{Case} & \textbf{Handling} \\
  \midrule
  \endhead
  \bottomrule
  \endfoot
  Band mode (\code{i\_band}~=~1) &
    Full zonal extent: \code{igstart(1)=1}, \code{igstop(1)=gi},
    \code{ntiles=1}. No longitude search needed. \\
  Global grid 0--360, domain crosses Greenwich ($\lambda=0$) &
    The western edge (\code{xlon(1,:)}) contains negative longitudes
    and the eastern edge (\code{xlon(xi,:)}) contains positive ones.
    Converts both edge longitudes to 0--360 and sets
    \code{ntiles=2}: tile~1 runs from the western limit to \code{gi},
    tile~2 from 1 to the eastern limit. \\
  Global grid 0--360, domain entirely within 0--360 &
    Maps negative \code{xlon} values to 0--360 via an \code{allocate}d
    temporary array, then computes a single tile. \\
  Global grid $-180$--180, domain does not cross $\pm180$ &
    Simple min/max search; single tile. \\
  Global grid $-180$--180, domain crosses the dateline ($\pm180$) &
    Detected by a large longitude jump ($>350°$) along the western or
    eastern boundary column, or by the western column being entirely
    positive while the eastern column is entirely negative. Sets
    \code{ntiles=2}: tile~1 covers from the limit to \code{gi},
    tile~2 from 1 to the other limit. \\
\end{longtable}

After the longitude tile computation, three polar cases are checked by the
internal functions \code{has\_north\_pole} and \code{has\_south\_pole}:

\begin{longtable}{L{4cm} L{10cm}}
  \toprule
  \textbf{Case} & \textbf{Handling} \\
  \midrule
  \endhead
  \bottomrule
  \endfoot
  Domain contains the North Pole &
    Sets \code{jgstop=gj} (or \code{jgstart=1} for a N-to-S grid),
    forces \code{ntiles=1} and selects the full zonal extent, because
    a polar domain always wraps zonally. \\
  Domain contains the South Pole &
    Symmetric treatment: \code{jgstart=1} (S-to-N) extended to the
    southernmost row, full zonal extent. \\
  Normal (no pole inside) &
    \code{jgstart}/\code{jgstop} from \code{minlat}/\code{maxlat},
    with a $\pm2$--$\pm3$ grid-point margin. For N-to-S global grids,
    the index arithmetic is reflected about the grid midpoint. \\
\end{longtable}

A final clipping step after both longitude and latitude computations:
\begin{itemize}
  \item If \code{igstart(1) < 1} with \code{ntiles=1}, the window wraps
        westward past the start of the global grid; \code{ntiles} is
        promoted to 2 and \code{igstart(1)} is shifted by \code{gi}.
  \item If \code{igstop(1) > gi} with \code{ntiles=1}, the window wraps
        eastward; similarly promoted to 2 tiles.
  \item \code{jgstart} and \code{jgstop} are clamped to $[1, \text{gj}]$.
\end{itemize}

\subsubsection{Internal Functions}

Both \code{get\_window\_r4} and \code{get\_window\_r8} contain two identical
internal functions (internal to each subroutine via \code{contains}):

\begin{longtable}{L{5cm} L{9cm}}
  \toprule
  \textbf{Function} & \textbf{Description} \\
  \midrule
  \endhead
  \bottomrule
  \endfoot
  \code{has\_north\_pole(l, i)} &
    Returns \code{.true.} if the latitude column at index \code{i}
    (the domain mid-column) contains values that decrease from south
    to north, indicating the North Pole is enclosed. Early exit if all
    latitudes in the column are negative. \\
  \code{has\_south\_pole(l, i)} &
    Symmetric: returns \code{.true.} if the mid-column latitudes
    decrease from north to south (South Pole enclosed), with early
    exit if all latitudes are positive. \\
\end{longtable}

\textbf{Note:} The polar-detection heuristic checks the middle column of
the domain (\code{xi/2}). It assumes that the domain is convex enough that
the central meridian is representative. Highly skewed or rotated domains
near a pole may in principle defeat the heuristic, but such configurations
are not used in standard RegCM setups.

%------------------------------------------------------------------------------
\subsection{Design Notes}
\label{subsec:earth_design}

\textbf{Double-precision arithmetic throughout.} Even when the generic
interface selects the \code{rk4} overload of \code{gcdist\_simple} or
\code{get\_window}, all trigonometric and distance calculations are performed
in \code{rk8}. This is important because trigonometric functions of
degree-valued arguments that have been converted to radians can lose
significant bits in single precision, especially near the poles.

\textbf{Vincenty vs.\ law-of-cosines.} \code{gcdist} uses the two-argument
arctangent form (Vincenty) while \code{gcdist\_simple} uses the law-of-cosines
form with a clamp. The Vincenty form is more accurate for small separations
(the kd-tree use case, where nearby points are searched), while the cosine
form is simpler and sufficient when only order-of-magnitude distances are
needed (e.g.\ domain-size checks).

\textbf{Grid-spacing normalisation in \texttt{gcdist}.} Returning the
distance in units of \code{ds} rather than kilometres allows the kd-tree
radius search (\code{mod\_kdinterp}) to express its search radius as a pure
number of grid spacings, which is resolution-independent and directly
comparable to the interpolation weight cutoffs.

\textbf{Two-tile window design.} The \code{global\_domain} type's
\code{ntiles} field abstracts the dateline/Greenwich wrap-around problem
into a single descriptor that the calling pre-processing code can handle
with a simple loop over \code{ntiles}. This avoids duplicating the
longitude-convention logic in every module that reads global input data.

\textbf{Margin of $\pm2$--$\pm3$ grid points.} The window bounds are
expanded by 2 points on the ``near'' side and 3 on the ``far'' side to
ensure that interpolation stencils at the domain boundary always have
global data available, even for high-order schemes.

%==============================================================================
\section{Spherical Barycentric Coordinates --- \texttt{mod\_spbarcoord.F90}}
\label{sec:spbarcoord}
%==============================================================================

\subsection{Overview}

\code{mod\_spbarcoord} computes the interpolation weights of a point on the
sphere with respect to a set of surrounding points, using the spherical
barycentric (mean value) coordinates of \citet{langer2006}. Given the query
point $\mathbf{p}$ and $n$ vertices $\mathbf{v}_1, \ldots, \mathbf{v}_n$ on
the unit sphere, it returns weights $\lambda_i \ge 0$ with
$\sum_i \lambda_i = 1$ such that $\sum_i \lambda_i \mathbf{v}_i$ points in the
direction of $\mathbf{p}$. A field known at the vertices is then
interpolated at $\mathbf{p}$ as $f(\mathbf{p}) = \sum_i \lambda_i f_i$.

The only user is \code{mod\_kdinterp}, which computes the weights of every
target grid point from the source points found by a kd-tree search
(\code{compwgt\_genlin}): the 4 nearest points when the target grid is
finer than the source grid, or the points within a search radius (up to
32) when it is coarser. The module is \code{private} by default and its
only public symbol is the subroutine \code{spherical\_barycentric}.

\subsection{Dependencies}

\begin{lstlisting}
mod_spbarcoord
    |-- mod_realkinds   (rk8)
    |-- mod_intkinds    (ik4)
    |-- mod_constants   (d_zero)
\end{lstlisting}

The points are Cartesian coordinates on the unit sphere, as produced by
\code{ll2xyz} of \code{mod\_earth} (Section~\ref{subsec:earth_ll2xyz}).

\subsection{Module Data}

\begin{longtable}{L{3.5cm} L{3.5cm} L{7cm}}
  \toprule
  \textbf{Name} & \textbf{Type} & \textbf{Description} \\
  \midrule
  \endhead
  \bottomrule
  \endfoot
  \code{vpoint} & derived type &
    A vertex for sorting: its index \code{idx} (\code{integer(ik4)}) in
    the input array and its coordinates \code{v(3)} (\code{real(rk8)}). \\
  \code{notzero} & \code{real(rk8)}, parameter &
    \code{epsilon(1.0\_rk8)}, the floor used to avoid divisions by zero. \\
\end{longtable}

The module has no variables: all the work arrays are local to
\code{spherical\_barycentric}, so it can be called concurrently, e.g.\ from
OpenMP threads.

%------------------------------------------------------------------------------
\subsection{\texttt{spherical\_barycentric} --- Interpolation Weights}
\label{subsec:spbar_weights}

\begin{lstlisting}
subroutine spherical_barycentric(np, p, v, lambda)
\end{lstlisting}

\begin{longtable}{L{2.5cm} L{1.5cm} L{3.5cm} L{6.5cm}}
  \toprule
  \textbf{Argument} & \textbf{Intent} & \textbf{Type} & \textbf{Description} \\
  \midrule
  \endhead
  \bottomrule
  \endfoot
  \code{np} & \code{in} & \code{integer(ik4)} &
    Number of vertices. The routine returns immediately, without setting
    \code{lambda}, if \code{np}~=~0. \\
  \code{p(3)} & \code{in} & \code{real(rk8)} &
    Query point, unit vector. \\
  \code{v(3,np)} & \code{in} & \code{real(rk8)} &
    Vertices, unit vectors, in any order. \\
  \code{lambda(np)} & \code{out} & \code{real(rk8)} &
    Weights, in the order of the input vertices, summing to 1. \\
\end{longtable}

\subsubsection{Algorithm}

\begin{enumerate}
  \item \textbf{Ordering.} The mean value coordinates are defined for a
        polygon, so the vertices must be taken in cyclic order around the
        query point. The routine copies them into an array of
        \code{vpoint}, keeping their original index, and sorts it by the
        angle $\varphi_i$ of each vertex around $\mathbf{p}$ in the plane
        tangent to the sphere at $\mathbf{p}$:
        \[
          \varphi_i = \text{atan2}(\mathbf{v}_i\cdot\mathbf{e}_2,\;
                                   \mathbf{v}_i\cdot\mathbf{e}_1)
        \]
        where $(\mathbf{e}_1, \mathbf{e}_2)$ is an orthonormal basis of
        the tangent plane: $\mathbf{e}_1$ is the normalised cross product
        of $\mathbf{p}$ with the coordinate axis least aligned with it, and
        $\mathbf{e}_2 = \mathbf{p}\times\mathbf{e}_1$. The resulting polygon
        is star-shaped as seen from $\mathbf{p}$, whatever the input order
        and the spatial arrangement of the vertices.
  \item \textbf{Angles.} For each vertex $i$ of the ordered polygon it
        computes
        \begin{itemize}
          \item $\theta_i$, the angle between $\mathbf{p}$ and
                $\mathbf{v}_i$, i.e.\ the great-circle distance from the
                query point to the vertex;
          \item $\alpha_i$, the angle at $\mathbf{p}$ between the
                great-circle arcs to $\mathbf{v}_i$ and to
                $\mathbf{v}_{i+1}$ (cyclically, $\mathbf{v}_{n+1} =
                \mathbf{v}_1$), computed as the angle between the normals
                $\mathbf{p}\times\mathbf{v}_i$ and
                $\mathbf{p}\times\mathbf{v}_{i+1}$ of the two arcs.
        \end{itemize}
        Both use \code{acos} of the normalised dot product, clamped to
        $[-1, 1]$.
  \item \textbf{Weights.} The unnormalised weight of vertex $i$ is
        \[
          w_i = \frac{\tan(\alpha_{i-1}/2) + \tan(\alpha_i/2)}
                     {\sin\theta_i}
        \]
        (with $\alpha_0 = \alpha_n$), the spherical form of the mean value
        coordinates. $\sin\theta_i$ is replaced by \code{notzero} when it is
        smaller, and the weights are divided by their sum (also floored at
        \code{notzero}). Each weight is stored at the original index of its
        vertex.
  \item \textbf{Normalisation check.} The weights are divided once more
        by their sum, if it is larger than \code{notzero}, to make the sum
        exactly one after rounding.
\end{enumerate}

A query point that coincides with a vertex would give $\theta_i = 0$. The
caller handles this case before calling the routine:
\code{compwgt\_genlin} in \code{mod\_kdinterp} gives weight 1 to a source
point at zero distance.

\subsubsection{Internal Procedures}

\begin{longtable}{L{4.5cm} L{9.5cm}}
  \toprule
  \textbf{Procedure} & \textbf{Description} \\
  \midrule
  \endhead
  \bottomrule
  \endfoot
  \code{set\_cyclic\_order} &
    Fills the \code{vpoint} array, computes the angles $\varphi_i$ and
    sorts the vertices by angle with an insertion sort (the number of
    vertices is small, at most 32 in \code{mod\_kdinterp}). \\
  \code{compute\_angles} &
    Computes $\theta_i$ and $\alpha_i$ for the ordered vertices. \\
  \code{angle\_between(x, y)} &
    Angle between two vectors, $\arccos$ of the clamped normalised dot
    product (the norm product is floored at \code{notzero}). Pure. \\
  \code{vecprod(x, y, z)} &
    Cross product $\mathbf{z} = \mathbf{x}\times\mathbf{y}$. \\
  \code{dotprod}, \code{norma}, \code{norma2} &
    Dot product, Euclidean norm and squared norm of 3-vectors. Pure. \\
\end{longtable}

%------------------------------------------------------------------------------
\subsection{Design Notes}
\label{subsec:spbar_design}

\textbf{Mean value coordinates on the sphere.} Planar barycentric
coordinates are defined only for triangles; mean value coordinates extend
them to polygons with any number of vertices, and the spherical version
works directly with points on the sphere, without a map projection. The
weights are positive inside a convex polygon, vary smoothly with the query
point, and reproduce the query point, so a field that is linear in the
Cartesian coordinates is interpolated exactly, up to the normalisation.

\textbf{Ordering around the query point.} Sorting by the angle in the
tangent plane at $\mathbf{p}$ always gives a cyclic order. In a test with
query points inside the polygon of their neighbours, the weights
reproduce the query point in every case, also with vertices scattered over
a disk as in the radius search of \code{mod\_kdinterp}.

\textbf{Query point outside the polygon.} If the query point is outside
the polygon of its neighbours (one angular gap between consecutive vertices
larger than $\pi$), no set of positive weights can reproduce it. The angle
$\alpha_i$ across the gap is computed without sign, i.e.\ as
$2\pi - \text{gap}$, so the weights are still positive and sum to one: the
result is a weighted average of the neighbours, not an extrapolation.

%==============================================================================
\section{Map Projections --- \texttt{mod\_projections.F90}}
\label{sec:projections}
%==============================================================================

\subsection{Overview}

\code{mod\_projections} maps between the model grid indices $(i, j)$ and
geographic coordinates (latitude, longitude) for the six grids RegCM
supports, computes the map factor, and rotates winds between the
geographic frame (eastward, northward) and the grid frame (along $i$, along
$j$). The projection is selected by the namelist variable \code{iproj}:

\begin{longtable}{L{2cm} L{4cm} L{8cm}}
  \toprule
  \textbf{Code} & \textbf{Grid} & \textbf{Notes} \\
  \midrule
  \endhead
  \bottomrule
  \endfoot
  \code{LAMCON} & Lambert conformal &
    Secant cone with two true latitudes, or tangent cone when they are
    equal (within $0.1^\circ$). \\
  \code{POLSTR} & Polar stereographic &
    True scale at \code{clat}. \\
  \code{NORMER} & Normal Mercator &
    True scale at \code{clat}. Used by the tropical band. \\
  \code{ROTMER} & Rotated Mercator &
    Mercator on a rotated sphere whose equator passes through
    (\code{plat}, \code{plon}); true scale on the rotated equator. \\
  \code{ROTLLR} & Rotated latitude--longitude &
    Regular grid in the coordinates of a sphere whose north pole is at
    (\code{plat}, \code{plon}). The only grid that accepts a spacing in
    degrees (negative \code{ds}). \\
  any other & Regular latitude--longitude &
    Fallback for an unknown code, with no wind rotation and a map factor
    of 1. \\
\end{longtable}

\noindent Note that \code{plat} and \code{plon} have different meanings for
the two rotated grids: for \code{ROTMER} they are a point of the rotated
equator (normally the domain centre, \code{plat = clat},
\code{plon = clon}), for \code{ROTLLR} they are the position of the rotated
north pole.

The module is used by the terrain program to compute the coordinates,
cell bounds and map factors of every grid (one projection object for each
staggering: cross, $u$, $v$ and dot points), by the ICBC programs to rotate
the input winds onto the grid, by the model output and the nesting code to
rotate them back, and by the GrADS post-processors.

\subsection{Dependencies}

\begin{lstlisting}
mod_projections
    |-- mod_intkinds    (ik4)
    |-- mod_realkinds   (rk8, rkx)
    |-- mod_constants   (degrad_r8, raddeg_r8, earthrad_r8, halfpi_r8,
    |                    mathpi_r8, dlowval_r8, deg00, deg45, deg90,
    |                    deg180, deg360)
    |-- mod_memutil     (getmem)
    |-- iso_fortran_env (REAL32, REAL64; with F2008)
\end{lstlisting}

\code{mod\_memutil} allocates the arrays of the wind rotation factors.

%------------------------------------------------------------------------------
\subsection{Derived Types}
\label{subsec:proj_types}

\subsubsection{\texttt{anyprojparams} --- Projection Parameters}

The input of \code{initialize}, filled by the caller from the namelist.

\begin{longtable}{L{3.5cm} L{3cm} L{7.5cm}}
  \toprule
  \textbf{Field} & \textbf{Type} & \textbf{Description} \\
  \midrule
  \endhead
  \bottomrule
  \endfoot
  \code{pcode} & \code{character(6)} & Projection code (\code{iproj}). \\
  \code{ds} & \code{real(rk8)} & Grid spacing [m]. \\
  \code{clat}, \code{clon} & \code{real(rk8)} &
    Geographic coordinates of the domain centre [degrees]. \\
  \code{plat}, \code{plon} & \code{real(rk8)} &
    \code{ROTMER}: point of the rotated equator; \code{ROTLLR}: rotated
    north pole [degrees]. \\
  \code{trlat1}, \code{trlat2} & \code{real(rk8)} &
    True latitudes of \code{LAMCON} [degrees]. \\
  \code{nlon}, \code{nlat} & \code{integer(ik4)} &
    Number of points along $i$ and $j$. \\
  \code{staggerx}, \code{staggery} & \code{logical} &
    Whether the grid is staggered by half a grid spacing along $i$ or $j$
    (the $u$, $v$ and dot grids). \\
  \code{rotparam} & \code{logical} &
    Compute the wind rotation factors. Without them the rotation routines
    must not be called. \\
\end{longtable}

\subsubsection{\texttt{regcm\_projection} --- The Projection Object}

\begin{lstlisting}
type regcm_projection
  private
  type(regcm_projdata), pointer :: p => null()
  procedure(transform),  pass(pj), pointer, public :: llij, ijll
  procedure(map_factor), pass(pj), pointer, public :: mapfac
  procedure(rotate2),    pass(pj), pointer, public :: uvrotate2, uvbkrotate2
  procedure(rotate3),    pass(pj), pointer, public :: uvrotate3, uvbkrotate3
  real(rk8), pointer, contiguous, dimension(:,:) :: f1, f2, f3, f4
contains
  procedure, public :: initialize, destruct, rotation_angle, rl00, conefac
  procedure, public :: wind_rotate, wind2_rotate
  procedure, public :: wind_antirotate, wind2_antirotate
end type regcm_projection
\end{lstlisting}

\code{initialize} points the procedure pointers to the routines of the
selected projection, so that a call such as \code{call pj\%ijll(i, j, lat,
lon)} needs no test on the projection. The private pointer \code{p} holds
the projection constants (\code{regcm\_projdata}, a private type with one
set of fields per projection). The arrays \code{f1} and \code{f2} hold
$\cos$ and $\sin$ of the wind rotation angle at every grid point, and
\code{f3} the angle itself (\code{LAMCON}, \code{POLSTR}, \code{ROTMER});
\code{f4} is not used.

%------------------------------------------------------------------------------
\subsection{Public Procedures}
\label{subsec:proj_public}

\begin{longtable}{L{5cm} L{9cm}}
  \toprule
  \textbf{Procedure} & \textbf{Description} \\
  \midrule
  \endhead
  \bottomrule
  \endfoot
  \code{initialize(pjpara)} &
    Sets up the projection from an \code{anyprojparams} and, if
    \code{rotparam} is true, computes the rotation factors at every grid
    point. \\
  \code{ijll(i, j, lat, lon)} &
    Geographic coordinates [degrees, longitude in $[-180, 180]$] of the
    (real) grid position $(i, j)$, with $i = 1, \ldots,$ \code{nlon}. \\
  \code{llij(lat, lon, i, j)} &
    Inverse: real grid position of a geographic point. \\
  \code{mapfac(xlat, xlon, xmap)} &
    Map factor at the points of the 2-D coordinate arrays. \\
  \code{wind\_rotate(u, v)}, \code{wind2\_rotate(u, v)} &
    3-D and 2-D: replace geographic wind components (eastward,
    northward) with grid components (along $i$, along $j$). \\
  \code{wind\_antirotate(u, v)}, \code{wind2\_antirotate(u, v)} &
    Inverse: grid components to geographic ones. \\
  \code{uvrotate2}, \code{uvrotate3}, \code{uvbkrotate2},
  \code{uvbkrotate3} &
    The same rotations through procedure pointers. \\
  \code{rotation\_angle(lon, lat, p)} &
    Rotation angle [radians] at the points of a regular latitude--longitude
    grid (1-D \code{lon} and \code{lat}), for the \code{pdef} file of the
    GrADS post-processors: $n(\lambda_0 - \lambda)$ for \code{LAMCON},
    $\lambda_0 - \lambda$ for \code{POLSTR}, $\delta$ for \code{ROTMER}
    (see below), and 0 otherwise. \code{ROTLLR} is not implemented; its
    callers skip it. \\
  \code{rl00(lat0, lon0)} &
    Rotated (or, for the fallback grid, geographic) coordinates of the
    first grid point, for \code{ROTLLR} and the fallback grid. \\
  \code{conefac()} &
    Cone factor $n$ of \code{LAMCON}; 1 for the other projections. \\
  \code{destruct()} &
    Releases the projection constants. \\
  \code{rounder(x, ltop)} &
    Module function: rounds \code{x} up (\code{ltop} true) or down to a
    multiple of $10^{-3}$. \\
\end{longtable}

The wind arrays are \code{pointer, contiguous} and must have the shape of
the grid of the projection (\code{nlon} $\times$ \code{nlat}), with the
first index along $i$; their bounds are taken from the arrays.

The kind parameters \code{rk4}, \code{rk8} and \code{rkx} are declared
again inside each abstract interface, with the same definitions as in
\code{mod\_realkinds}, because the interfaces cannot rely on host
association.

%------------------------------------------------------------------------------
\subsection{Grid Position and Conventions}
\label{subsec:proj_grid}

\code{initialize} computes the offset of the domain centre from the first
grid point,
\[
  c_i = \frac{n_\text{lon}}{2} - \frac{1}{2}\,[\text{not staggered in } x],
  \qquad
  c_j = \frac{n_\text{lat}}{2} - \frac{1}{2}\,[\text{not staggered in } y],
\]
and all the projections place (\code{clat}, \code{clon}) at the grid
position $(c_i + 1, c_j + 1)$, the middle of the grid: with
$n_\text{lon} = 61$ cross points, at $i = 31$. A staggered grid is shifted by
half a grid spacing with respect to the cross grid. The projections working
with the grid index (\code{LAMCON}, \code{POLSTR}, \code{NORMER},
\code{ROTMER}) receive the index $c_i + 1$, the latitude--longitude grids
the offset $c_i$. Before October 2026 the first group placed the centre one
grid point earlier, at $(c_i, c_j)$; domains created with those versions are
shifted by one grid point with respect to new ones. The namelist variables
\code{cntri} and \code{cntrj} are not used by the projection.

Below, $R$ is the Earth radius (\code{earthrad\_r8}), $\Delta s$ the grid
spacing, $\varphi$ and $\lambda$ latitude and longitude, $\varphi_c$ and
$\lambda_c$ those of the domain centre, and $h = \pm1$ the hemisphere of
the centre.

%------------------------------------------------------------------------------
\subsection{Lambert Conformal (\texttt{LAMCON})}
\label{subsec:proj_lambert}

The standard longitude $\lambda_0$ is \code{clon}. The cone factor is
\[
  n = \frac{\log\cos\varphi_1 - \log\cos\varphi_2}
           {\log\tan(45^\circ - |\varphi_1|/2) -
            \log\tan(45^\circ - |\varphi_2|/2)}
\]
for a secant cone, and $n = \sin|\varphi_1|$ for a tangent cone
($|\varphi_1 - \varphi_2| \le 0.1^\circ$). The distance from the pole, in
grid spacings, is
\[
  r(\varphi) = \frac{R}{\Delta s}\,\frac{\cos\varphi_1}{n}
    \left(\frac{\tan\bigl((90^\circ h - \varphi)/2\bigr)}
               {\tan\bigl((90^\circ h - \varphi_1)/2\bigr)}\right)^{n}
\]
and, with the pole position $(P_i, P_j)$ computed so that the centre falls
at $(c_i, c_j)$,
\[
  i = h\,\bigl(P_i + h\,r\sin\bigl(n(\lambda - \lambda_0)\bigr)\bigr),
  \qquad
  j = h\,\bigl(P_j - r\cos\bigl(n(\lambda - \lambda_0)\bigr)\bigr).
\]
\code{ijll\_lc} inverts these relations. The map factor is
\[
  m = \frac{\sin\chi_2}{\sin\chi}
      \left(\frac{\tan(\chi/2)}{\tan(\chi_2/2)}\right)^{n'}
\]
with $\chi = 90^\circ - \varphi$ the colatitude and $n'$ the cone factor
computed from the colatitudes (\code{nfac}); for a tangent cone $\chi_1$
and $\cos\chi_1$ replace $\chi_2$ and $n'$. It is 1 on the true
latitudes.

\textbf{Wind rotation.} The grid $j$ axis makes the angle
$\alpha = n(\lambda_0 - \lambda)$ with the local meridian. In the northern
hemisphere
\[
  u_g = u\cos\alpha + v\sin\alpha, \qquad
  v_g = -u\sin\alpha + v\cos\alpha,
\]
and in the southern hemisphere the same with $-\alpha$. The back rotation
is the inverse.

%------------------------------------------------------------------------------
\subsection{Polar Stereographic (\texttt{POLSTR})}
\label{subsec:proj_polar}

With the reference longitude $\lambda_r = \lambda_0 + 90^\circ$ and the
scale factor $k = 1 + h\sin\varphi_c$, which makes the grid true to scale at
\code{clat},
\[
  r(\varphi) = \frac{R}{\Delta s}\,\frac{k\cos\varphi}{1 + h\sin\varphi},
  \qquad
  i = P_i + r\cos(\lambda - \lambda_r), \qquad
  j = P_j + h\,r\sin(\lambda - \lambda_r).
\]
The map factor is $m = k/(1 + h\sin\varphi)$. The wind is rotated as for
\code{LAMCON} with $n = 1$, i.e.\ $\alpha = \lambda_0 - \lambda$, and the
sign of $\alpha$ reversed in the southern hemisphere.

%------------------------------------------------------------------------------
\subsection{Normal Mercator (\texttt{NORMER})}
\label{subsec:proj_mercator}

The longitude step is $\Delta\lambda = \Delta s/(R\cos\varphi_c)$
[radians], so that the grid is true to scale at \code{clat}, and
\[
  i = c_i + \frac{\lambda - \lambda_c}{\Delta\lambda}, \qquad
  j = c_j + \frac{\ln\tan(45^\circ + \varphi/2) -
                  \ln\tan(45^\circ + \varphi_c/2)}{\Delta\lambda}.
\]
The map factor is $m = \cos\varphi_c/\cos\varphi$. The grid axes are
aligned with the parallels and meridians, so there is no wind rotation.

%------------------------------------------------------------------------------
\subsection{Rotated Mercator (\texttt{ROTMER})}
\label{subsec:proj_rotmer}

The sphere is rotated so that (\code{plat}, \code{plon}) lies on the
rotated equator, with the rotated pole at latitude $90^\circ -$
\code{plat} and longitude \code{plon}$\,+180^\circ$. The rotated
coordinates $(\varphi', \lambda')$ of a point follow the COSMO formulas
for the rotated pole, and the grid is a Mercator grid in them, true to
scale on the rotated equator:
\[
  x = \lambda', \qquad y = \ln\tan(45^\circ + \varphi'/2), \qquad
  i = c_i + \frac{x - x_c}{\Delta}, \qquad
  j = c_j + \frac{y - y_c}{\Delta}
\]
with $\Delta = \Delta s/R$ (in degrees, as $x$ and $y$), and $(x_c, y_c)$
the coordinates of the domain centre, computed with the same formulas. The
map factor is $m = 1/\cos\varphi'$.

\textbf{Wind rotation.} \code{uvrot\_rc} computes the angle $\delta$
between the geographic and the rotated meridians, from the rotated pole
$(\varphi_p, \lambda_p)$:
\[
  \tan\delta = \frac{\cos\varphi_p\,\sin(\lambda_p - \lambda)}
                    {\sin\varphi_p\cos\varphi -
                     \cos\varphi_p\sin\varphi\cos(\lambda_p - \lambda)},
  \qquad
  u_g = u\cos\delta - v\sin\delta, \quad
  v_g = u\sin\delta + v\cos\delta.
\]

%------------------------------------------------------------------------------
\subsection{Rotated Latitude--Longitude (\texttt{ROTLLR})}
\label{subsec:proj_rotllr}

The rotated north pole is at (\code{plat}, \code{plon}). \code{get\_equator}
converts it into the rotation angle $\varphi_0$ (the colatitude of the
pole, with the sign of its hemisphere) and the longitude $\lambda_0$ of the
point where the rotated prime meridian crosses the rotated equator. The
private subroutines \code{geo\_to\_rotated} and \code{rotated\_to\_geo}
rotate the unit vector of a point around the axis through the longitudes
$\lambda_0 \pm 90^\circ$:
\[
  \begin{aligned}
    \sin\varphi' &= -\sin\varphi_0\cos\varphi\cos(\lambda-\lambda_0)
                    + \cos\varphi_0\sin\varphi, \\
    \lambda' &= \text{atan2}\bigl(\cos\varphi\sin(\lambda-\lambda_0),\;
                \cos\varphi_0\cos\varphi\cos(\lambda-\lambda_0)
                + \sin\varphi_0\sin\varphi\bigr),
  \end{aligned}
\]
and the inverse with $-\varphi_0$. When $|\varphi_0| \le 10^{-7}$ (pole at
the geographic pole) the coordinates are not rotated.

The grid is regular in the rotated coordinates, with
$\Delta\varphi' = \Delta\lambda' = \Delta s/R$ in degrees:
\[
  \varphi' = \varphi'_1 + (j-1)\Delta\varphi', \qquad
  \lambda' = \lambda'_1 + (i-1)\Delta\lambda',
\]
where $(\varphi'_1, \lambda'_1)$ (\code{rlat0}, \code{rlon0}, returned by
\code{rl00}) are the rotated coordinates of the first point, computed from
those of the domain centre. The domain may cross the rotated longitude
$180^\circ$: \code{llij\_rl} takes $\lambda'$ within $180^\circ$ of the
rotated longitude of the centre before computing $i$. The map factor is
$m = 1/\cos\varphi'$, the factor of the $i$ direction (the grid is not
conformal: the $j$ direction has a factor of 1).

\textbf{Wind rotation.} The angle $\delta$ between the geographic and the
rotated meridians at each point is
\[
  \delta = \text{atan2}\bigl(\cos\varphi_p\sin(\lambda_p - \lambda),\;
           \cos\varphi\sin\varphi_p -
           \sin\varphi\cos\varphi_p\cos(\lambda_p - \lambda)\bigr)
\]
with $(\varphi_p, \lambda_p)$ the rotated pole, and the wind is rotated as
for \code{ROTMER}. If the rotated pole is at the geographic north pole
(\code{plat} within $0.001^\circ$ of $90^\circ$) the rotation is skipped
(\code{skiprot}).

%------------------------------------------------------------------------------
\subsection{Fallback Latitude--Longitude Grid}
\label{subsec:proj_ll}

For an unknown code, \code{setup\_ll} defines a regular grid in geographic
coordinates with $\Delta\varphi = \Delta\lambda = \Delta s/R$ in degrees,
centred as \code{ROTLLR}. The map factor is set to 1 and the winds are not
rotated. Across a pole \code{ijll\_ll} reflects the latitude
($180^\circ - \varphi$) and moves to the opposite meridian; the same
reflection is applied in the rotated coordinates by \code{ijll\_rl} and
\code{llij\_rl}.

%------------------------------------------------------------------------------
\subsection{Design Notes}
\label{subsec:proj_design}

\textbf{Procedure pointers.} The projection is chosen once, in
\code{initialize}; afterwards the calls go directly to the routines of the
projection through the procedure pointers. The type-bound
\code{wind\_rotate} family selects the routine with a \code{select case}
on the code instead, and does nothing for the projections without
rotation.

\textbf{Precomputed rotation factors.} The rotation angle depends only on
the grid point, so $\cos$ and $\sin$ are computed once in
\code{initialize} and the rotation of a 3-D field costs four
multiplications per point. This is why each staggering needs its own
projection object, and why the wind arrays must have the grid shape.

\textbf{Double precision.} All the computations are done in \code{rk8},
whatever \code{rkx}, and only the results are converted. The rotated
latitude--longitude conversions use \code{atan2} instead of \code{acos},
which loses precision near $\cos = \pm1$ and needs the sign of the
longitude from a separate test.

\textbf{Map factor of non-conformal grids.} RegCM uses one map factor per
point. For \code{ROTLLR} it is the factor of the $i$ direction,
$1/\cos\varphi'$; the $j$ direction has a factor of 1. The fallback grid
uses 1 for both.

\textbf{Numerical checks.} The projections were checked on grids of
$61\times51$ points (and $61\times201$ for the Mercator grids) for:
centre position; round trip \code{ijll} $\to$ \code{llij}; map factor
against $\Delta s$ divided by the great-circle distance of neighbouring
points; and wind rotation against the direction of the grid axes computed
by finite differences, plus rotation followed by back rotation. The cases
were \code{LAMCON} (north, south, tangent), \code{POLSTR} (north, south),
\code{NORMER} (equator, $30^\circ$N), \code{ROTMER} (\code{plat} equal to
and different from \code{clat}, south) and \code{ROTLLR} (EUR and AUS
CORDEX poles, South America, a domain across the dateline). All agree to
round-off, or to the accuracy of the finite differences.

%==============================================================================
\section{Heights and Pressures --- \texttt{mod\_hgt.F90}}
\label{sec:hgt}
%==============================================================================

\subsection{Overview}

\code{mod\_hgt} collects the hydrostatic computations used to move between
heights and pressures in the pre- and post-processing:

\begin{itemize}
  \item heights of the model levels, from temperature and pressure, for the
        hydrostatic (\code{hydrost}, \code{htsig\_2}, \code{htsig\_o}) and
        non-hydrostatic cores (\code{nonhydrost}), and for input data on
        pressure or hybrid levels (\code{htsig\_1}, \code{htsig\_3});
  \item heights of pressure levels, with extrapolation above the top and
        below the surface (\code{height}, \code{height\_o});
  \item pressures of height levels (\code{psig}, \code{psig1});
  \item mean sea level pressure with the ECMWF method (\code{mslp}) and a
        smoothing filter for it (\code{gs\_filter});
  \item conversion between pressure and the Exner function (\code{p2pai},
        \code{pai2p}).
\end{itemize}

The routines are used by the ICBC programs (\code{mod\_cmip6},
\code{mod\_cmip7}, \code{mod\_gnhnc} and \code{mod\_nest}) and by the
\code{sigma2p} and \code{sigma2z} post-processors. The model has its own
copy of the sea level pressure computation, in \code{mod\_lm\_interface}
(see Section~\ref{subsec:hgt_mslp}).

\subsection{Dependencies}

\begin{lstlisting}
mod_hgt
    |-- mod_dynparam    (maxvert)
    |-- mod_intkinds    (ik4)
    |-- mod_realkinds   (rk4, rk8, rkx)
    |-- mod_constants   (rovg, rgas, egrav, lrate, ep1, rovcp, cpovr, p00,
    |                    d_zero, d_half, d_one, d_two)
    |-- mod_message     (not used)
\end{lstlisting}

\code{maxvert} (256) sizes the local work arrays of one column, so the
number of levels must not exceed it. The private constants \code{srovg},
\code{slrate}, \code{segrav} and \code{srgas} are the single precision
values of $R/g$, $\Gamma$, $g$ and $R$ for the \code{rk4} variants.

\subsection{Conventions}

\begin{itemize}
  \item Heights are in m above sea level unless stated otherwise, pressures
        in any consistent unit (the callers use Pa, hPa or kPa), and
        temperatures in K.
  \item The arrays are \code{(i, j, k)}. Most routines assume the vertical
        index increasing downward, from the top; \code{height},
        \code{htsig\_1} and \code{htsig\_3} detect the order from the
        pressures of the first column, and \code{psig}/\code{psig1} need it
        increasing upward, from the bottom.
  \item The thickness of a layer is computed with the hypsometric equation,
        \[
          z_2 - z_1 = \frac{R}{g}\,\bar T\,\ln\frac{p_1}{p_2},
        \]
        with $\bar T$ the mean of the temperatures of the two levels
        (weighted with their sigma thickness in \code{hydrost}).
  \item A surface pressure below $-9995$ marks a missing value in
        \code{height} and \code{htsig\_1}; the result is then $-9999$.
  \item Generic interfaces select the precision: the \code{\_double},
        \code{\_single} and \code{\_nonhy} variants differ only in the kind
        of the arrays and in the way the level pressures are given.
\end{itemize}

%------------------------------------------------------------------------------
\subsection{Heights of the Model Levels}
\label{subsec:hgt_levels}

\begin{longtable}{L{5.5cm} L{8.5cm}}
  \toprule
  \textbf{Routine} & \textbf{Description} \\
  \midrule
  \endhead
  \bottomrule
  \endfoot
  \code{hydrost(h, t, topo, ps, ptop, sigmah, ni, nj, nk)} &
    Heights of the half sigma levels of the hydrostatic core, consistent
    with the model. \code{ps} is $p^* = p_s - p_\text{top}$: the pressure
    of level $k$ is $\sigma_k p^* + p_\text{top}$. The mean temperature of a
    layer is weighted with the sigma thickness of the two levels. \\
  \code{nonhydrost(h, t0, p0, ps, topo, ...)} &
    Heights of the levels of the non-hydrostatic core from the reference
    temperature and pressure \code{t0}, \code{p0}. Three variants: double
    and single precision with bounds $1{:}n$, and \code{nonhydrost\_single\_2},
    which is double precision (despite the name) with bounds
    \code{i1:i2}, \code{j1:j2}. \\
  \code{htsig\_2}, \code{htsig\_o(t, h, pstar, ht, sig, ptop, ...)} &
    Heights of sigma levels; here \code{pstar} is the surface pressure and
    the pressure of level $k$ is $\sigma_k(p_s - p_\text{top}) +
    p_\text{top}$. \code{htsig\_o} has double and single precision
    variants. \\
  \code{htsig\_1(t, h, p3d, ps, ht, im, jm, km)} &
    Heights of levels with given 3-D pressure, in either vertical order,
    with missing values. \\
  \code{htsig\_3(z, t, p, q, ps, ht)} &
    As \code{htsig\_1} with the virtual temperature
    $T_v = T(1 + \epsilon_1 q)$, $\epsilon_1 = M_d/M_w - 1$, on pointer
    arrays whose sizes give the dimensions. \\
\end{longtable}

All integrate upward from the surface: the first layer goes from the
surface height to the lowest level with the temperature of the lowest
level.

%------------------------------------------------------------------------------
\subsection{Heights of Pressure Levels}
\label{subsec:hgt_plevels}

\code{height(hp, h, t, ps, p3d, ht, im, jm, km, p, kp)} computes the height
\code{hp} of the \code{kp} pressure levels \code{p} from height, temperature
and pressure on the model levels; \code{height\_o} does the same for
sigma levels (\code{sig}, \code{ptop}) or for given 3-D pressures
(\code{\_nonhy}). For each pressure $p$, with $p_t$ and $p_b$ the pressures
of the top and lowest levels and $p_s$ the surface pressure:

\begin{longtable}{L{3.5cm} L{10.5cm}}
  \toprule
  \textbf{Case} & \textbf{Height} \\
  \midrule
  \endhead
  \bottomrule
  \endfoot
  $p < p_t$ (above the top) &
    Isothermal with the top temperature:
    $z = z_t + (R/g)\,T_t\ln(p_t/p)$. \\
  between two levels &
    From the lower level $b$, with the mean of $T_b$ and the temperature
    interpolated linearly in $\ln p$ at $p$:
    $z = z_b + (R/g)\,\bar T\ln(p_b/p)$. \\
  $p_b \le p \le p_s$ &
    From the surface, with the lowest level temperature:
    $z = z_s + (R/g)\,T_b\ln(p_s/p)$. \\
  $p > p_s$ (below the surface) &
    Constant lapse rate $\Gamma$ (\code{lrate}) from the surface, starting
    from the mean temperature $T_0$ of the two lowest levels (R. Errico):
    $z = z_s + (T_0/\Gamma)\bigl(1 - (p/p_s)^{R\Gamma/g}\bigr)$. \\
\end{longtable}

\code{height} has OpenACC and \code{do concurrent} (\code{STDPAR})
versions of the column loop.

%------------------------------------------------------------------------------
\subsection{Pressures of Height Levels}
\label{subsec:hgt_psig}

\code{psig(t, h, p3d, ps, ht, im, jm, km)} computes the pressure of levels
of given height, ordered from the bottom, by integrating the hypsometric
equation upward from the surface:
\[
  p_1 = p_s \exp\!\left(-\frac{z_1 - z_s}{(R/g)\,T_1}\right), \qquad
  p_k = p_{k-1}\exp\!\left(-\frac{z_k - z_{k-1}}{(R/g)\,\bar T_k}\right)
\]
with $\bar T_k$ the mean of $T_{k-1}$ and $T_k$. \code{psig1} uses the
temperature of the upper level, $T_k$, instead of the mean. They are used
for input data on height levels (HadGEM3 in \code{mod\_cmip6},
\code{mod\_gnhnc}).

%------------------------------------------------------------------------------
\subsection{Mean Sea Level Pressure: \texttt{mslp}}
\label{subsec:hgt_mslp}

\begin{lstlisting}
subroutine mslp(t, pl, ps, ht, slp, im, jm, kz)   ! single precision
\end{lstlisting}

\code{t} is the temperature on the model levels (\code{kz} is the lowest),
\code{pl} the pressure of the lowest level, \code{ps} the surface pressure
and \code{ht} the surface height. The method is the ECMWF one, as described
in section~3.1.b of Trenberth, Berry and Buja (1993, NCAR/TN-396+STR), and
also used in the NCAR CAM model:

\begin{enumerate}
  \item Surface temperature, extrapolated from the lowest level with the
        standard lapse rate, $\alpha = \Gamma R/g$:
        $T^* = T_{k_z}\bigl(1 + \alpha(p_s/p_l - 1)\bigr)$.
  \item Sea level temperature $T_0 = T^* + \Gamma z_s$, and limits on the
        extrapolated column:
        \begin{itemize}
          \item $T^* \le 290.5 < T_0$: the lapse rate is reduced to reach
                $290.5$~K at sea level, $\alpha' = R(290.5 - T^*)/(g z_s)$;
          \item $T^* > 290.5$ and $T_0 > 290.5$: $\alpha' = 0$ and
                $T^* = (290.5 + T^*)/2$;
          \item otherwise $\alpha' = \alpha$, and if $T^* < 255$~K,
                $T^* = (255 + T^*)/2$.
        \end{itemize}
  \item With $\beta = g z_s/(R T^*)$ and $x = \alpha'\beta$,
        \[
          p_\text{msl} = p_s\exp\!\left[\beta\left(1 - \frac{x}{2} +
                         \frac{x^2}{3}\right)\right],
        \]
        the expansion of $p_s(1 + \Gamma' z_s/T^*)^{g/(R\Gamma')}$.
\end{enumerate}

\noindent Where $|z_s| < 10^{-4}$~m the sea level pressure is the surface
pressure. The limits avoid unrealistic values over high and hot or cold
terrain, where the extrapolated column does not exist; in a standard
atmosphere the result is exact to the truncation of the series (0.3~hPa for
a surface at 5000~m).

The model computes the sea level pressure for its output in
\code{mslp} of \code{mod\_lm\_interface}, with the same method on the model
fields (surface geopotential, lowest level temperature and pressure),
followed by a smoothing like \code{gs\_filter} but with 20 Jacobi
iterations, across the MPI subdomains.

%------------------------------------------------------------------------------
\subsection{Smoothing: \texttt{gs\_filter}}
\label{subsec:hgt_gsfilter}

\code{gs\_filter(v, vm, im, jm)} smooths the 2-D field \code{v} (the sea
level pressure) with 20 Gauss--Seidel sweeps of
\[
  v_{i,j} = \tfrac14\bigl(v_{i-1,j} + v_{i+1,j} + v_{i,j-1} + v_{i,j+1}
            - m_{i,j}\bigr),
\]
keeping the boundary values. The source term $m$ is the Laplacian of
\code{vm} (the surface pressure) divided by half its range, so that the
smoothed field keeps the structure of the surface pressure. If \code{vm}
is constant the source term is zero.

%------------------------------------------------------------------------------
\subsection{Exner Function}
\label{subsec:hgt_exner}

The elemental functions \code{p2pai(p)} $= (p/p_{00})^{R/c_p}$ and
\code{pai2p(pai)} $= p_{00}\,\pi^{c_p/R}$, with $p_{00} = 10^5$~Pa, convert
between pressure [Pa] and the Exner function used by the MOLOCH core.

%------------------------------------------------------------------------------
\subsection{Design Notes}
\label{subsec:hgt_design}

\textbf{Extrapolation below the surface.} Two kinds of extrapolation are
used and must not be confused. Between the lowest level and the surface
the air exists, and the actual temperature of the lowest level is used.
Below the surface (pressure levels with $p > p_s$, sea level pressure) the
column is fictitious: \code{height} uses a constant lapse rate, and
\code{mslp} the ECMWF limits on the column temperature. The same distinction
applies to the surface pressure extrapolation in \code{mod\_vertint}
(\code{intzps}, below the lowest level of the input data, with the ECMWF
limits) and in \code{mod\_write}/\code{mod\_moloch} (between the lowest
model level and the surface, without limits).

\textbf{Numerical checks.} The routines were checked against an exact
hydrostatic atmosphere with a constant lapse rate, built with the
constants of the code. The level heights agree to the accuracy of the
mean layer temperature (centimetres for layers of a few hundred metres),
the heights of pressure levels between levels to round-off, and
\code{psig} to $10^{-6}$ near the surface.

%==============================================================================
\section{Sigma Levels --- \texttt{mod\_sigma.F90}}
\label{sec:sigma}
%==============================================================================

\subsection{Overview}

\code{mod\_sigma} builds the vertical sigma coordinate of the hydrostatic
and non-hydrostatic cores. It is used by the terrain program, which writes
the levels in the domain file, and by the \code{sigma2p}/\code{sigma2z}
post-processors. The levels are stored in three public arrays:

\begin{longtable}{L{5cm} L{9cm}}
  \toprule
  \textbf{Array} & \textbf{Content} \\
  \midrule
  \endhead
  \bottomrule
  \endfoot
  \code{sigma\_coordinate(1:nk+1)} &
    Full (interface) levels, from 0 at the top to 1 at the surface. \\
  \code{half\_sigma\_coordinate(1:nk)} &
    Half levels, the mid points of the full levels, where the model
    variables are defined. \\
  \code{sigma\_delta(1:nk)} &
    Layer thicknesses $\Delta\sigma_k = \sigma_{k+1} - \sigma_k$. \\
\end{longtable}

In the hydrostatic core the pressure is $p = \sigma p^* + p_\text{top}$
with $p^* = p_s - p_\text{top}$.

\subsection{Dependencies}

\begin{lstlisting}
mod_sigma
    |-- mod_intkinds, mod_realkinds
    |-- mod_constants   (d_zero, d_half, d_one, d_two)
    |-- mod_memutil     (getmem)
    |-- mod_message     (die), mod_stdio (stdout, stderr)
\end{lstlisting}

\subsection{\texttt{init\_sigma}}

\begin{lstlisting}
subroutine init_sigma(nk, dmax, dmin, lzeta)
\end{lstlisting}

\begin{itemize}
  \item \textbf{Predefined sets.} For $n_k$ = 2 (CLM stand-alone), 14
        (RegCM2), 18 (RegCM3), 23 (MM5~V3) and 41 (high resolution), the
        levels are hard-coded, unless \code{lzeta} is true.
  \item \textbf{Computed sets.} Otherwise the layer thicknesses follow
        \[
          \Delta\sigma_k = \Delta\sigma_{\max}\,a^{k-1}\,b^{(k-1)(k-2)/2},
        \]
        so that they decrease roughly exponentially from
        $\Delta\sigma_{\max}$ at the top to $\Delta\sigma_{\min}$ near the
        surface. Given $b$, $a$ is fixed by
        $\Delta\sigma_{n_k} = \Delta\sigma_{\min}$, and $b$ is found by a
        search with adaptive step until $\sum\Delta\sigma_k = 1$ within
        $10^{-5}$ (at most $10^6$ steps; the run stops if it does not
        converge). The defaults are $\Delta\sigma_{\max} = 0.05$ and
        $\Delta\sigma_{\min} = 0.01$ (0.01 and 0.005 with \code{lzeta}),
        overridden by \code{dmax} and \code{dmin}. They must satisfy
        $n_k\Delta\sigma_{\max} \ge 1$ and $n_k\Delta\sigma_{\min} < 1$.
\end{itemize}

In both cases the last full level is set to exactly 1, and the half levels
and thicknesses are computed from the full levels.

%==============================================================================
\section{Humidity Conversions --- \texttt{mod\_humid.F90}}
\label{sec:humid}
%==============================================================================

\subsection{Overview}

\code{mod\_humid} converts between specific humidity, mixing ratio and
relative humidity, in place. It includes \code{pfwsat.inc}
(Section~\ref{sec:includes}), the saturation mixing ratio over water
($T \ge 0^\circ$C) or ice, $q_s = \epsilon\,e_s/(p - e_s)$ with $e_s$ from
an eighth-order polynomial and limited to $0.15\,p$. Relative humidity is
defined with the mixing ratio, $r = q/q_s$.

\subsection{Dependencies}

\begin{lstlisting}
mod_humid
    |-- mod_intkinds, mod_realkinds
    |-- mod_constants   (tzero, ep2, govr, d_zero, d_half, d_one, d_1000)
    |-- pfwsat.inc
\end{lstlisting}

\subsection{Public Procedures}

\begin{longtable}{L{5.5cm} L{8.5cm}}
  \toprule
  \textbf{Procedure} & \textbf{Description} \\
  \midrule
  \endhead
  \bottomrule
  \endfoot
  \code{sph2mxr(q, ni, nj, nk)} &
    Specific humidity to mixing ratio, $q/(1-q)$; double and single
    precision. \\
  \code{mxr2sph(q, ni, nj, nk)} &
    Mixing ratio to specific humidity, $q/(1+q)$. \\
  \code{mxr2rh(t, q, ps, sigma, ptop, ...)} &
    Mixing ratio to relative humidity on sigma levels, with $p_s$ the
    surface pressure: $p = \sigma(p_s - p_\text{top}) + p_\text{top}$;
    double and single precision. \\
  \code{mxr2rh(t, q, p3d, ni, nj, nk)} &
    The same with the 3-D pressure (non-hydrostatic and MOLOCH output). \\
  \code{rh2mxr(t, q, p, nk)}, \code{rh2mxr(t, q, p, ni, nj, nk)} &
    Relative humidity to mixing ratio on pressure levels, $p$ in Pa (one
    column, or a field with 1-D pressure levels). \\
  \code{rh2mxr(t, q, ps, ptop, sigma, ni, nj, nk)} &
    The same on sigma levels, with \code{ps} $= p^*$ and \code{ptop} in
    kPa: $p = 1000(p_\text{top} + \sigma p^*)$ Pa. \\
  \code{rh2mxr\_z(t, q, z, ps, ni, nj, nk)} &
    The same on height levels (MOLOCH): \code{z} is the height above the
    surface, ordered from the top, and \code{ps} the surface pressure in
    kPa. The pressure of the levels is integrated hydrostatically from the
    surface with the mean layer temperature. \\
\end{longtable}

The conversions to mixing ratio and relative humidity keep the result
non-negative.

%==============================================================================
\section{MOLOCH Vertical Coordinate --- \texttt{mod\_zita.F90}}
\label{sec:zita}
%==============================================================================

\subsection{Overview}

\code{mod\_zita} defines the height-based, terrain-following coordinate
$\zeta$ of the MOLOCH core, from 0 at the surface to $\zeta_\text{top}$
(\code{mo\_ztop}), and the related functions. The height of a point is
\[
  z = b(\zeta) + h\,g(\zeta)
\]
with $h$ the orography, so that $z = h$ at $\zeta = 0$ and
$z = \zeta_\text{top}$ at the top, where the effect of the orography has
vanished.

\subsection{Dependencies}

\begin{lstlisting}
mod_zita
    |-- mod_realkinds, mod_intkinds
    |-- mod_constants   (egrav, rgas, regrav)
\end{lstlisting}

\subsection{Levels}

\code{model\_zitaf(zitaf, ztop)} and \code{model\_zitah(zitah, ztop)}
compute equally spaced full levels ($\zeta_\text{top}$ to 0, $n_z+1$
values) and half levels (their mid points), ordered from the top. The
functions \code{sigmazita} $= 1 - \zeta/\zeta_\text{top}$ and its inverse
\code{zitasigma} map $\zeta$ to a nominal sigma, used as the vertical axis
of the files and for data prepared on pseudo-sigma levels
($p = \sigma p_s$, converted to the real pressures of the $\zeta$ levels by
\code{zita\_interp}). It is not the pressure ratio of the levels.

\subsection{Functions of the Coordinate}

\begin{longtable}{L{5cm} L{9cm}}
  \toprule
  \textbf{Function} & \textbf{Definition} \\
  \midrule
  \endhead
  \bottomrule
  \endfoot
  \code{bzita(zita, ztop, zh)} &
    $b(\zeta) = Z_f\,(e^{\zeta/z_h} - 1)$ with
    $Z_f = \zeta_\text{top}/(e^{\zeta_\text{top}/z_h} - 1)$
    (\code{zfz}): $b(0) = 0$, $b(\zeta_\text{top}) = \zeta_\text{top}$,
    with the vertical scale $z_h$ (\code{mo\_h}). \code{md\_ak} is the same
    function. \\
  \code{gzita(zita, ztop, a0)} &
    Cubic decay of the orography, with $r = \zeta/\zeta_\text{top}$:
    $g = 1 - a_0 r - (3 - 2a_0)r^2 + (2 - a_0)r^3$, so that $g(0) = 1$
    and $g(\zeta_\text{top}) = 0$; $a_0$ (\code{mo\_a0}) sets the slope at
    the surface. \code{md\_bk} is the same function. \\
  \code{bzitap}, \code{gzitap} &
    Their derivatives with respect to $\zeta$ (private). \\
  \code{md\_zeta\_h(zita, orog, ...)} &
    Height above the orography, $z - h = b + h(g - 1)$. \\
  \code{md\_fmz\_h(zita, orog, ...)} &
    Metric factor $1/(\partial z/\partial\zeta) = 1/(b' + h g')$. \\
  \code{md\_zeta}, \code{md\_fmz} &
    The same with the surface geopotential instead of the orography. \\
\end{longtable}

\subsection{\texttt{zita\_interp}}

\begin{lstlisting}
call zita_interp(nx1, nx2, ny1, ny2, nz, [nn,] f, zeta, tvirt, sigmah, ps, imet)
\end{lstlisting}

Moves the fields \code{f} (3-D or 4-D, double or single precision), given
on the pseudo-sigma levels with pressure $p_\sigma = \sigma_k p_s$, to the
$\zeta$ levels, whose pressure is estimated as
$p_\zeta = p_s\exp\bigl(-g z/(R T_v)\bigr)$ with the height above ground
\code{zeta} and the virtual temperature of the level. The interpolation is
linear in $\ln p$ (\code{imet} = 1) or in $p$; outside the range the end
values are kept. It is used for the chemistry and aerosol boundary
conditions.

%==============================================================================
\section{One-Dimensional Interpolation --- \texttt{mod\_interp.F90}}
\label{sec:interp}
%==============================================================================

\subsection{Overview}

\code{mod\_interp} provides \code{interp1d}, the one-dimensional
interpolation with a tension spline of the MOLOCH pre-processing. It is
used by \code{mod\_vertint} (\code{intz1}, \code{intz3}, \code{intp1},
\code{intp3}) for the vertical interpolation of the initial and boundary
conditions, and by \code{mod\_rad\_aerosol} for the aerosol profiles.

\subsection{Dependencies}

\begin{lstlisting}
mod_interp
    |-- mod_intkinds    (ik4)
    |-- mod_realkinds   (rk4, rk8)
\end{lstlisting}

\subsection{\texttt{interp1d}}

\begin{lstlisting}
call interp1d(xi, g, xo, f, alfa, ex1, ex2)
call interp1d(xi, g, xo, f, alfa, ex1, ex2, zi, zg, sxi, sxo)
\end{lstlisting}

\begin{longtable}{L{2.2cm} L{1.5cm} L{10.3cm}}
  \toprule
  \textbf{Argument} & \textbf{Intent} & \textbf{Description} \\
  \midrule
  \endhead
  \bottomrule
  \endfoot
  \code{xi(:)}, \code{g(:)} & \code{in} &
    Coordinates and values of the input profile, increasing or
    decreasing. \\
  \code{xo}, \code{f} & \code{in}, \code{out} &
    Coordinates where the profile is wanted and the result: arrays, or
    scalars for one point. \\
  \code{alfa} & \code{in} &
    Tension: 1 gives linear interpolation, 0 the cubic, values in between
    a weighted mean of the two. \\
  \code{ex1}, \code{ex2} & \code{in} &
    Extrapolation below the smallest and above the largest coordinate:
    0 keeps the end value, 1 extrapolates linearly with the slope of the
    end interval, values in between a fraction of that slope. \\
  \code{zi}, \code{zg}, \code{sxi}, \code{sxo} & work, \code{in} &
    Only in the \code{\_new} variants: work arrays and explicit sizes, so
    that the routine needs no automatic arrays and can run on the GPU as an
    OpenACC \code{routine seq} (used by \code{intp1} with pointer
    arrays). \\
\end{longtable}

The generic interface has six variants: arrays and scalar \code{xo} in
single and double precision, and the two \code{\_new} variants. They
differ only in the kinds and in the way the arrays are passed.

\subsubsection{Algorithm}

\begin{enumerate}
  \item \textbf{Ordering.} A decreasing profile is reversed, so that the
        coordinates $z_k$ (\code{zi}) increase.
  \item \textbf{Coincident points.} The private \code{merge\_coincident}
        merges every coordinate not larger than the previous one, with its
        value, into one point with the mean coordinate and value, and checks
        the result again against the previous point, until the coordinates
        are strictly increasing. One point left gives that value everywhere;
        with two points the interpolation between them is linear (the
        extrapolation follows the next step).
  \item \textbf{Extrapolation.} For $x < z_1$,
        $f = g_1 + e_1\,s_1(x - z_1)$ with $s_1$ the slope of the first
        interval; for $x \ge z_n$ the same with $e_2$ and the last interval.
  \item \textbf{Interval.} The interval $z_r \le x < z_{r+1}$ is found,
        with $f_m = g_r$, $f_p = g_{r+1}$ and the neighbours
        $f_{mm} = g_{r-1}$, $f_{pp} = g_{r+2}$. In the first and last
        intervals the missing neighbour is a ghost point on the straight
        line through the two end points, e.g.\ $z_0 = 2z_1 - z_2$,
        $g_0 = 2g_1 - g_2$.
  \item \textbf{Interpolation.} With $\Delta = z_p - z_m$,
        $\Delta_m = z_m - z_{mm}$, $\Delta_p = z_{pp} - z_p$,
        $d_1 = x - z_m$ and $d_2 = z_p - x$, the cubic estimate is
        \[
          \begin{aligned}
            s ={}& f_m\left(\frac{d_2}{\Delta} + \frac{d_1 d_2^2}{\Delta^2\Delta_m}
                   - \frac{d_1^2 d_2}{(\Delta+\Delta_p)\Delta^2}\right)
                 + f_p\left(\frac{d_1}{\Delta} + \frac{d_1^2 d_2}{\Delta^2\Delta_p}
                   - \frac{d_1 d_2^2}{(\Delta+\Delta_m)\Delta^2}\right) \\
                 & - f_{mm}\frac{d_1 d_2^2}{(\Delta+\Delta_m)\Delta\Delta_m}
                   - f_{pp}\frac{d_1^2 d_2}{(\Delta+\Delta_p)\Delta\Delta_p},
          \end{aligned}
        \]
        the linear one is $l = (f_m d_2 + f_p d_1)/\Delta$, and the result
        is $f = \alpha\,l + (1 - \alpha)\,s$.
\end{enumerate}

\subsubsection{Properties}

\begin{itemize}
  \item A linear profile is reproduced exactly for any $\alpha$, $e_1$,
        $e_2$ (round-off), also with coincident points.
  \item With $\alpha = 0$ a quadratic profile is reproduced exactly in the
        inner intervals; in the first and last intervals the ghost points,
        which are on a straight line, make the result less accurate.
  \item A decreasing profile gives the same result as the same profile
        increasing, and the scalar and array forms agree.
\end{itemize}

\subsection{Design Notes}

\textbf{Coincident coordinates.} The routine accepts profiles with
repeated or slightly unordered coordinates. Earlier versions moved coincident points apart by
$10^{-6}$ of the range; the cubic then divides by these tiny spacings in
the neighbouring intervals, and in a test a linear profile with one
repeated level was off by about 0.3 of its values there. Merging the points
keeps the interpolation exact, and does not change the result for profiles
without coincident points.

\textbf{Removed code.} The module used to contain horizontal
interpolations (bilinear, distance-weighted, Cressman, kernel smoothing)
that are no longer used; they were removed together with their module
state.

%==============================================================================
\section{Cubic Splines --- \texttt{mod\_spline.F90}}
\label{sec:spline}
%==============================================================================

\subsection{Overview}

\code{mod\_spline} provides two implementations of one-dimensional cubic
spline interpolation:
\begin{itemize}
  \item the routines of \citet{press1992}: \code{spline} computes the
        second derivatives of the spline at the input points once, then
        \code{splint} interpolates and \code{splini} integrates at any
        number of points. They are used by the MSIS atmosphere of the
        radiation (\code{mod\_atm\_msis});
  \item the routine of Z.~Janjic (Yugoslav Federal Hydrometeorological
        Institute, Belgrade), from the Eta model, which solves for the
        second derivatives and interpolates in one call. It exists in two
        versions: \code{spline} with ten arguments, used by
        \code{mod\_capecin} for the CAPE and CIN diagnostics, and its Fortran
        90 recoding \code{spline1d}, used for the vertical profile of the
        boundary nudging (\code{exponential\_nudging} in
        \code{mod\_runparams}).
\end{itemize}

\subsection{Dependencies}

\begin{lstlisting}
mod_spline
    |-- mod_intkinds, mod_realkinds
    |-- mod_constants (d_zero, d_half, d_one, d_two, d_three, d_six, ...)
\end{lstlisting}

\subsection{Interfaces}

\begin{longtable}{L{5.2cm} L{8.8cm}}
  \toprule
  \textbf{Interface} & \textbf{Description} \\
  \midrule
  \endhead
  \bottomrule
  \endfoot
  \code{spline(x, y, yp1, ypn, y2)} &
    Second derivatives \code{y2} of the spline through the points
    $(x_i, y_i)$, $x$ increasing. \code{yp1} and \code{ypn} are the first
    derivatives at the two ends; a value above $0.99\cdot10^{30}$ selects a
    natural spline (zero second derivative) at that end. \\
  \code{splint(xa, ya, y2a, x, y)} &
    Value \code{y} of the spline at \code{x}, with the interval found by
    bisection. Outside the points the cubic of the first or last interval is
    extrapolated. Coincident points around \code{x} return $10^{-30}$. \\
  \code{splini(xa, ya, y2a, x, yi)} &
    Integral \code{yi} of the spline from \code{xa(1)} to \code{x}, the sum of
    the exact integrals of the cubics of the intervals; beyond the last
    point the last cubic is integrated, before the first point the result
    is 0. \\
  \code{spline(jtb, nold, xold, yold, y2, nnew, xnew, ynew, p, q)} &
    Janjic spline: values \code{ynew} at the \code{nnew} points \code{xnew}
    of the spline through the \code{nold} $\geq 3$ points
    $(x^\text{old}, y^\text{old})$. On input \code{y2(1)} and
    \code{y2(nold)} are the second derivatives at the ends (0 for a natural
    spline); on output \code{y2} holds all of them. \code{p} and \code{q}
    are work arrays; all the arrays have dimension \code{jtb}. Points above
    the last one get \code{yold(nold)}; points below the first one get the
    extrapolated cubic of the first interval. \\
  \code{spline1d(nold, xold, yold, y2, nnew, xnew, ynew)} &
    The same algorithm with arrays of their own size. Points outside the
    input range get the value of the nearest end point. \\
\end{longtable}

\subsection{Algorithms}

Both implementations solve the tridiagonal system for the second
derivatives $M_k$ of a cubic spline,
\begin{equation}
  \Delta x_{k-1} M_{k-1} + 2(\Delta x_{k-1}+\Delta x_k) M_k +
  \Delta x_k M_{k+1} = 6\left(\frac{y_{k+1}-y_k}{\Delta x_k} -
  \frac{y_k-y_{k-1}}{\Delta x_{k-1}}\right),
\end{equation}
by forward elimination and back substitution. The Numerical Recipes
version adds the end conditions on the first derivative to the system; the
Janjic version takes the end second derivatives as given. In the interval
$[x_k, x_{k+1}]$ the Janjic routines evaluate the cubic as
$a t^3 + b t^2 + c t + y_k$ with $t = x - x_k$,
$a = (M_{k+1}-M_k)/(6\Delta x_k)$, $b = M_k/2$ and
$c = (y_{k+1}-y_k)/\Delta x_k - \Delta x_k(M_{k+1}+2M_k)/6$, keeping the
coefficients while successive points fall in the same interval.

\subsection{Design Notes}

\textbf{Corrections.} The following were corrected in October 2026, and
are checked by \code{Testing/share/test\_spline.F90} (a spline with the
exact end conditions reproduces a cubic polynomial):
\begin{itemize}
  \item \code{spline1d} multiplied the right-hand side of the system by
        $1/6$ instead of 6, so the second derivatives were 36 times too
        small and the spline nearly piecewise linear. With three points, as
        in \code{exponential\_nudging}, the nudging profile changes by up to
        2\,\%. For more than three points the elimination loop also read one
        element beyond the end of the input arrays;
  \item the Janjic \code{spline} used $0.1666667$ for $1/6$, an error of
        about $4\cdot10^{-6}$ on a cubic; it now uses the exact value;
  \item in the Janjic \code{spline}, a point above the last input point
        followed by one in the first interval reused coefficients that had
        never been computed;
  \item \code{splie2} and \code{splin2} (two-dimensional splines) were
        removed: they were not used, and \code{splin2} wrote to an array it
        never allocated.
\end{itemize}

%==============================================================================
\section{Vertical Interpolation --- \texttt{mod\_vertint.F90}}
\label{sec:vertint}
%==============================================================================

\subsection{Overview}

\code{mod\_vertint} interpolates fields between vertical grids: from the
input data (pressure levels, or model levels with given height or
pressure) to the model levels in the ICBC programs, from model levels to
pressure or height levels in the post-processors, and the ozone
climatology in the model. It also computes the surface pressure on the
model orography from the input profiles. The input are ``GCM'' (\code{ccm})
fields, the output ``RCM'' (\code{rcm}) fields.

\subsection{Dependencies}

\begin{lstlisting}
mod_vertint
    |-- mod_intkinds, mod_realkinds, mod_stdio, mod_message
    |-- mod_constants   (rgas, regrav, lrate, govr, govcp, egrav, ...)
    |-- mod_stdatm
    |-- mod_interp      (interp1d)
    |-- mod_dynparam    (maxvert)
\end{lstlisting}

Missing values are $-9999$ (\code{missl}). The private constants are
\code{rgas2} $= R/2$, \code{rglrog} $= R\Gamma/g$ and \code{b1} $= -g/\Gamma$.

\subsection{Interpolation Between Pressure Levels}

\begin{longtable}{L{4.5cm} L{9.5cm}}
  \toprule
  \textbf{Interface} & \textbf{Description} \\
  \midrule
  \endhead
  \bottomrule
  \endfoot
  \code{intlin(fp, f, p3d, ..., p, kp)} &
    From model levels with 3-D pressure \code{p3d} to the pressure levels
    \code{p}, linear in $p$; constant extrapolation. Double and single
    precision; the \code{\_o} variants take sigma levels (\code{ps},
    \code{sig}, \code{ptop}) instead of \code{p3d}. \\
  \code{intlog(fp, f, ...)} &
    As \code{intlin}, linear in $\ln p$. Below the lowest level the field
    (temperature) is extrapolated with the standard lapse rate from the
    mean of the two lowest levels,
    $f = \bar f_1\,(p/p_1)^{R\Gamma/g}$; above the top it is constant. \\
  \code{intlinz(fz, f, hz, ..., z, kz)} &
    From model levels with 3-D height \code{hz} to the height levels
    \code{z}, linear in $z$, in either vertical order. \\
  \code{intlinreg(fp, f, ps, p, ..., kp, p3d, km)} &
    From regular pressure levels \code{p} to the 3-D pressures \code{p3d},
    linear in $p$ (through the normalised coordinate
    $(p - p_\text{top})/(p_s - p_\text{top})$), constant extrapolation;
    double, single and mixed precision. Columns with $p_s < 1$ (missing)
    are left missing. \code{intlinprof} does the same for
    a single profile \code{f(kp)}. \\
\end{longtable}

\noindent All the logarithmic interpolations and extrapolations
(\code{intlog}, \code{intv2}, \code{intv3}) use the logarithm of the
pressure itself. Earlier versions took it of the normalised coordinate
$(p - p_\text{top})/(p_\text{bottom} - p_\text{top})$, which is equivalent
only for a negligible top pressure: with a top at 100~hPa the error reached
5~K at 150~hPa, and \code{intlog\_o} gave NaN near the top.

\subsection{Interpolation to the Model Levels (ICBC)}

The input data are on pressure levels described as a sigma coordinate:
$p = \sigma_k p_{ss} + p_{st}$, with \code{sccm} the level values (1 at the
lowest level, ordered from the bottom), \code{pss} and \code{pst}. The model
levels are sigma levels \code{srcm} with $p^*$ \code{psrcm} and
$p_\text{top}$ \code{pt}.

\begin{longtable}{L{4.5cm} L{9.5cm}}
  \toprule
  \textbf{Routine} & \textbf{Description} \\
  \midrule
  \endhead
  \bottomrule
  \endfoot
  \code{intv1} &
    Linear in $p$, constant extrapolation (winds, humidity); method 2
    keeps at least $10^{-8}$ (mixing ratio). \\
  \code{intv2} &
    Linear in $\ln p$; below the lowest level the temperature follows the
    standard lapse rate, $T = T_1(b_1 - a_1)/(b_1 + a_1)$ with
    $a_1 = (R/2)\ln(\sigma/\sigma_1)$, a Pad\'e form of
    $T_1(p/p_1)^{R\Gamma/g}$. \\
  \code{intv3} &
    The same for the surface level only (surface temperature). \\
  \code{intz1} &
    MOLOCH: from levels with given height to the model heights (height
    above ground plus orography), with the tension spline of
    \code{interp1d} (\code{a}: 1 linear, 0 spline; \code{e1}, \code{e2}:
    extrapolation below and above, 0 constant, 1 linear). Three variants:
    3-D input heights, 1-D input heights with geopotential orography, one
    column. \\
  \code{intz3}, \code{intp3} &
    The same for the surface value only, in height or in pressure. \\
  \code{intp1} &
    The same as \code{intz1} with pressure as vertical coordinate,
    normalised between the top and bottom input levels. \\
\end{longtable}

\subsection{Surface Pressure on the Model Orography}
\label{subsec:vertint_ps}

The input surface does not match the model orography, so the surface
pressure is computed from the input profiles of height \code{zp} and
temperature \code{tp}, ordered from the bottom, at the model surface height
\code{zrcm}:

\begin{longtable}{L{4.5cm} L{9.5cm}}
  \toprule
  \textbf{Routine} & \textbf{Description} \\
  \midrule
  \endhead
  \bottomrule
  \endfoot
  \code{intgtb} + \code{intpsn} &
    Hydrostatic and non-hydrostatic cores. \code{intgtb} finds the level
    \code{pa}, \code{za} above the model surface and the mean temperature
    of the layer between them (the temperature at the model surface is
    interpolated linearly in height); below the lowest level it uses the
    lowest level and its temperature. \code{intpsn} then gives
    $p^* = p_a\exp\bigl(-g(z_\text{rcm} - z_a)/(R\bar T)\bigr) -
    p_\text{top}$. \\
  \code{intzps} &
    MOLOCH: the surface pressure in one step, for input on pressure levels
    (\code{pss}, \code{sccm}, \code{pst}) or with 3-D pressure. Between two
    levels the layer from the model surface to the level above has the
    mean of the two temperatures. Below the lowest level the temperature
    follows the ECMWF method: the standard lapse rate with the limits on
    the extrapolated column (\code{below\_ground\_tlayer}). Above the top input level the
    result is missing for input on pressure levels, and the top pressure
    for input with 3-D pressure. \\
\end{longtable}

\noindent \code{below\_ground\_tlayer(tb, zb, dz)} returns the mean
temperature of the layer of depth \code{dz} below a level at height
\code{zb} with temperature \code{tb}. With $T^* = $ \code{tb}, the standard
lapse rate $\Gamma$ (\code{lrate}) and the sea level temperature
$T_0 = T^* + \Gamma z_b$, the rules of Section~\ref{subsec:hgt_mslp} are
applied: if $T^* \le 290.5 < T_0$ the lapse rate is reduced to have
$T_0 = 290.5$~K; if both exceed 290.5~K the lapse rate is zero and
$T^* = (290.5 + T^*)/2$; if $T^* < 255$~K, $T^* = (255 + T^*)/2$.

\subsection{Search Functions}

\code{find\_from\_top} and \code{find\_from\_btm} (private, double and
single precision) return the first level $n \ge 2$ at which the coordinate
is below (above) the requested value, or the last level.

\subsection{Design Notes}

\textbf{Two kinds of extrapolation.} As for \code{mod\_hgt}
(Section~\ref{subsec:hgt_design}), the extrapolation below the lowest
input level of \code{intzps} is the only surface pressure extrapolation
that uses the ECMWF limits: the extrapolated column there may be under the
input orography. The
surface pressure of the MOLOCH initial and boundary conditions is
recomputed from the lowest model level in \code{mod\_write}, over the real
layer between that level and the surface, without limits.

\textbf{Numerical checks.} \code{intzps} was checked against an exact
atmosphere on ERA5-like pressure levels: the surface pressure is exact for
surfaces between levels, and differs only where the ECMWF limits apply
(0.4~hPa for a surface below the 1000~hPa level in a 303~K atmosphere).
\code{intgtb} reproduces the mean layer temperature of a linear profile.

%==============================================================================
\section{Nearest-Neighbour Search --- \texttt{mod\_kdtree2\_priority.F90} and
         \texttt{mod\_kdtree2.F90}}
\label{sec:kdtree2}
%==============================================================================

\subsection{Overview}

These two modules are the KDTREE2 library of \citet{kennel2004}, distributed
under the Academic Free License 1.1, adapted to the RegCM kinds and error
handling. They build a $k$-d tree on a set of points in a Euclidean space and
find, for a query point, its $n$ nearest neighbours or all the points within
a given distance. RegCM uses them, through \code{mod\_kdinterp}, for the
horizontal interpolation of the pre-processing and of the climatological
input of the model, with the points as unit vectors on the sphere
(\code{ll2xyz} of \code{mod\_earth}): there the Euclidean distance is the
chord, a monotonic function of the great-circle distance.

\code{mod\_kdtree2\_priority} provides the bounded max-heap used to keep the
$n$ closest points during a search; \code{mod\_kdtree2} the tree.

\subsection{Dependencies}

\begin{lstlisting}
mod_kdtree2_priority
    |-- mod_realkinds   (rk4, rk8)
    |-- mod_message     (die)
mod_kdtree2
    |-- mod_realkinds, mod_message, mod_stdio (stderr)
    |-- mod_kdtree2_priority
\end{lstlisting}

The coordinates and distances are \code{real(kdkind)}, with \code{kdkind}
$=$ \code{rk8}.

\subsection{Types}

\begin{longtable}{L{3.5cm} L{10.5cm}}
  \toprule
  \textbf{Type} & \textbf{Content} \\
  \midrule
  \endhead
  \bottomrule
  \endfoot
  \code{kdtree2\_result} &
    One result: \code{dis}, the \emph{squared} distance from the query
    point, and \code{idx}, the index of the point in the input array. \\
  \code{pq} &
    Priority queue: \code{heap\_size} elements of a \code{kdtree2\_result}
    array, in max-heap order of \code{dis}. \\
  \code{kdtree2} &
    The tree: dimension \code{dimen}, number of points \code{n}, a pointer
    to the input data \code{the\_data(1:dimen,1:n)} (not copied: it must
    stay allocated while the tree is used), the permutation \code{ind} of
    the points, an optional rearranged copy of the data, the options
    \code{sort} and \code{rearrange}, and the root node. \\
  \code{tree\_node} &
    A node: the cut dimension and value, the range \code{l:u} of the
    permuted points it holds, its bounding box and the two children (none
    for a terminal node). \\
\end{longtable}

\subsection{Public Procedures}

\begin{longtable}{L{5.8cm} L{8.2cm}}
  \toprule
  \textbf{Procedure} & \textbf{Description} \\
  \midrule
  \endhead
  \bottomrule
  \endfoot
  \code{tp => kdtree2\_create(data [, indim, balanced, sort, rearrange])} &
    Builds the tree on \code{data(1:d,1:n)}, with $n > d$. \code{indim}
    restricts the search to the first \code{indim} coordinates;
    \code{balanced} splits the nodes at the median instead of the mean;
    \code{sort} returns the results in increasing distance;
    \code{rearrange} keeps a copy of the data ordered as the tree, for a
    faster search. \code{mod\_kdinterp} uses \code{sort} and
    \code{rearrange}. \\
  \code{kdtree2\_n\_nearest(tp, qv, nn, results)} &
    The \code{nn} points nearest to the vector \code{qv}, in
    \code{results(1:nn)}; \code{nn} returns the number found. \\
  \code{kdtree2\_r\_nearest(tp, qv, r2, nfound, nalloc, results)} &
    The points within the squared distance \code{r2} (boundary included);
    at most \code{nalloc} are stored, \code{nfound} returns how many were
    found. If \code{nfound} $>$ \code{nalloc} a warning is printed and the
    stored points are not the closest ones. \\
  \code{kdtree2\_r\_count(tp, qv, r2)} &
    The number of points within the squared distance \code{r2}, used to
    size the results of \code{kdtree2\_r\_nearest}. \\
  \code{kdtree2\_sort\_results(nfound, results)} &
    Sorts \code{results(1:nfound)} by increasing distance (heap sort). \\
  \code{kdtree2\_n\_nearest\_brute\_force},
  \code{kdtree2\_r\_nearest\_brute\_force} &
    The same searches by exhaustive comparison, for testing. \\
  \code{kdtree2\_destroy(tp)} &
    Releases the tree (not the input data). \\
  \code{pq\_create}, \code{pq\_insert}, \code{pq\_replace\_max} &
    Priority queue on a given result array: insertion, and replacement of
    the largest element, each returning the new largest distance. \\
\end{longtable}

\subsection{Algorithm}

\textbf{Building.} The points are split recursively. Each node takes the
coordinate with the largest spread of its points and cuts it at the mean of
the points (or the median with \code{balanced}); nodes with at most 14
points (\code{bucket\_size} $+ 1$) are terminal. If all the points fall on
one side of the cut, as for coincident points with no spread, the node is
split at the median index, so that the recursion always ends. After building
the children, the cut is moved between the two child boxes and the node box
is set to their union.

\textbf{Searching.} The search descends first into the child on the side of
the query point, then visits the other child only if the distance from the
query point to its box is not larger than the current search radius:
\begin{itemize}
  \item for $n$ nearest neighbours, the radius is the largest distance in
        the priority queue once it holds $n$ points, so that it shrinks
        while closer points are found;
  \item for a fixed radius, all the points within it are kept.
\end{itemize}
In a terminal node the squared distance of each point is accumulated over
the coordinates, stopping as soon as it exceeds the radius.

\subsection{Design Notes}

\textbf{Global search record.} The state of a search is kept in the module
variable \code{sr}, so the searches are not reentrant: two searches must not
run at the same time (for example in OpenMP threads). The callers in RegCM
search serially.

\textbf{Squared distances.} All the distances in the interface are squared,
including the radius \code{r2}; \code{mod\_kdinterp} converts them.

\textbf{Changes from the original library.} The searches around a point of
the data set with a decorrelation window, and the priority queue operations
not used by the tree (maximum, extraction, deletion), were removed. The
following were corrected: the \code{balanced} option was applied whenever
present, whatever its value; the coincident points made the tree building
recurse without end; with more points found than stored,
\code{kdtree2\_r\_nearest} sorted beyond the stored results; the
brute-force radius search excluded the points at the radius, unlike the
tree.

\textbf{Numerical checks.} On 20\,000 random points on the sphere and 2\,000
query points the tree searches agree exactly with the brute-force ones, for
4, 8 and 32 nearest neighbours (also with a balanced tree), for the
fixed-radius search and for the count, with the results sorted; a tree with
199 coincident points is built and searched correctly.

%==============================================================================
\section{Horizontal Interpolation --- \texttt{mod\_kdinterp.F90}}
\label{sec:kdinterp}
%==============================================================================

\subsection{Overview}

\code{mod\_kdinterp} interpolates fields between horizontal grids of any
shape: regular latitude--longitude grids, curvilinear grids (the model
domain, rotated or projected input grids) and unstructured point lists. It
is used by the terrain program, the ICBC programs, the CLM pre-processing and
the model (ozone, aerosol and CCN climatologies, slab ocean, CLM input). The
interpolation is in two steps:
\begin{enumerate}
  \item \code{h\_interpolator\_create} finds, for every target point, the
        source points to use and their weights, and stores them in an
        \code{h\_interpolator};
  \item the \code{h\_interpolate\_*} routines apply the weights to any
        number of fields on the same grids.
\end{enumerate}

\subsection{Dependencies}

\begin{lstlisting}
mod_kdinterp
    |-- mod_realkinds, mod_intkinds, mod_constants (erkm, d_zero, ...)
    |-- mod_stdio, mod_message
    |-- mod_earth       (ll2xyz, gcdist_simple)
    |-- mod_spbarcoord  (spherical_barycentric)
    |-- mod_kdtree2     (kd-tree search)
\end{lstlisting}

\subsection{Types}

\begin{longtable}{L{3.2cm} L{10.8cm}}
  \toprule
  \textbf{Type} & \textbf{Content} \\
  \midrule
  \endhead
  \bottomrule
  \endfoot
  \code{h\_interpolator} &
    The source shape \code{sshape(2)} and the target \code{tg}. \\
  \code{ftarget} &
    The target shape \code{tshape(2)} and, for each target point, an
    \code{ijwgt}. \\
  \code{ijwgt} &
    The number \code{np} of source points (0: no value) and their
    \code{pwgt}. \\
  \code{pwgt} &
    The indices \code{i}, \code{j} of a source point and its weight
    \code{wgt}. \\
\end{longtable}

The missing value is $-9999$ (\code{h\_missing\_value}); input values below
$-9990$ are treated as missing.

\subsection{Creating an Interpolator}

\begin{lstlisting}
call h_interpolator_create(h_i, slat, slon, tlat, tlon [, ds, roi, mxp])
call h_interpolator_create(h_i, ni, slat, slon, no, tlat, tlon)
\end{lstlisting}

\begin{longtable}{L{3.6cm} L{3.4cm} L{7cm}}
  \toprule
  \textbf{Variant} & \textbf{Source $\to$ target} & \textbf{Search} \\
  \midrule
  \endhead
  \bottomrule
  \endfoot
  \code{interp\_create\_ll\_g} &
    regular (1-D \code{slat}, \code{slon}) $\to$ curvilinear (2-D) &
    with \code{ds} [km] \\
  \code{interp\_create\_g\_g} &
    curvilinear $\to$ curvilinear &
    with \code{ds} [km] \\
  \code{interp\_create\_ll\_ll} &
    regular $\to$ regular &
    from the cell sizes \\
  \code{interp\_create\_g\_ll} &
    curvilinear $\to$ regular &
    from the cell sizes \\
  \code{interp\_create\_ll\_ll\_1d} &
    point list $\to$ point list &
    4 nearest points \\
\end{longtable}

The source points are converted to unit vectors (\code{ll2xyz}) and put in a
$k$-d tree (Section~\ref{sec:kdtree2}). The size of a grid cell, $d$, is
the largest diagonal of its cells (\code{maxedis}, in km). The search for
each target point is then one of:

\begin{itemize}
  \item \textbf{Interpolation} (the source is at least as fine as the
        target): the 4 nearest source points. This is the default, and the
        only mode of the point-list variant.
  \item \textbf{Averaging} (the source is much finer than the target: its
        cell size is less than half of \code{ds}, or less than the target
        cell size for the variants without \code{ds}): all the source
        points within the radius $r = $\,\code{ds}$\cdot$\code{roi}, or
        $r = d_\text{target}\cdot$\code{roi} without \code{ds}
        (\code{roi} defaults to 1), up to \code{mxp} (default 32; 0 for no
        limit), the nearest ones if there are more. A target point with
        fewer than 4 source points in the radius gets no value.
\end{itemize}

\noindent The kd-tree distances are chords on the unit sphere, so the radius
is $d/R$ with $R$ the Earth radius in km (\code{erkm}).

\textbf{Weights.} If a source point coincides with the target point
(squared chord below machine epsilon), it gets weight 1. Otherwise the
weights are the spherical barycentric coordinates of the target point with
respect to the source points found (Section~\ref{sec:spbarcoord}).

\subsection{Interpolating}

\begin{longtable}{L{4.2cm} L{9.8cm}}
  \toprule
  \textbf{Interface} & \textbf{Description} \\
  \midrule
  \endhead
  \bottomrule
  \endfoot
  \code{h\_interpolate\_cont(h\_i, g, f)} &
    Continuous fields, 1-D to 5-D (the extra dimensions are interpolated
    level by level), single or double precision. The value is the weighted
    mean of the valid source points, with the weights normalised over them;
    missing if none is valid. On 2-D and 3-D fields a smoothing--desmoothing
    filter follows (\code{smtdsmt}: four passes of a 1-2-1 filter with
    coefficients $\nu = 0.5$ and $\nu = -0.52$ along each direction,
    skipping missing values), which damps the shortest wavelengths. The
    result is limited to the range of the source field. \\
  \code{h\_interpolate\_cont(h\_i, g, f, ftmp, gnj, gni, nj, ni, n3)} &
    The 3-D variant for OpenACC (\code{interp\_3d\_new}), processing all
    the levels together, with the work array \code{ftmp}. Same result as
    the level-by-level version. \\
  \code{h\_interpolate\_nn(h\_i, g, f)} &
    The value of the valid source point with the largest weight (1-D and
    2-D). \\
  \code{h\_interpolate\_class(h\_i, g, f)} &
    Categories. For real fields (2-D to 5-D), the value of the source point
    with the largest weight. For integer fields, the most frequent class
    among the source points; a target point with no source points gets the
    lowest class. \\
  \code{h\_interpolate\_class(h\_i, g, f, iw, pct)} &
    Land use: the most frequent class, but if it is the water class
    \code{iw} and covers less than \code{pct}\,\% of the source points, the
    next most frequent one; missing without source points. \\
  \code{h\_interpolate\_sum(h\_i, g, f)} &
    The mean of the valid source points, without weights (class fractions
    of the terrain program); missing if none is valid. \\
  \code{h\_interpolator\_destroy(h\_i)} &
    Releases the weights. \\
\end{longtable}

\subsection{Design Notes}

\textbf{Reuse.} Finding the neighbours is the expensive part; the
interpolator is computed once and applied to every field and time step on
the same grids.

\textbf{Smoothing and range.} The smoothing of the 2-D and 3-D continuous
fields removes the small-scale noise that the interpolation from a few
neighbours may create, and the final limit to the source range prevents
overshoots of the barycentric weights and of the filter. The 1-D
interpolation is neither smoothed nor (since October 2026) limited on
missing points.

\textbf{Corrections.} The following were corrected in October 2026: the
progress dots of the constructors divided by zero for targets with fewer
than 10 rows or points; the 1-D interpolation replaced missing values with
the source minimum (or $10^{20}$); \code{h\_interpolate\_sum} averaged the
missing values ($-9999$) of the source into the result; the land use class
of a target without source points was the lowest class instead of missing.

\textbf{Numerical checks.} From a global $1^\circ$ grid to a
$40\times30$ curvilinear grid: a constant field is reproduced exactly; a
field $280 + 20\cos\varphi\cos 2\lambda$~K is interpolated within
0.003~K; the level-by-level and OpenACC 3-D variants give the same result;
missing values stay missing in the 1-D interpolation and are excluded from
the class fraction means.
These checks are in \code{Testing/share/test\_kdinterp.F90}, and the
searches of Section~\ref{sec:kdtree2} are compared with the brute force
ones in \code{Testing/share/test\_kdtree2.F90}.

%==============================================================================
\section{Dynamic Memory --- \texttt{mod\_space.F90} and \texttt{mod\_memutil.F90}}
\label{sec:memutil}
%==============================================================================

\subsection{Overview}

All the arrays of RegCM whose size is known only at run time are pointers
allocated by \code{getmem} and released by \code{relmem}. The two modules
share the work:
\begin{itemize}
  \item \code{mod\_space} defines a container type for each kind and rank,
        holding one allocatable pointer, and the generic \code{getspc} that
        allocates it with the given bounds and initial value;
  \item \code{mod\_memutil} keeps every allocated container in a list (a
        \emph{pool}) for each kind and rank, so that an array can be
        released by its pointer alone and all the arrays released at the end
        of the run.
\end{itemize}

\noindent Pointers rather than allocatable arrays let several names refer to
the same storage (\code{assignpnt}), and let the model data types hold
references to arrays allocated elsewhere.

\subsection{Dependencies}

\begin{lstlisting}
mod_space
    |-- mod_intkinds, mod_realkinds
    |-- mod_date (rcm_time_and_date)
mod_memutil
    |-- mod_intkinds, mod_realkinds, mod_constants
    |-- mod_space
    |-- mod_message (checkalloc)
    |-- mod_date (rcm_time_and_date)
\end{lstlisting}

\subsection{Kinds, Ranks and Initial Values}

\begin{longtable}{L{1.4cm} L{4.2cm} L{2.6cm} L{5.4cm}}
  \toprule
  \textbf{Suffix} & \textbf{Type} & \textbf{Ranks} &
  \textbf{Initial value} \\
  \midrule
  \endhead
  \bottomrule
  \endfoot
  \code{l} & \code{logical} & 1--5 & \code{.false.} \\
  \code{s} & \code{integer(ik2)} & 1--5 & $-1$ \\
  \code{i} & \code{integer(ik4)} & 1--5 & $-1$ \\
  \code{r} & \code{real(rk4)} & 1--5 & 0 \\
  \code{d} & \code{real(rk8)} & 1--5 & 0 \\
  \code{t} & \code{rcm\_time\_and\_date} & 1 & default of the type \\
\end{longtable}

\noindent The containers of \code{mod\_space} are named after the kind and
rank (\code{larr1d}, \code{sarr2d}, \code{iarr3d}, \code{r4arr4d},
\code{r8arr5d}, \code{rtarr1d}, \ldots); \code{bounds} holds the lower and
upper bound of one dimension.

\subsection{Interfaces}

\begin{longtable}{L{4.6cm} L{9.4cm}}
  \toprule
  \textbf{Interface} & \textbf{Description} \\
  \midrule
  \endhead
  \bottomrule
  \endfoot
  \code{memory\_init} &
    Creates the empty pools. Called once at the start of every program. \\
  \code{getmem(a, l1, h1, [\ldots, l5, h5,] vn)} &
    Allocates \code{a(l1:h1, \ldots)} with the initial value of its kind
    and adds it to the pool. If \code{a} is already associated it is first
    released. \code{vn} names the array in the error message if the
    allocation fails (\code{checkalloc}, which stops the program). \\
  \code{relmem(a)} &
    Deallocates the array \code{a} points to, removes it from the pool and
    nullifies \code{a}. Nothing happens if \code{a} is not associated or
    does not point to a pool array. \\
  \code{memory\_destroy} &
    Deallocates all the arrays still in the pools, and the pools. \\
  \code{assignpnt(a, b)} &
    Points \code{b} to the whole of \code{a}, with the same bounds (ranks
    1--5). \\
  \code{assignpnt(a, b, k [, l])} &
    Points \code{b} to the slice \code{a(\ldots, k)} or
    \code{a(:, :, k, l)}, with the bounds of the remaining dimensions:
    1-D from 2-D (real), 2-D from 3-D or 4-D, 3-D from 4-D, 4-D from
    5-D. \\
  \code{remappnt4(a, b)} &
    Points \code{b} to the 4-D real array \code{a} with the third
    dimension starting at 1. \\
\end{longtable}

\noindent The slice and the null cases of \code{assignpnt} are the same for
all kinds: if \code{a} is not associated, \code{b} is nullified.

\subsection{Pool Structure}

Each pool is a singly linked list from a root \code{r} to a last element
\code{l}. The last element is always empty: \code{getmem} allocates the new
array in it and appends a new empty element, which becomes \code{l}. Every
element before the last one therefore holds an array, and \code{relmem}
walks the list from the root up to the last element, comparing the address
of each array with its argument, and unlinks the one it finds. The cost of
\code{relmem} (and of \code{getmem} on an associated pointer) grows with the
number of arrays of that kind and rank; most arrays are allocated once at
the start and never released.

\subsection{Design Notes}

\textbf{Zero size arrays.} The address comparison
(\code{associated(a, target)}) is always false for an array of zero size, so
before October 2026 such arrays were never released: \code{relmem} left
the pointer associated and the element in the pool, and every
\code{getmem} on it added a new element and made the following searches of
that pool slower. A zero size pointer now releases any zero size array of
the pool; as such arrays hold no data, which one is released does not
matter.

\textbf{Pointers to pool arrays.} \code{relmem} nullifies only its
argument, and \code{memory\_destroy} no pointer at all; any other pointer
to the same array (from \code{assignpnt}, or a pointer assignment) is left
dangling. A pointer obtained from \code{assignpnt} must not be passed to
\code{relmem}: it is not found in the pool and nothing is released.

\textbf{Corrections.} The dead code of the list removal (a branch for the
last element, which is never a pool array) and the commented-out
\code{memshare} interface were removed in October 2026.

\textbf{Testing.} \code{Testing/share/test\_memutil.F90} checks the bounds
and initial values, the release of the first, a middle and the last
element of a pool, the reallocation of an associated pointer, zero size
arrays, \code{assignpnt} and \code{remappnt4}; it runs clean under
\code{valgrind}. The tests in \code{Testing/share} are built and run by
\code{make check}.

%==============================================================================
\section{NetCDF Helpers --- \texttt{mod\_nchelper.F90}}
\label{sec:nchelper}
%==============================================================================

\subsection{Overview}

\code{mod\_nchelper} wraps the netCDF-Fortran calls used by the
pre-processing programs, the post-processing tools and the model to open
and create files, define and check dimensions and variables, read and write
whole variables, and read and write attributes. Every wrapper checks the
status of the library call and stops the program with the source file, line
and a message (\code{checkncerr}) on error, unless an optional argument asks
to tolerate a missing name. The model output streams use the separate
\code{mod\_ncstream} (Section~\ref{sec:ncstream}).

\subsection{Dependencies}

\begin{lstlisting}
mod_nchelper
    |-- netcdf
    |-- mod_intkinds, mod_realkinds, mod_constants, mod_stdio
    |-- mod_memutil   (getmem for the coordinate arrays)
    |-- mod_dynparam  (jx, iy, kz, ds, nsg, iomode, debug_level,
    |                  projection parameters)
    |-- mod_message   (die)
\end{lstlisting}

\subsection{Conventions}

\begin{itemize}
  \item \textbf{Variable index lists.} The \code{define\_*}, \code{add\_*} and
        \code{write\_var*} routines take an index \code{ipnt} into an array
        of dimension or variable ids and advance it after each definition
        or write, so that a sequence of calls fills the array in order and
        the writes follow the same order.
  \item \textbf{Optional missing names.} The numeric readers,
        \code{check\_var} and
        \code{ncd\_inqdim} accept \code{lerror}: if it is \code{.true.} on
        input and the name is missing, it is set to \code{.false.} and the
        routine returns; otherwise a missing name stops the program.
        \code{ncd\_inqdim} also accepts \code{lexist}, set to whether the
        dimension exists; a missing dimension returns \code{dlen} $= -1$.
  \item \textbf{File format.} \code{createfile\_withname} uses
        \code{iomode} of \code{mod\_dynparam} (netCDF-4, CDF5 or classic
        format, chosen at configure time) and disables the fill values.
        \code{regcm\_vartype} is the type of the variables of
        \code{add\_variable}: double precision, or single when the code is
        built with \code{SINGLE\_PRECISION\_REAL}, as the working real kind
        \code{rkx}. With netCDF-4 and
        compression enabled, the variables get the configured filter.
  \item \textbf{Grid coordinates.} \code{jx} and \code{iy} hold the
        projection coordinates in km, centred on the domain:
        $x_i = (i - (n+1)/2)\,\Delta s$, with the subgrid distance
        $\Delta s/n_\text{sg}$ when the dimensions are not those of the
        model grid. \code{kz} holds $\sigma$ (hydrostatic and
        non-hydrostatic cores, with \code{ptop} in hPa) or $\zeta$ with the
        terms \code{a} and \code{b} of the height formula (MOLOCH).
\end{itemize}

\subsection{Interfaces}

\begin{longtable}{L{5.4cm} L{8.6cm}}
  \toprule
  \textbf{Routine} & \textbf{Description} \\
  \midrule
  \endhead
  \bottomrule
  \endfoot
  \multicolumn{2}{l}{\textit{Files}} \\
  \code{createfile\_withname(fname, ncid)} &
    Creates a file in define mode. \\
  \code{openfile\_withname(fname, ncid)} &
    Opens a file read only. \\
  \code{closefile(ncid)} & Closes a file. \\
  \code{checkncerr(ival, file, line, msg)} &
    Stops with the netCDF error message if \code{ival} is an error. \\
  \multicolumn{2}{l}{\textit{Definitions}} \\
  \code{add\_common\_global\_params(ncid, prgname, lsub)} &
    Global attributes of every RegCM file: title, institution, source,
    conventions, history (creation time), references, git revision,
    experiment (domain name), projection and its parameters, grid
    distance (of the subgrid if \code{lsub}). \\
  \code{define\_basic\_dimensions(ncid, nx, ny, nz, ipnt, idims)} &
    Dimensions \code{jx}, \code{iy} and \code{kz}. \\
  \code{add\_dimension(ncid, name, n, ipnt, idims)} &
    Another dimension; $n = -1$ for the unlimited one. \\
  \code{define\_horizontal\_coord}, \code{write\_horizontal\_coord} &
    The \code{jx} and \code{iy} coordinate variables (see Conventions). \\
  \code{define\_vertical\_coord}, \code{write\_vertical\_coord},
  \code{write\_vertical\_coord\_zita} &
    The \code{kz} coordinate, with the CF attributes of the vertical
    coordinate of the dynamical core (\code{idynamic}), and \code{ptop} or
    \code{a}, \code{b}. \\
  \code{define\_cross\_geolocation\_coord(ncid, idims, ipnt, ivar)} &
    The \code{xlat} and \code{xlon} variables. \\
  \code{add\_variable(ncid, name, long\_name, units, idims, ipnt, ivars)} &
    A variable of type \code{regcm\_vartype} with its long name and
    units. \\
  \code{add\_attribute(ncid, name, value [, ivar])} &
    A text attribute, global without \code{ivar}. \\
  \multicolumn{2}{l}{\textit{Writing and reading data}} \\
  \code{write\_var1d\_static}, \code{write\_var2d\_static},
  \code{write\_var3d\_static} &
    Write a whole variable: 1-D of any type, 2-D and 3-D single precision.
    With \code{debug\_level} $> 2$ the file is synchronised after each
    write. \\
  \code{read\_var1d\_static}, \code{read\_var2d\_static},
  \code{read\_var3d\_static} &
    Read a variable into an array, by name: single, double or integer,
    passing the array (pointer or assumed shape) or its sizes;
    \code{istart} and \code{icount} read a part of a 2-D or 3-D
    variable. \\
  \multicolumn{2}{l}{\textit{Inquiry}} \\
  \code{check\_dims(ncid)} &
    Stops if \code{jx}, \code{iy} or \code{kz} of the file differ from the
    namelist. \\
  \code{check\_dimlen(ncid, name, n)} &
    Whether the dimension exists with length \code{n}. \\
  \code{ncd\_inqdim(ncid, name [, dlen, lerror, lexist])} &
    Length or existence of a dimension. \\
  \code{check\_var(ncid, name [, lerror])} &
    Whether a variable exists. \\
  \code{get\_attribute(ncid, name, value [, ivar])} &
    Reads a text, integer, single or double attribute, global without
    \code{ivar}. \\
  \code{nf90\_myget\_att\_text(ncid, ivar, name, value)} &
    Reads a text attribute and returns the netCDF status. \\
  \code{cdumlogical(cdum, l)} & \code{'Yes'} or \code{'No'}. \\
\end{longtable}

\subsection{Design Notes}

\textbf{String attributes.} netCDF-4 has two kinds of text attributes: the
classic character arrays and variable length strings (\code{NC\_STRING}),
which Python tools such as \code{netCDF4} and \code{xarray} write by
default. The Fortran interface reads only the first kind, so
\code{nf90\_myget\_att\_text} reads a single string attribute through the C
function \code{nc\_get\_att\_string}, copies it, and frees the C string;
\code{get\_attribute} uses it for text attributes. The C variable ids are
those of Fortran minus one, so the global id 0 becomes \code{NC\_GLOBAL}
($-1$).

\textbf{Corrections.} The following were corrected in October 2026, and
are checked by \code{Testing/share/test\_nchelper.F90} (a file written and
read back):
\begin{itemize}
  \item \code{get\_attribute} on a global string attribute passed an
        undefined variable id to the C library and stopped with
        \emph{variable not found}; it now uses
        \code{nf90\_myget\_att\_text}, which also frees the string it
        reads;
  \item \code{ncd\_inqdim} with \code{lexist} returned before setting
        \code{dlen};
  \item \code{define\_horizontal\_coord} used the model grid size and
        distance for any dimension, so the subgrid coordinates of
        \code{mksurfdata} spanned $n_\text{sg}$ times the domain;
  \item the \code{formula\_terms} of the $\sigma$ coordinate were written
        as \code{kz: sigma} instead of the CF form \code{sigma: kz};
  \item two error messages named the wrong attribute or variable, and
        the definitions of \code{a} and \code{b} were not checked.
\end{itemize}
In \code{Main/mod\_ncio.F90} the parallel reading of the domain file took
the soil temperature initialisation flag from the
\code{initialized\_soil\_moisture} attribute; it now reads
\code{initialized\_soil\_temperature}.

%==============================================================================
\section{Output Streams --- \texttt{mod\_ncstream\_types.F90} and
         \texttt{mod\_ncstream.F90}}
\label{sec:ncstream}
%==============================================================================

\subsection{Overview}

\code{mod\_ncstream} writes the netCDF files of RegCM: the domain file of
\code{terrain}, the ICBC and SST files of the pre-processing, and all the
model output (\code{mod\_ncout}). A file is an \emph{output stream}: the
caller describes each variable once with a derived type that points to its
data, and the module defines the dimensions, the coordinate variables, the
grid mapping and the CF attributes from the variable descriptions and the
model parameters of \code{mod\_dynparam}. \code{mod\_ncstream\_types} holds
the types. The writes go through \code{mod\_async\_netcdf}, which can
perform them in a background thread; with \code{PNETCDF} the module uses
the PnetCDF library instead.

\subsection{Dependencies}

\begin{lstlisting}
mod_ncstream_types
    |-- mod_intkinds, mod_realkinds, mod_constants, mod_date
    |-- mpi_f08 (communicator and info types)
mod_ncstream
    |-- mod_ncstream_types
    |-- mod_dynparam (grid, projection and vertical parameters)
    |-- mod_date, mod_zita (MOLOCH vertical coordinate)
    |-- mod_stdio, mod_message
    |-- netcdf and mod_async_netcdf, or pnetcdf
\end{lstlisting}

\subsection{Using an Output Stream}

\begin{lstlisting}
type(nc_output_stream) :: ncout
type(ncoutstream_params) :: opar
type(ncvariable2d_mixed) :: vtopo
type(ncvariable3d_mixed) :: vta

opar%fname = 'out.nc'          ! file name
opar%pname = 'myprog'          ! program name in the history attribute
opar%zero_date = idate1        ! start of the first time bound
opar%l_bound = .true.          ! grid with the boundary points
call outstream_setup(ncout,opar)          ! create the file

vtopo%vname = 'topo'                       ! describe the variables
vtopo%vunit = 'm'
vtopo%long_name = 'Surface height'
vtopo%standard_name = 'surface_altitude'
vtopo%lrecords = .false.                   ! static
vtopo%rval => topo                         ! data (pointer)
vta%vname = 'ta'
...
vta%lrecords = .true.                      ! one record per time
vta%rval => ta
call outstream_addvar(ncout,vtopo)         ! define them
call outstream_addvar(ncout,vta)
call outstream_addatt(ncout,ncattribute_string('name','value'))
call outstream_enable(ncout,sigma)         ! end of definitions

call outstream_writevar(ncout,vtopo)       ! static data
do                                         ! each output time
  call outstream_addrec(ncout,idate)       !   new record
  call outstream_writevar(ncout,vta)       !   its data
end do
call outstream_dispose(ncout)              ! close the file
\end{lstlisting}

\noindent Definitions (\code{addvar}, \code{addatt}, \code{addvaratt}) are
ignored after \code{outstream\_enable}, and writes before it. Variables are
written from the arrays their pointers point to at the time of the call.

\subsection{Stream Parameters}

\begin{longtable}{L{3.2cm} L{10.8cm}}
  \toprule
  \textbf{Field} & \textbf{Meaning} \\
  \midrule
  \endhead
  \bottomrule
  \endfoot
  \code{fname}, \code{pname} & File name; program name for \code{history}. \\
  \code{zero\_date} & Start of the first time bound. \\
  \code{l\_bound} & The horizontal dimensions include the boundary points
    ($jx \times iy$); otherwise they are the internal cross points
    ($(jx-3) \times (iy-3)$), as in the model output. \\
  \code{l\_band}, \code{l\_crm} & Periodic in $x$ or in both directions:
    full $jx$ or $iy$. \\
  \code{l\_subgrid} & Subgrid file: horizontal sizes times $n_\text{sg}$,
    grid distance divided by $n_\text{sg}$. \\
  \code{l\_full\_sigma} & $kz+1$ levels (interfaces) instead of $kz$. \\
  \code{l\_plev} & Pressure levels (\code{npgwlev}) instead of $\sigma$. \\
  \code{l\_specint} & Spectral interval coordinates. \\
  \code{l\_keep} & Open an existing file and find its definitions instead
    of creating them (appending records). \\
  \code{l\_sync} & \code{outstream\_sync} flushes the file. \\
  \code{mpicom}, \code{mpinfo}, \code{mpi\_iotype} & Parallel I/O: the
    communicator of the writing processes and the netCDF parallel mode. \\
  \code{global\_jstart} \ldots\ \code{global\_iend} & With parallel I/O,
    the part of the grid written by this process. \\
\end{longtable}

\subsection{Variables}

The variable types combine a rank and a kind:
\code{ncvariable}$n$\code{d\_}\textit{kind}, $n = 0, \ldots, 4$, with kind
\code{real} (single precision data and file), \code{double} (double
precision), \code{integer}, and \code{mixed}: data of the working kind
\code{rkx} written as single precision. All variables have \code{vname},
\code{vunit}, \code{long\_name}, \code{standard\_name}, \code{notes},
\code{cell\_method} (default \code{time: point}), \code{lrecords} (time
dependent), \code{lfillvalue} with \code{rmissval} or \code{imissval}, and,
from rank 1, the string \code{axis} that gives the dimensions in order, one
character each:

\begin{longtable}{L{1.2cm} L{4.6cm} L{1.2cm} L{6.6cm}}
  \toprule
  \textbf{Code} & \textbf{Dimension} & \textbf{Code} & \textbf{Dimension} \\
  \midrule
  \endhead
  \bottomrule
  \endfoot
  \code{x}, \code{y} & \code{jx}, \code{iy} &
  \code{s} & \code{soil\_layer} \\
  \code{z} & \code{kz} &
  \code{d} & \code{depth} (lake levels) \\
  \code{t} & \code{time} (added for \code{lrecords}) &
  \code{M} & \code{month} \\
  \code{b} & \code{bnds} (time bounds) &
  \code{S}, \code{B} & spectral intervals, their bounds \\
  \code{T} & \code{ntex} (soil textures) &
  \code{n} & \code{nv} (cell corners) \\
  \code{2}, \code{w}, \code{q}, \code{W}, \code{Q} &
    heights 2, 10, 50, 100, 150~m & & \\
\end{longtable}

\noindent The default axes are \code{xy}, \code{xyz} and \code{xyzd} for
ranks 2 to 4. Variables with \code{x} and \code{y} are \emph{gridded}: they
get the \code{coordinates} (cross, dot or staggered points by name) and
\code{grid\_mapping} attributes. Further fields select what is written:
\code{j1:j2}, \code{i1:i2} (\code{k1:k2}) a window of the data array,
\code{is\_slice} with \code{rval\_slice} and the argument \code{is} of
\code{outstream\_writevar} a slice of an array with one more dimension, and
\code{rval\_level} one level of a 3-D variable.

\subsection{What the Module Adds}

\begin{itemize}
  \item \textbf{Global attributes}: title, institution, source,
        conventions (CF-1.10), history, references, model revision,
        experiment, projection and its parameters, grid size (of the
        subgrid), index of the projection origin, dynamical core and, for
        MOLOCH, the parameters of the $\zeta$ coordinate.
  \item \textbf{Coordinates}: \code{jx}, \code{iy} in m centred on the
        domain, or \code{rlon}, \code{rlat} in degrees for \code{ROTLLR};
        \code{kz} ($\sigma$, with \code{ptop} in hPa) or $\zeta$ with the
        terms \code{a} and \code{b}, or \code{plev}; the height, soil and
        spectral coordinates for the dimensions used.
  \item \textbf{Grid mapping}: a scalar \code{crs} with the CF
        \code{grid\_mapping\_name} and parameters of the projection, its
        PROJ string (\code{proj4\_params}) and WKT (\code{crs\_wkt}) on a
        sphere of radius \code{earthrad}.
  \item \textbf{Time}: \code{time} in hours since 1950-01-01 00:00:00~UTC
        with the model calendar; \code{outstream\_addrec} adds a record at a
        date. If a variable has a \code{cell\_method} with a statistic over
        time (\code{time: mean}, \code{max}, \code{min}, \code{sum}), the file
        gets \code{time\_bnds}: each record spans from the previous record
        (from \code{zero\_date} for the first) to its own time.
  \item \textbf{Compression}: with netCDF-4 and filters, variables with
        more than three dimensions (time included) in a serial file get the
        configured filter.
\end{itemize}

\subsection{Other Interfaces}

\begin{longtable}{L{5.2cm} L{8.8cm}}
  \toprule
  \textbf{Routine} & \textbf{Description} \\
  \midrule
  \endhead
  \bottomrule
  \endfoot
  \code{outstream\_addvaratt(ncout, var, att)} &
    An attribute of a variable; attributes are
    \code{ncattribute\_string}, \code{\_logical} (written \code{Yes} or
    \code{No}), \code{\_integer}, \code{\_real4}, \code{\_real8} and their
    arrays (up to 8 values). \\
  \code{outstream\_writevar(ncout, var [, lcopy, is])} &
    Writes a variable, in the current record if time dependent. With
    \code{lcopy = .false.} the internal buffer is written as it is (used for
    the coordinates). \\
  \code{outstream\_writevar\_rank0(ncout, var, myrank)} &
    A static 3-D \code{mixed} variable that only process 0 holds (the
    cell corner coordinates), in independent mode in parallel files. \\
  \code{outstream\_sync(ncout)} & Flushes the file if \code{l\_sync}. \\
  \code{outstream\_async\_flush}, \code{outstream\_async\_finalize} &
    Wait for the pending writes; stop the writer thread. \\
  \code{outstream\_netcdf\_lock}, \code{\_unlock},
  \code{outstream\_async\_defer\_wake\_begin}, \code{\_end} &
    Exclusive access to the netCDF library and batching of the writes of
    an output step (Section~\ref{sec:async_netcdf}). \\
\end{longtable}

\subsection{Design Notes}

\textbf{One buffer per stream.} \code{outstream\_enable} allocates one
integer, one single and one double buffer per stream, as large as the
largest variable of the stream; every write copies the data (window or
slice) into it, converting \code{mixed} data to single precision, and
writes the buffer. With OpenACC the buffer is page-locked host memory, and
the copy is done on the device. With \code{BITSHAVE} the single precision
values are truncated to 15 bits of mantissa before writing, which improves
the compression.

\textbf{Corrections.} The following were corrected in October 2026, and
are checked by \code{Testing/share/test\_ncstream.F90} (a file written and
read back with the netCDF library):
\begin{itemize}
  \item the \code{\_FillValue} of integer variables was written as a single
        precision attribute, which the library rejects;
  \item the time bounds were triggered by a test of the characters of the
        first nine of \code{cell\_methods} (\code{verify}), not of the
        words; they now follow the presence of \code{time: mean},
        \code{max}, \code{min} or \code{sum};
  \item the name of a variable in an error message was taken from an
        absent optional argument for global attributes;
  \item the input streams (\code{instream\_*}, \code{nc\_input\_stream}),
        about 1700 lines, were removed: they were not used, and their
        record search was wrong for files with more than two records;
  \item an example program under \code{TESTNCSTREAM}, which no longer
        compiled, and the unused \code{bitshave\_nb} were removed.
\end{itemize}
In \code{PreProc/ICBC/mod\_write.F90} the CLM forcing file of
\code{clmbc} was enabled after the writes of its static fields
(coordinates, mask, topography, cell bounds), which were therefore not
written; the stream is now enabled first.

\textbf{Notes on the OpenACC path} (not changed): with both
\code{OPENACC} and \code{BITSHAVE} defined, the device pointer of the
\code{mixed} 2-D and 3-D data is not set before it is used; and the
device copy writes the whole data array, ignoring the
\code{j1:j2}, \code{i1:i2} window.

%==============================================================================
\section{Asynchronous NetCDF I/O --- \texttt{mod\_async\_netcdf.F90} and
         \texttt{regcm\_async\_netcdf\_thread.c}}
\label{sec:async_netcdf}
%==============================================================================

\subsection{Overview}

\code{mod\_async\_netcdf} lets the model hand its netCDF writes, and the
prefetch of the next boundary conditions, to a background thread, so that
the time step loop does not wait for the file system. All the writes of
\code{mod\_ncstream} go through it. The thread, a mutex and condition
variables come from the small C file \code{regcm\_async\_netcdf\_thread.c}
(POSIX threads).

The background thread exists only if both
\begin{itemize}
  \item RegCM was configured with \code{--enable-async-netcdf}
        (\code{-DASYNC\_NETCDF}), and
  \item the environment variable \code{RCM\_ASYNC\_OUTPUT\_GB} gives a
        positive memory cap, in GB, for the data waiting to be written.
\end{itemize}
Otherwise the module is \emph{disabled}: every call performs the netCDF
operation at once and returns its status, and the locks do nothing. This is
the default. With \code{PNETCDF} the module is not used.

\subsection{Dependencies}

\begin{lstlisting}
mod_async_netcdf
    |-- iso_c_binding, netcdf
    |-- mod_intkinds, mod_realkinds
    |-- cudafor (OPENACC only)
    `-- regcm_async_netcdf_thread.c (ASYNC_NETCDF only)
\end{lstlisting}

\subsection{How It Works}

\begin{itemize}
  \item \textbf{Queue.} A write copies the data into a new buffer, adds an
        item (file, variable, start, count, data) to a first-in first-out
        queue and returns. The worker thread takes the items in order and
        performs them with the netCDF library; sync and close requests are
        queued in the same way, so they act after the writes before them.
        Read requests have their own queue, served first.
  \item \textbf{Memory cap.} The bytes of the queued data are counted; a
        write that would exceed the cap waits until the worker has freed
        enough. A single item larger than the cap is accepted when the
        queue is empty, so that nothing waits for ever.
  \item \textbf{One user of the library at a time.} The netCDF library is
        not thread safe. The worker takes the netCDF lock for each item,
        and the main thread must take it with \code{async\_netcdf\_lock} /
        \code{async\_netcdf\_unlock} around its own netCDF calls (create,
        define, open, close); \code{mod\_ncstream} does so. The main thread
        has priority: a waiting main thread is served before the next item.
        The lock is not recursive.
  \item \textbf{Errors.} The status of a queued write is not known when the
        call returns. The first failure is kept and returned by the next
        call (\code{wait\_all}, \code{sync}, \code{close}, \code{shutdown},
        or any write), and every later operation fails with it: the output
        of a run stops at the first error instead of continuing silently.
  \item \textbf{Batching.} Between \code{async\_netcdf\_defer\_wake\_begin}
        and \code{\_end} the worker is not woken for each item; the items of
        an output step are then written together.
  \item \textbf{OpenACC.} The cap is allocated at start as one page-locked
        host pool, and the item buffers are slices of it. The device
        variants copy the data from the device into the slice (converting
        \code{rkx} to single precision on the device for the \code{rkx\_as\_r4}
        variants), so the time step can continue while the host writes.
\end{itemize}

\subsection{Interfaces}

\begin{longtable}{L{5.6cm} L{8.4cm}}
  \toprule
  \textbf{Routine} & \textbf{Description} \\
  \midrule
  \endhead
  \bottomrule
  \endfoot
  \code{async\_netcdf\_put\_var(ncid, varid, values [, start, count])} &
    Writes a 1-D array or a scalar of type single, double or integer. \\
  \code{async\_netcdf\_put\_var\_device},
  \code{\_rkx\_as\_r4\_device}, \code{\_rkx\_as\_r4\_device\_3d} &
    The same from device data (OpenACC); the 3-D variant writes a block of
    an \code{rkx} array as single precision, and can be forced to be
    synchronous. \\
  \code{async\_netcdf\_sync(ncid)}, \code{async\_netcdf\_close(ncid)} &
    Flush or close a file after the writes queued before. \\
  \code{async\_netcdf\_wait\_all()} &
    Waits until the queues are empty; returns the first error. \\
  \code{async\_netcdf\_shutdown()} &
    Waits for all the work, stops the thread and frees the pool. \\
  \code{async\_netcdf\_lock}, \code{async\_netcdf\_unlock} &
    The netCDF lock of the main thread. \\
  \code{async\_netcdf\_defer\_wake\_begin}, \code{\_end} &
    Batch the items queued in between (they can be nested). \\
  \code{async\_netcdf\_initialize([limit])}, \code{async\_netcdf\_is\_enabled},
  \code{async\_netcdf\_configured\_cap\_bytes} &
    Start the module (lowering the cap to \code{limit} at the first
    initialisation only), whether the worker is running, the cap from the
    environment. Any call starts the module. \\
  \code{async\_netcdf\_get\_var\_rkx(ncid, varid, values, start, count,
  counter)} &
    Queues the read of a record into \code{values}, counted on
    \code{counter}; fails when the module is disabled. \\
  \code{async\_netcdf\_counter\_reset}, \code{\_poll}, \code{\_wait},
  \code{\_fail} &
    Pending reads and status of a counter. \\
  \code{async\_netcdf\_try\_acquire\_buffer}, \code{\_release\_buffer} &
    A buffer of \code{rkx} values counted in the cap (from the pool with
    OpenACC), for the reads; fails when the module is disabled or the cap
    is reached. \\
\end{longtable}

\noindent The reads, counters and buffers are used by \code{mod\_ncio} to
read the next ICBC record while the model integrates.

\subsection{Design Notes}

\textbf{Testing.} \code{Testing/share/test\_async\_netcdf.F90} writes 200
records of three types and a scalar, with a sync in the middle and
batching, closes the file, and checks it; with the worker it also checks a
queued read, the buffers, and that an error is reported. \code{make check}
runs it with a cap of 50~kB, so that in a build with
\code{--enable-async-netcdf} the writes wait for space. The same test was
run with the worker under \code{valgrind --tool=drd} without reports in
RegCM code.

\textbf{Changes.} In October 2026 three functions that only forwarded to
the direct netCDF writes were removed, and the lines longer than 80
characters were wrapped; no change of behaviour.
